%% CSoNet 2026 Journal Track submission (Journal of Combinatorial Optimization)
%DIF LATEXDIFF DIFFERENCE FILE
%DIF DEL main_old.tex   Mon Oct  5 07:48:20 2026
%DIF ADD main.tex       Mon Oct  5 07:48:20 2026
%% Build: pdflatex main && bibtex main && pdflatex main && pdflatex main

\documentclass[pdflatex,sn-mathphys-ay]{sn-jnl}

\usepackage{amsmath,amssymb,amsthm}
\usepackage{booktabs}
\usepackage{graphicx}
\usepackage{algorithm}
\usepackage{algpseudocode}

\theoremstyle{thmstyleone}
\newtheorem{theorem}{Theorem}
\newtheorem{proposition}[theorem]{Proposition}
\newtheorem{corollary}[theorem]{Corollary}
\newtheorem{lemma}[theorem]{Lemma}
%DIF 17a17
\newtheorem{fact}{Fact} %DIF > 
%DIF -------
\theoremstyle{thmstyletwo}
\newtheorem{assumption}{Assumption}
\newtheorem{remark}{Remark}
\newtheorem{definition}{Definition}
\theoremstyle{thmstylethree}
\newtheorem{example}{Example}

\raggedbottom
%DIF PREAMBLE EXTENSION ADDED BY LATEXDIFF
%DIF UNDERLINE PREAMBLE %DIF PREAMBLE
\RequirePackage[normalem]{ulem} %DIF PREAMBLE
\RequirePackage{color}\definecolor{RED}{rgb}{1,0,0}\definecolor{BLUE}{rgb}{0,0,1} %DIF PREAMBLE
\providecommand{\DIFadd}[1]{{\protect\color{blue}\uwave{#1}}} %DIF PREAMBLE
\providecommand{\DIFdel}[1]{{\protect\color{red}\sout{#1}}}                      %DIF PREAMBLE
%DIF SAFE PREAMBLE %DIF PREAMBLE
\providecommand{\DIFaddbegin}{} %DIF PREAMBLE
\providecommand{\DIFaddend}{} %DIF PREAMBLE
\providecommand{\DIFdelbegin}{} %DIF PREAMBLE
\providecommand{\DIFdelend}{} %DIF PREAMBLE
\providecommand{\DIFmodbegin}{} %DIF PREAMBLE
\providecommand{\DIFmodend}{} %DIF PREAMBLE
%DIF FLOATSAFE PREAMBLE %DIF PREAMBLE
\providecommand{\DIFaddFL}[1]{\DIFadd{#1}} %DIF PREAMBLE
\providecommand{\DIFdelFL}[1]{\DIFdel{#1}} %DIF PREAMBLE
\providecommand{\DIFaddbeginFL}{} %DIF PREAMBLE
\providecommand{\DIFaddendFL}{} %DIF PREAMBLE
\providecommand{\DIFdelbeginFL}{} %DIF PREAMBLE
\providecommand{\DIFdelendFL}{} %DIF PREAMBLE
\newcommand{\DIFscaledelfig}{0.5}
%DIF HIGHLIGHTGRAPHICS PREAMBLE %DIF PREAMBLE
\RequirePackage{settobox} %DIF PREAMBLE
\RequirePackage{letltxmacro} %DIF PREAMBLE
\newsavebox{\DIFdelgraphicsbox} %DIF PREAMBLE
\newlength{\DIFdelgraphicswidth} %DIF PREAMBLE
\newlength{\DIFdelgraphicsheight} %DIF PREAMBLE
% store original definition of \includegraphics %DIF PREAMBLE
\LetLtxMacro{\DIFOincludegraphics}{\includegraphics} %DIF PREAMBLE
\newcommand{\DIFaddincludegraphics}[2][]{{\color{blue}\fbox{\DIFOincludegraphics[#1]{#2}}}} %DIF PREAMBLE
\newcommand{\DIFdelincludegraphics}[2][]{% %DIF PREAMBLE
\sbox{\DIFdelgraphicsbox}{\DIFOincludegraphics[#1]{#2}}% %DIF PREAMBLE
\settoboxwidth{\DIFdelgraphicswidth}{\DIFdelgraphicsbox} %DIF PREAMBLE
\settoboxtotalheight{\DIFdelgraphicsheight}{\DIFdelgraphicsbox} %DIF PREAMBLE
\scalebox{\DIFscaledelfig}{% %DIF PREAMBLE
\parbox[b]{\DIFdelgraphicswidth}{\usebox{\DIFdelgraphicsbox}\\[-\baselineskip] \rule{\DIFdelgraphicswidth}{0em}}\llap{\resizebox{\DIFdelgraphicswidth}{\DIFdelgraphicsheight}{% %DIF PREAMBLE
\setlength{\unitlength}{\DIFdelgraphicswidth}% %DIF PREAMBLE
\begin{picture}(1,1)% %DIF PREAMBLE
\thicklines\linethickness{2pt} %DIF PREAMBLE
{\color[rgb]{1,0,0}\put(0,0){\framebox(1,1){}}}% %DIF PREAMBLE
{\color[rgb]{1,0,0}\put(0,0){\line( 1,1){1}}}% %DIF PREAMBLE
{\color[rgb]{1,0,0}\put(0,1){\line(1,-1){1}}}% %DIF PREAMBLE
\end{picture}% %DIF PREAMBLE
}\hspace*{3pt}}} %DIF PREAMBLE
} %DIF PREAMBLE
\LetLtxMacro{\DIFOaddbegin}{\DIFaddbegin} %DIF PREAMBLE
\LetLtxMacro{\DIFOaddend}{\DIFaddend} %DIF PREAMBLE
\LetLtxMacro{\DIFOdelbegin}{\DIFdelbegin} %DIF PREAMBLE
\LetLtxMacro{\DIFOdelend}{\DIFdelend} %DIF PREAMBLE
\DeclareRobustCommand{\DIFaddbegin}{\DIFOaddbegin \let\includegraphics\DIFaddincludegraphics} %DIF PREAMBLE
\DeclareRobustCommand{\DIFaddend}{\DIFOaddend \let\includegraphics\DIFOincludegraphics} %DIF PREAMBLE
\DeclareRobustCommand{\DIFdelbegin}{\DIFOdelbegin \let\includegraphics\DIFdelincludegraphics} %DIF PREAMBLE
\DeclareRobustCommand{\DIFdelend}{\DIFOaddend \let\includegraphics\DIFOincludegraphics} %DIF PREAMBLE
%DIF COLORLISTINGS PREAMBLE %DIF PREAMBLE
\RequirePackage{listings} %DIF PREAMBLE
\RequirePackage{color} %DIF PREAMBLE
\lstdefinelanguage{DIFcode}{ %DIF PREAMBLE
%DIF DIFCODE_UNDERLINE %DIF PREAMBLE
  moredelim=[il][\color{red}\sout]{\%DIF\ <\ }, %DIF PREAMBLE
  moredelim=[il][\color{blue}\uwave]{\%DIF\ >\ } %DIF PREAMBLE
} %DIF PREAMBLE
\lstdefinestyle{DIFverbatimstyle}{ %DIF PREAMBLE
	language=DIFcode, %DIF PREAMBLE
	basicstyle=\ttfamily, %DIF PREAMBLE
	columns=fullflexible, %DIF PREAMBLE
	keepspaces=true %DIF PREAMBLE
} %DIF PREAMBLE
\lstnewenvironment{DIFverbatim}{\lstset{style=DIFverbatimstyle}}{} %DIF PREAMBLE
\lstnewenvironment{DIFverbatim*}{\lstset{style=DIFverbatimstyle,showspaces=true}}{} %DIF PREAMBLE
%DIF END PREAMBLE EXTENSION ADDED BY LATEXDIFF

\begin{document}

\title[Minimum Weighted Hazard-Exposure Dispatch]{Minimum Weighted
Hazard-Exposure Dispatch: Complexity, an Exact Algorithm, and an FPTAS}

\author*[1]{\fnm{Quang-Vinh} \sur{Dang}}\email{vinh.dq4@buv.edu.vn}
\author[2]{\fnm{Minh Ngoc} \sur{Dinh}}\email{minh.dinh@maeducation.com}
\author[1]{\fnm{Hoang-Viet} \sur{Vu}}\email{30066555@st.buv.edu.vn}
\author[3]{\fnm{Phuc-Son} \sur{Nguyen}}\email{son.np3393@ueh.edu.vn}

\affil*[1]{\orgname{British University Vietnam}, \orgaddress{\city{Hung
Yen}, \country{Vietnam}}}
\affil[2]{\orgname{Millennia Education}, \orgaddress{\city{Ho Chi Minh
City}, \country{Vietnam}}}
\affil[3]{\orgname{UEH University}, \orgaddress{\city{Ho Chi Minh
City}, \country{Vietnam}}}

\DIFdelbegin %DIFDELCMD < \abstract{We study a dispatch-scheduling problem arising when a single
%DIFDELCMD < depot must send a vehicle or crew to a set of sites before a spreading
%DIFDELCMD < hazard---a flood, wildfire, contamination plume, or growing congestion
%DIFDELCMD < zone---reaches them. Each site has a fixed round-trip dispatch time, a
%DIFDELCMD < deadline at which the hazard arrives, and a criticality weight; the
%DIFDELCMD < objective is a dispatch order maximizing the total weight of sites
%DIFDELCMD < reached before the hazard arrives. We call this problem Minimum
%DIFDELCMD < Weighted Hazard-Exposure Dispatch (MWHED) and show it is exactly the
%DIFDELCMD < classical single-machine problem of maximizing weighted on-time jobs,
%DIFDELCMD < reparametrized in transportation terms. We prove MWHED is weakly
%DIFDELCMD < NP-hard via a reduction showing its equal-deadline special case is
%DIFDELCMD < exactly the 0/1 knapsack problem; give a pseudo-polynomial exact
%DIFDELCMD < dynamic program from an earliest-deadline-first exchange argument;
%DIFDELCMD < adapt value-scaling to obtain a fully polynomial-time approximation
%DIFDELCMD < scheme; and identify a second tractable special case---equal dispatch
%DIFDELCMD < times---solvable in $O(n\log n)$ time by a matroid greedy algorithm,
%DIFDELCMD < showing hardness requires heterogeneity in both dispatch cost and
%DIFDELCMD < criticality, not either alone. We further show the naive fallback of
%DIFDELCMD < dispatching in deadline order with no reconsideration has unbounded
%DIFDELCMD < worst-case loss, so an exact or near-exact algorithm is necessary, not
%DIFDELCMD < merely convenient. Numerical experiments, robust across alternative
%DIFDELCMD < weight and deadline distributions, confirm the approximation scheme
%DIFDELCMD < tracks the true optimum as instance size grows, while the naive policy
%DIFDELCMD < degrades. The results give planners exact and near-exact tools for a
%DIFDELCMD < provably hard scheduling problem, and delineate which structural
%DIFDELCMD < features of a dispatch network make hazard-exposure minimization
%DIFDELCMD < tractable.}
%DIFDELCMD < %%%
\DIFdelend \DIFaddbegin \abstract{We study a dispatch-scheduling problem arising when a single
depot must send a vehicle or crew to sites before a spreading hazard
(flood, wildfire, contamination plume) reaches them. Each site has a
round-trip dispatch time, a hazard-arrival deadline and a criticality
weight; the objective is a dispatch order maximizing the total weight of
sites served in time. We call this Minimum Weighted Hazard-Exposure
Dispatch (MWHED) and show it is exactly the classical problem
$1\|\sum w_jU_j$, so its basic picture is inherited from scheduling
theory: weak NP-hardness (the equal-deadline case is 0/1 knapsack), a
Lawler--Moore-type pseudo-polynomial dynamic program, and a value-scaling
FPTAS. We present these in a self-contained form for the transportation
audience and say precisely which are classical. The framing adds a
classification: restricted to the instances in which some of the
dispatch times, weights and deadlines are constant, the problem is
polynomial exactly when dispatch times or weights are constant and
weakly NP-hard otherwise, so heterogeneous deadlines alone never cause
hardness; and a matroid-greedy algorithm for equal dispatch times that
extends to any number of identical vehicles. Natural heuristics (deadline
order with or without skipping, a weighted greedy repair) have unbounded
worst-case ratio, and the FPTAS analysis is tight up to a factor
$1-\epsilon^2$. A numerical study with active FPTAS rounding, verified
against exhaustive search, and a Camp Fire case study with a sensitivity
analysis complete the paper.}
\DIFaddend 

\keywords{combinatorial optimization, scheduling under deadlines,
computational complexity, approximation algorithms, transportation
networks, disaster response}

\maketitle

% ══════════════════════════════════════════════════════════════════
\section{Introduction}\label{sec:intro}

A recurring operational problem in infrastructure and emergency-response
networks is deciding, from a single depot, in what order to dispatch a
vehicle or crew to a set of sites that are each racing against a
spreading hazard. A wildfire crew choosing which structures to defend
before the fire front arrives, a utility company dispatching repair
trucks to substations before a cascading outage propagates to them, and
a flood-response team evacuating or servicing neighborhoods before
rising water reaches them all face the same underlying decision: given
a known rate at which a hazard is spreading, and a known cost of
reaching each site from a common depot, in what order should sites be
visited to save as much value as possible before the hazard reaches
each one in turn?

This decision is a scheduling problem, not merely a routing one. Once
every visit is a self-contained round trip from a single depot---the
standard topology for a technician, repair-crew, or single-vehicle
emergency-response dispatch base---the cost of visiting a site does not
depend on which site was visited immediately before it, only on the
cumulative time already spent. The question is then purely one of
\emph{order}: which sites to reach early, which to defer, and which to
concede to the hazard, given that each site $i$ becomes ``exposed''
once the cumulative dispatch time exceeds its known hazard-arrival
deadline $d_i$, at a loss of $w_i$.

The three motivating settings above share a further feature that makes
the \emph{ordering} decision, rather than the underlying routing or
hazard-forecasting problem, the right object of study. In each case,
the hazard's own spread is typically produced by a separate, already
well-developed forecasting step---a wildfire-spread simulator driven by
fuel load, wind, and terrain; a cascading-failure model for a power
network; a hydrological model for a flood plain---so that, by the time
a dispatcher is choosing a visiting order, each site's arrival deadline
$d_i$ is already a given input, not something to be modeled jointly
with the dispatch decision. Similarly, the round-trip dispatch time
$p_i$ is typically fixed by the depot's location and the local road or
access network, and is known before dispatch begins. What remains
undetermined, and squarely a combinatorial decision, is the
\emph{order} in which a single crew or vehicle should visit sites once
$p_i$ and $d_i$ are fixed---exactly the question MWHED isolates. This
separation of concerns (forecast the hazard and fix the travel network
first; optimize the dispatch order second) mirrors how the wildfire
operations-research literature we discuss in Section~\ref{sec:related}
is typically organized, and is what lets us give a complete
complexity and algorithmic answer to the ordering question alone,
rather than a heuristic for the joint problem.

We formalize this as Minimum Weighted Hazard-Exposure Dispatch
(MWHED)\DIFdelbegin \DIFdel{and give a complete complexity and algorithmic picture. Our
contributions are:
}%DIFDELCMD < 

%DIFDELCMD < \begin{enumerate}
%DIFDELCMD < \item %%%
\DIFdel{We show MWHED is exactly }\DIFdelend \DIFaddbegin \DIFadd{. It is worth saying at the outset, and Section~\ref{sec:model}
proves it, that MWHED is }\emph{\DIFadd{exactly}} \DIFaddend the classical single-machine
\DIFdelbegin \DIFdel{scheduling
problem }\DIFdelend \DIFaddbegin \DIFadd{problem $1\|\sum w_jU_j$ }\DIFaddend of maximizing the weighted number of on-time
jobs\DIFdelbegin \DIFdel{, and prove it
}\DIFdelend \DIFaddbegin \DIFadd{. Consequently the basic complexity picture (weak NP-hardness, a
pseudo-polynomial dynamic program, an approximation scheme) is
inherited from scheduling theory, and we do not claim a new algorithmic
technique for the general problem. Our contributions are therefore of
two kinds. The first is a self-contained development of the classical
results in the transportation vocabulary, with complete proofs and a
running example, aimed at a reader who would not otherwise connect a
dispatch problem to $1\|\sum w_jU_j$:
}\begin{enumerate}
\item \DIFadd{MWHED }\DIFaddend is weakly NP-hard\DIFdelbegin \DIFdel{by an elementary, self-contained argument: the special case in which every site shares the same hazard deadline is }\emph{\DIFdel{exactly}} %DIFAUXCMD
\DIFdelend \DIFaddbegin \DIFadd{, by a direct reduction from
}\textsc{\DIFadd{Partition}} \DIFadd{that uses only the equal-deadline special case,
which is exactly }\DIFaddend the 0/1 knapsack problem
(Theorem~\ref{thm:hardness}).
\item \DIFdelbegin \DIFdel{We give an }\DIFdelend \DIFaddbegin \DIFadd{An }\DIFaddend exact pseudo-polynomial dynamic program \DIFdelbegin \DIFdel{(Theorem~\ref{thm:dp}), derived directly }\DIFdelend \DIFaddbegin \DIFadd{derived }\DIFaddend from an
earliest-deadline-first exchange argument \DIFdelbegin \DIFdel{that we prove from first
principles rather than assert}\DIFdelend \DIFaddbegin \DIFadd{(Theorem~\ref{thm:dp}); this
is the Lawler--Moore dynamic program \mbox{%DIFAUXCMD
\citep{lawlermoore1969}}\hskip0pt%DIFAUXCMD
}\DIFaddend .
\item \DIFdelbegin \DIFdel{We adapt the classical }\DIFdelend \DIFaddbegin \DIFadd{A }\DIFaddend value-scaling \DIFdelbegin \DIFdel{technique to obtain a fully
polynomial-time approximation scheme (FPTAS ) for MWHED
}\DIFdelend \DIFaddbegin \DIFadd{FPTAS }\DIFaddend (Theorem~\ref{thm:fptas})\DIFdelbegin \DIFdel{, so that a solution within any desired factor
$1-\epsilon$ of optimal is computable in time polynomial in the number
of sites and $1/\epsilon$.
}\DIFdelend \DIFaddbegin \DIFadd{; approximation
schemes for this problem are classical \mbox{%DIFAUXCMD
\citep{sahni1976,gensl1981}}\hskip0pt%DIFAUXCMD
, and
we present the construction together with a complete proof and an
explicit treatment of the case in which the scaling is vacuous.
}\end{enumerate}
\DIFadd{The second kind is what we believe the transportation framing adds,
and where we have tried to go beyond a relabelling:
}\begin{enumerate}
\setcounter{enumi}{3}
\DIFaddend \item \DIFdelbegin \DIFdel{We identify a second special case---equal round-trip dispatch
times, i.e., sites equidistant from the depot---that is solvable
}\emph{\DIFdel{exactly}} %DIFAUXCMD
\DIFdel{in $O(n \log n)$ time via a matroid greedy algorithm
}\DIFdelend \DIFaddbegin \emph{\DIFadd{A complete classification by homogeneity}}
\DIFaddend (Theorem~\DIFdelbegin \DIFdel{\ref{thm:equalp}), showing that heterogeneity in dispatch cost
}\emph{\DIFdel{alone}} %DIFAUXCMD
\DIFdel{does not cause hardness : it is the
combination of
heterogeneous dispatch cost and heterogeneous criticality, under a shared deadline structure, that does}\DIFdelend \DIFaddbegin \DIFadd{\ref{thm:classification}). For every restriction obtained by
requiring some of the dispatch times, weights, and deadlines to be
constant, MWHED is polynomial if dispatch times or weights are
constant and weakly NP-hard otherwise; heterogeneity of the hazard
deadlines never causes hardness by itself. Together with the
equal-cost result below this turns the informal question ``which
network structure makes dispatch easy'' into a precise statement}\DIFaddend .
\item \DIFdelbegin \DIFdel{We show the natural fallback for a planner who declines to
optimize at all---dispatching sites in increasing-deadline order and
simply accepting whichever remain on time---has }\DIFdelend \emph{\DIFdelbegin \DIFdel{unbounded}%DIFDELCMD < \MBLOCKRIGHTBRACE
%DIFDELCMD < %%%
\DIFdel{worst-case loss relative to $W^\ast$ (Proposition~\ref{prop:naive}), so
Theorems~\ref{thm:dp} and~\ref{thm:fptas} are not merely convenient but
necessary for any }\DIFdelend \DIFaddbegin \DIFadd{Equal dispatch times, with identical vehicles}}
\DIFadd{(Theorem~\ref{thm:equalp}). A matroid-greedy algorithm solves the
equal-cost case in $O(n\log n)$ time for one vehicle and, by the same
proof, for any number $m$ of identical vehicles. We prove the required
matroid property from scratch rather than cite it. In particular, the
equal-cost case does }\emph{\DIFadd{not}} \DIFadd{become hard when several vehicles are
available (it does once release dates are added; Section~\ref{sec:disc-multi}).
}\item \emph{\DIFadd{Worst-case results for natural heuristics}}
\DIFadd{(Propositions~\ref{prop:naive}, \ref{prop:repair}, \ref{prop:tight}).
Dispatching in deadline order, with or without skipping sites that
would be late, and a weighted greedy repair all have unbounded
}\DIFaddend worst-case \DIFdelbegin \DIFdel{guarantee at all. }\DIFdelend \DIFaddbegin \DIFadd{ratio; and the standard FPTAS analysis is tight up to a
factor $1-\epsilon^2$. These results say that some form of
reconsideration is needed for a guarantee; they do }\emph{\DIFadd{not}} \DIFadd{say that
the dynamic program or the FPTAS is the only way to obtain one.
}\DIFaddend \end{enumerate}
\DIFdelbegin %DIFDELCMD < 

%DIFDELCMD < %%%
\DIFdel{Together, Theorems~\ref{thm:hardness} and~\ref{thm:equalp} give a sharp
structural picture of where the difficulty in this problem lives, Theorems~\ref{thm:dp} and ~\ref{thm:fptas} give planners exact and provably near-exact tools for the general, hard case, and
Proposition~\ref{prop:naive} shows those tools are not optional: the obvious do-nothing alternative can be made to fail arbitrarily badly.
A
numerical illustration in }\DIFdelend \DIFaddbegin \DIFadd{Finally, }\DIFaddend Section~\ref{sec:experiments} \DIFdelbegin \DIFdel{confirms the FPTAStracks the true optimum closely well before its worst-case
guarantee requires, while a naive dispatch order (no optimization at
all, just visiting sites in deadline order) degrades as instances grow}\DIFdelend \DIFaddbegin \DIFadd{reports a numerical study
designed so that the FPTAS's rounding is verifiably active (with small
weights it is vacuous, as Section~\ref{sec:acc} explains), together
with exhaustive-search verification of every algorithm, and a
real-wildfire case study with a sensitivity analysis over its uncertain
inputs}\DIFaddend . To keep the argument
concrete\DIFdelbegin \DIFdel{rather than purely symbolic}\DIFdelend , Sections~\ref{sec:model} and~\ref{sec:theory} carry a single
four-site running example through \DIFdelbegin \DIFdel{every result}\DIFdelend \DIFaddbegin \DIFadd{the results}\DIFaddend , worked out by hand
alongside each proof.

The remainder of this paper is organized as follows.
Section~\ref{sec:related} positions MWHED against the scheduling,
real-time-systems, and disaster-response \DIFdelbegin \DIFdel{routing literatures }\DIFdelend \DIFaddbegin \DIFadd{literatures and states
precisely what is classical}\DIFaddend . Section~\ref{sec:model} defines MWHED \DIFdelbegin \DIFdel{formally }\DIFdelend and
establishes the earliest-deadline-first \DIFdelbegin \DIFdel{structural }\DIFdelend lemma every later result
depends on. Section~\ref{sec:theory} proves the \DIFdelbegin \DIFdel{four theorems and the proposition
}\DIFdelend \DIFaddbegin \DIFadd{results }\DIFaddend above.
Section~\ref{sec:experiments} reports \DIFdelbegin \DIFdel{four complementary
numerical experiments}\DIFdelend \DIFaddbegin \DIFadd{the numerical study and the case
study}\DIFaddend . Section~\ref{sec:discussion} discusses \DIFdelbegin \DIFdel{three
natural extensions }\DIFdelend \DIFaddbegin \DIFadd{extensions (several
vehicles, a moving depot, and uncertain parameters) }\DIFaddend and precisely what
each breaks in our analysis, and Section~\ref{sec:conclusion}
concludes. Two appendices give an exhaustive hand-verification of the
running example and a \DIFdelbegin \DIFdel{fully
}\DIFdelend spelled-out derivation of the FPTAS \DIFdelbegin \DIFdel{'s approximation-ratio }\DIFdelend bound.

% ══════════════════════════════════════════════════════════════════
\section{Related work}\label{sec:related}

\textbf{Scheduling with deadlines.} MWHED, once its transportation
framing is stripped away, is the classical single-machine problem of
maximizing the weighted number of on-time jobs, equivalently minimizing
the weighted number of late jobs, $1\|\sum w_jU_j$ in the standard
three-field notation \DIFaddbegin \DIFadd{(job $j$ has processing time $p_j$, due date $d_j$
and weight $w_j$; $U_j=1$ if it is late)}\DIFaddend . \citet{moore1968} gives an
$O(n\log n)$ algorithm for the \emph{unweighted} case
($1\|\sum U_j$): process jobs in earliest-due-date order, and whenever
the running schedule becomes infeasible, permanently discard the
\DIFaddbegin \DIFadd{scheduled }\DIFaddend job with the largest processing time \DIFdelbegin \DIFdel{among those scheduled so far. This procedure is commonly called
}\DIFdelend \DIFaddbegin \DIFadd{(}\DIFaddend the Moore--Hodgson
algorithm\DIFdelbegin \DIFdel{. }\DIFdelend \DIFaddbegin \DIFadd{). For the weighted case, }\DIFaddend \citet{lawlermoore1969} give a
\DIFdelbegin \DIFdel{general
functional-equation (dynamic programming) technique for resource
allocation and sequencing problems that specializes to a
pseudo-polynomial algorithm for the weighted case; the classification
of $1\|\sum w_jU_j$ itself as (weakly ) }\DIFdelend \DIFaddbegin \DIFadd{dynamic program over jobs in due-date order and total processing time
that runs in $O(nP)$ time, $P=\sum_jp_j$; this is our
Theorem~\ref{thm:dp}. The problem is weakly }\DIFaddend NP-hard\DIFdelbegin \DIFdel{is standard in the
scheduling-complexity literature, surveyed by
\mbox{%DIFAUXCMD
\citet{lenstra1977}}\hskip0pt%DIFAUXCMD
, and follows directly from the same reduction we
give in Section~\ref{sec:theory}: }\DIFdelend \DIFaddbegin \DIFadd{, since }\DIFaddend an instance
with a common due date is \DIFdelbegin \DIFdel{exactly }\DIFdelend a 0/1 knapsack instance, whose NP-hardness is
classical \citep{gareyjohnson1979}\DIFdelbegin \DIFdel{. We are not aware of a prior statement of this
reduction, or of }\DIFdelend \DIFaddbegin \DIFadd{; complexity surveys include
\mbox{%DIFAUXCMD
\citet{lenstra1977}}\hskip0pt%DIFAUXCMD
. Approximation schemes are also classical:
\mbox{%DIFAUXCMD
\citet{sahni1976} }\hskip0pt%DIFAUXCMD
gives rounding-based approximation algorithms for
single-processor job sequencing with deadlines, which is the same
problem, and \mbox{%DIFAUXCMD
\citet{gensl1981} }\hskip0pt%DIFAUXCMD
give a faster approximation algorithm for
it. Our Theorem~\ref{thm:fptas} is therefore a re-derivation of a known
type of result in our notation, not a new technique, and the same holds
for the common-deadline hardness argument and }\DIFaddend the equal-processing-time
\DIFdelbegin \DIFdel{tractable }\DIFdelend special case of Theorem~\ref{thm:equalp}, \DIFdelbegin \DIFdel{in transportation-network language}\DIFdelend \DIFaddbegin \DIFadd{which is the classical
unit-time job sequencing with profits. What we are not aware of, and
what Sections~\ref{sec:theory}--\ref{sec:experiments} add, is the
classification of Theorem~\ref{thm:classification}, the extension to
identical vehicles of Theorem~\ref{thm:equalp}, the heuristic lower
bounds of Propositions~\ref{prop:naive}--\ref{prop:tight}, and the
verified computational study}\DIFaddend .

This classical scheduling problem remains an active research topic:
\citet{hejl2022} give a refined pseudo-polynomial dynamic program for
the practically important ``strongly correlated'' regime (where
processing times and weights are close in value), and
\citet{hermelin2024} improve the Lawler--Moore-style DP's running time
via $(\max,+)$-convolution techniques---an algorithmic refinement
orthogonal to, and compatible with, Theorem~\ref{thm:dp}'s DP.
\citet{heegerhermelin2024} sharpen the \DIFdelbegin \DIFdel{complexity picturefurther}\DIFdelend \DIFaddbegin \DIFadd{parameterized picture}\DIFaddend ,
proving $1\|\sum w_jU_j$ is $W[1]$-hard parameterized by \DIFdelbegin \DIFdel{either }\DIFdelend the number of
distinct processing times \DIFdelbegin \DIFdel{or }\DIFdelend \DIFaddbegin \DIFadd{and also by }\DIFaddend the number of distinct weights\DIFdelbegin \DIFdel{in the instance---a result that, translated }\DIFdelend \DIFaddbegin \DIFadd{.
Translated }\DIFaddend into MWHED's language\DIFdelbegin \DIFdel{,
says that even a small }\DIFdelend \DIFaddbegin \DIFadd{: there is, unless $\mathrm{FPT}=W[1]$,
no algorithm running in $f(k)\cdot\mathrm{poly}(n)$ time where $k$ is
the }\DIFaddend number of distinct dispatch \DIFdelbegin \DIFdel{costs }\DIFdelend \DIFaddbegin \DIFadd{times (}\DIFaddend or criticality levels\DIFdelbegin \DIFdel{does not, by itself, guarantee tractability, sharpening
the boundary Theorem~\ref{thm:equalp} and Remark~\ref{rem:hardness-boundary}
draw between the fully-heterogeneous and equal-cost cases}\DIFdelend \DIFaddbegin \DIFadd{), so a
small number of distinct values does not make the problem easy in the
fixed-parameter sense. For a fixed $k$ the problem remains polynomial
(a dynamic program over the counts of selected sites of each type has
$n^{O(k)}$ states) and for $k=1$ it is the tractable case of
Theorem~\ref{thm:equalp} (or Moore--Hodgson, for equal weights), so
these results are consistent with, and refine, the boundary that
Theorem~\ref{thm:classification} draws}\DIFaddend . On the
heuristic side, \citet{antonov2026} give a data-driven (machine-learning-guided)
heuristic for the same objective, a more sophisticated relative of the
weighted greedy-repair heuristic we use as an empirical baseline in
Section~\ref{sec:experiments}.

\textbf{Real-time and deadline-driven systems.} Earliest-deadline-first
ordering is also the central object of classical real-time scheduling
theory, where \citet{liulayland1973} show that, for periodic tasks on
a single processor, earliest-deadline-first achieves full processor
utilization while any fixed-priority policy (such as rate-monotonic
scheduling) cannot. That literature's EDF result and our
Lemma~\ref{lem:edd} are structurally analogous---both establish that
ordering by deadline is necessary and sufficient for a natural notion
of feasibility---but the objectives differ fundamentally: real-time
scheduling asks whether \emph{every} task can be met, a feasibility
question with a clean necessary-and-sufficient utilization test,
whereas MWHED presupposes infeasibility is unavoidable when dispatch
costs are large relative to deadlines, and asks instead which
\emph{weighted subset} of sites to sacrifice---the optimization
question that drives Theorem~\ref{thm:hardness}'s hardness result and
has no analogue in the feasibility-testing literature. It is worth
noting why the two literatures' core feasibility results still agree
despite this difference in objective: \citeauthor{liulayland1973}'s
utilization test is a necessary-and-sufficient \emph{aggregate}
condition (a single inequality over all tasks) precisely because their
tasks recur periodically forever, so any transient slack evens out over
a long enough horizon; MWHED's sites are each served at most once, so
no such aggregate test can substitute for Lemma~\ref{lem:edd}'s
site-by-site exchange argument, and no analogous single-inequality
feasibility test exists for MWHED's one-shot dispatch problem---indeed,
Theorem~\ref{thm:hardness} shows that even deciding the size of the
best \emph{weighted} feasible subset, let alone a simple aggregate
test for it, is NP-hard. More recent real-time scheduling work extends
EDF-style analysis to richer task models: \citet{gunzel2022} give a
unified schedulability test for self-suspending periodic tasks that
recovers EDF (among other policies) as a special case, and
\citet{wang2025} refine the schedulability analysis for that same
self-suspending, EDF-like setting. Both retain the periodic,
run-forever structure that makes an aggregate utilization test
possible; MWHED's one-shot, single-visit-per-site structure is exactly
what rules out an analogous test here.

\textbf{Approximation.} The 0/1 knapsack problem admits a fully
polynomial-time approximation scheme via scaling and rounding of item
values, due to \citet{ibarrakim1975}; \DIFdelbegin \DIFdel{we adapt this technique directly
to MWHED's dynamic program in }\DIFdelend \DIFaddbegin \DIFadd{the same technique applies to job
sequencing with deadlines \mbox{%DIFAUXCMD
\citep{sahni1976,gensl1981}}\hskip0pt%DIFAUXCMD
, and
}\DIFaddend Theorem~\ref{thm:fptas} \DIFdelbegin \DIFdel{. The adaptation
is not entirely mechanical: knapsack 's FPTAS scales values and reuses
a single time-indexed capacity constraint, whereas MWHED has }\DIFdelend \DIFaddbegin \DIFadd{is this construction written for MWHED. The one
point at which the presentation differs from knapsack is that there is
}\DIFaddend no single shared \DIFdelbegin \DIFdel{capacity---each site carries its own deadline---so }\DIFdelend \DIFaddbegin \DIFadd{capacity (each site has its own deadline), so }\DIFaddend the
value-indexed \DIFdelbegin \DIFdel{dual table of Theorem~\ref{thm:fptas}'s proof must track
}\DIFdelend \DIFaddbegin \DIFadd{table must test }\DIFaddend feasibility against each site's \DIFdelbegin \DIFdel{individual }\DIFdelend $d_i$\DIFdelbegin \DIFdel{rather than a single
global bound}\DIFdelend \DIFaddbegin \DIFadd{;
this is easy in deadline order}\DIFaddend , and the case split \DIFdelbegin \DIFdel{we make explicit }\DIFdelend in
Appendix~\ref{app:fptas-detail} (\DIFdelbegin \DIFdel{between the scaling being vacuousor
not) has no counterpart in the standard knapsack presentation, where
$w_{\max}$ and $n$ play a comparatively simpler role}\DIFdelend \DIFaddbegin \DIFadd{whether the scaling is vacuous) is
simply careful bookkeeping, not a new idea}\DIFaddend . The classical
$O(n^3/\epsilon)$-style \DIFdelbegin \DIFdel{FPTAS we adapt }\DIFdelend \DIFaddbegin \DIFadd{scheme }\DIFaddend is not the fastest \DIFdelbegin \DIFdel{knapsack FPTAS
known today: \mbox{%DIFAUXCMD
\citet{jin2019} }\hskip0pt%DIFAUXCMD
improve the running time via a more
careful scaling scheme, }\DIFdelend \DIFaddbegin \DIFadd{known even for
knapsack: \mbox{%DIFAUXCMD
\citet{jin2019} }\hskip0pt%DIFAUXCMD
}\DIFaddend and \citet{chen2025fptas} give \DIFdelbegin \DIFdel{a nearly
quadratic-time FPTAS, $\tilde O(n + 1/\epsilon^2)$, that is conditionally
near-optimal. We adapt the classical (rather than the fastest known)
}\DIFdelend \DIFaddbegin \DIFadd{faster
schemes, and \mbox{%DIFAUXCMD
\citet{gensl1981} }\hskip0pt%DIFAUXCMD
give a faster scheme for job sequencing
with deadlines itself. We keep the classical }\DIFaddend construction because its
\DIFdelbegin \DIFdel{case structure is what Theorem~\ref{thm:fptas}'s
proof and Appendix~\ref{app:fptas-detail}'s derivation build on
directly; porting the faster constructions to MWHED's per-site-deadline
setting (rather than knapsack's single shared capacity) is a natural
avenue for tightening Theorem~\ref{thm:fptas}'s runtime bound}\DIFdelend \DIFaddbegin \DIFadd{analysis is the cleanest and because Proposition~\ref{prop:tight} shows
that analysis is tight}\DIFaddend .

\textbf{Matroids and greedy algorithms.} The equal-processing-time
special case of Theorem~\ref{thm:equalp} is an instance of the
classical problem of scheduling unit-time jobs with deadlines to
maximize total profit, whose feasible job sets form a matroid, on which
a greedy algorithm is provably optimal; we follow the treatment in
\citet{kortevygen2018}. This matroid structure is exactly what breaks
once dispatch costs become heterogeneous, which is why
Theorem~\ref{thm:equalp}'s greedy algorithm has no counterpart for
general MWHED. Concretely, in the equal-deadline restriction (which
Theorem~\ref{thm:hardness} identifies with 0/1 knapsack of capacity
$D$), take three sites with dispatch times $2,1,5$ and a common
deadline $D=5$: $\{3\}$ (the size-$5$ site alone) is feasible and
already uses the entire capacity, while $\{1,2\}$ (total size $3$) is
also feasible; but neither remaining site can be added to $\{3\}$ ($5+2$
and $5+1$ both exceed $D=5$), so a feasible set of size $1$ fails to
extend toward a feasible set of size $2$ by any element of the latter---precisely
the matroid exchange property failing. This is why a single large,
low-value commitment can block a strictly better pair of smaller
commitments, and is the structural reason Theorem~\ref{thm:hardness}'s
hardness survives even though the matroid case (Theorem~\ref{thm:equalp})
is easy. Matroid algorithms themselves remain an active area:
\citet{terao2025} give faster matroid partition algorithms, and
\citet{ma2026} prove a greedy algorithm is simultaneously optimal and
fair for a multiagent matroid-upgrading problem, reinforcing that
``greedy is optimal on a matroid'' continues to be a productive
template well beyond the classical setting Theorem~\ref{thm:equalp}
uses.

\textbf{Routing and dispatch under a spreading hazard.} Wildfire
suppression is the application area where the ``race a spreading
hazard'' structure has been studied most directly in the operations
research literature, and it illustrates the point at which our
abstraction departs from that literature's usual level of detail.
\citet{harris2023} model fire spread via a minimum-travel-time
principle---the fire's arrival time at any point is governed by the
fastest path a flame front can take through spatially varying
fuel and terrain, akin to a shortest-path computation over a
continuous landscape---and decompose the resulting joint
suppression-and-routing decision via logic-based Benders decomposition,
trading exact fire-spread geometry against tractable large-scale
optimization. \citet{avci2024} extend this style of problem to multiple
heterogeneous resource types (e.g., ground crews and aircraft with
different suppression capabilities and travel speeds), which multiplies
the routing decision's combinatorial complexity without changing its
underlying deadline structure. \citet{alvelos2025} go further, jointly
optimizing dispatching, the positioning of resources in anticipation of
the fire's advance, and their routing, in a single integrated model.
More recently, \citet{granda2025} give a mathematical programming
formulation for wildfire suppression resource allocation, and, closest
in spirit to our own complexity results, \citet{delazeri2026} prove
NP-completeness for a landscape-graph model of wildfire suppression
resource allocation, alongside a competitive mixed-integer program and
a physics-grounded (Rothermel fire-spread model) instance generator---an
independent hardness result in essentially our application domain,
though for a differently structured (graph-based, simultaneous
multi-resource) model than MWHED's single-vehicle dispatch-ordering
question. Two recent systematic reviews position the broader
literature: \citet{tezcan2025} survey optimization approaches to
forest-fire suppression, and \citet{rostamian2026} review 177 papers on
wildfire management optimization under uncertainty and complexity,
both confirming that pure complexity-theoretic characterizations of a
dispatch-ordering subproblem---as opposed to joint suppression,
positioning, and routing models---remain rare in this literature.
Across all of these, the fire's own spatial spread is modeled in detail
and coupled tightly to the routing and positioning decisions; MWHED
instead abstracts the hazard's spread into a per-site deadline (a
natural reduction once the hazard's spread model is fixed and known, as
is standard once a fire, flood, or contamination-plume simulator has
produced arrival-time estimates for each site) and asks the sharper
question of what
complexity and algorithmic guarantees are available for the resulting
pure dispatch-\emph{ordering} decision. To our knowledge, no prior work
in this literature establishes the knapsack-equivalence hardness
result, the exact dynamic program, the FPTAS, or the equal-cost
tractable case we give here.

\textbf{Positioning.} MWHED \DIFdelbegin \DIFdel{sits at the intersection of these four
literatures without duplicating any of them: it }\DIFdelend inherits its complexity classification from
\DIFdelbegin \DIFdel{classical scheduling theory
, both the founding
results and their recent refinements
(\mbox{%DIFAUXCMD
\citealp{moore1968,lawlermoore1969,lenstra1977,hejl2022,hermelin2024,heegerhermelin2024,antonov2026}}\hskip0pt%DIFAUXCMD
}\DIFdelend \DIFaddbegin \DIFadd{scheduling theory
(\mbox{%DIFAUXCMD
\citealp{moore1968,lawlermoore1969,lenstra1977,sahni1976,gensl1981,hejl2022,hermelin2024,heegerhermelin2024,antonov2026}}\hskip0pt%DIFAUXCMD
}\DIFaddend ),
its algorithmic toolkit from knapsack approximation and matroid theory
\DIFdelbegin \DIFdel{,
including recent speedups
}\DIFdelend (\citealp{ibarrakim1975,kortevygen2018,jin2019,chen2025fptas,terao2025,ma2026}),
its deadline-ordering intuition from real-time systems
(\citealp{liulayland1973,gunzel2022,wang2025}), and its \DIFdelbegin \DIFdel{motivating
application }\DIFdelend \DIFaddbegin \DIFadd{motivation }\DIFaddend from
disaster-response operations research
(\citealp{harris2023,avci2024,alvelos2025,granda2025,delazeri2026,tezcan2025,rostamian2026})\DIFdelbegin \DIFdel{,
but the specific problem --- and every result we prove about it --- is
new. }\DIFdelend \DIFaddbegin \DIFadd{.
Its theoretical results are, with the exceptions listed above, classical
ones restated in transportation language; }\DIFaddend Table~\ref{tab:summary} in
Section~\ref{sec:theory} \DIFdelbegin \DIFdel{summarizes how each of our five main
results maps onto this classification. The overall }\DIFdelend \DIFaddbegin \DIFadd{marks which. We believe the }\DIFaddend pattern---a
transportation reparametrization of a classical scheduling problem \DIFdelbegin \DIFdel{,
whose
hardness, exact algorithm, and approximation scheme all follow
by directly adapting standard techniques, while a natural structural
restriction (equal dispatch cost) recovers tractability---is, we
believe, a template that could apply to other
transportation-network
problems with a similar }\DIFdelend \DIFaddbegin \DIFadd{whose
tractability boundary has a network interpretation---applies to other
}\DIFaddend deadline-driven \DIFdelbegin \DIFdel{structure, not only to the
specific hazard-dispatch motivation we use here}\DIFdelend \DIFaddbegin \DIFadd{transportation problems}\DIFaddend ; Section~\ref{sec:discussion}
discusses \DIFdelbegin \DIFdel{three such extensionsconcretely}\DIFdelend \DIFaddbegin \DIFadd{extensions}\DIFaddend .

% ══════════════════════════════════════════════════════════════════
\section{Model}\label{sec:model}

\begin{definition}[Minimum Weighted Hazard-Exposure Dispatch]\label{def:mwhed}
An instance of MWHED consists of $n$ sites $1,\dots,n$ and a depot. Each
site $i$ has a round-trip dispatch time $p_i \in \mathbb Z_{>0}$ (the
time to travel from the depot to site $i$ and back, or more generally
the time the depot's single vehicle or crew is occupied while serving
site $i$), a hazard-arrival deadline $d_i \in \mathbb Z_{>0}$ (the known
time at which a spreading hazard reaches site $i$), and a criticality
weight $w_i \in \mathbb Z_{>0}$. A \emph{dispatch order} is a
permutation $\sigma$ of $\{1,\dots,n\}$; the \emph{completion time} of
site $\sigma(k)$ is $C_{\sigma(k)}(\sigma) := \sum_{j=1}^k
p_{\sigma(j)}$, and site $i$ is \emph{on time} under $\sigma$ if $C_i(\sigma)
\le d_i$, otherwise \emph{exposed}. The objective is to choose $\sigma$
maximizing the total on-time weight
\begin{equation}\label{eq:objective}
W(\sigma) := \sum_{i=1}^n w_i \cdot \mathbb 1\bigl[C_i(\sigma) \le
d_i\bigr].
\end{equation}
\end{definition}

Table~\ref{tab:notation} collects the notation introduced in this
definition and used throughout the rest of the paper, for reference.

\begin{table}[h]
\caption{Notation used throughout the paper.}\label{tab:notation}
\begin{tabular}{@{}ll@{}}
\toprule
Symbol & Meaning \\
\midrule
$n$ & number of sites \\
$p_i$ & round-trip dispatch time for site $i$ \\
$d_i$ & hazard-arrival deadline for site $i$ \\
$w_i$ & criticality weight for site $i$ \\
$\sigma$ & a dispatch order (permutation of sites) \\
$C_i(\sigma)$ & completion time of site $i$ under order $\sigma$ \\
$W(\sigma)$ & total on-time weight of order $\sigma$, Eq.~\eqref{eq:objective} \\
$W^\ast$ & the optimal (maximum) value of $W(\sigma)$ \\
$S$ & a subset of sites dispatched on time \\
$P$ & $\sum_{i=1}^n p_i$, the total dispatch time (Theorem~\ref{thm:dp}) \\
$f(i,t)$ & exact DP table: max weight using sites $1..i$, time exactly $t$ \\
$\epsilon$ & FPTAS accuracy parameter, $\epsilon\in(0,1)$ (Theorem~\ref{thm:fptas}) \\
$K$ & FPTAS weight-scaling factor \\
$w_i'$ & scaled weight, $\lfloor w_i/K\rfloor$ \\
$g(i,v)$ & FPTAS's value-indexed dual DP table \\
$D_i$ & latest feasible dispatch position, equal-cost case, $\lfloor d_i/p\rfloor$ \\
\DIFaddbeginFL \DIFaddFL{$m$ }& \DIFaddFL{number of identical vehicles (MWHED-$m$, Section~\ref{sec:equalp}) }\\
\DIFaddendFL \bottomrule
\end{tabular}
\end{table}

Because every visit returns to the depot before the next dispatch
begins, $p_i$ does not depend on which site precedes or follows site
$i$ in the order---the defining structural feature that distinguishes
MWHED from a routing problem and places it in the scheduling family.
Minimizing total \emph{exposure cost} $\sum_i w_i(1 - \mathbb
1[C_i(\sigma)\le d_i])$ is equivalent to maximizing $W(\sigma)$, since
$\sum_i w_i$ is a constant of the instance; we work with the
maximization form throughout.

\begin{assumption}[Individually feasible sites]\label{as:feasible}
Every site satisfies $p_i \le d_i$.
\end{assumption}

This is without loss of generality: a site with $p_i > d_i$ cannot be
on time under any dispatch order, even served first, so it can be
deleted from the instance in a preprocessing pass without changing the
optimal value achievable on the remaining sites. One immediate
consequence we use later: the single highest-weight site, dispatched
alone, is itself a feasible subset, so $W^\ast \ge w_{\max} := \max_i
w_i$.

\DIFaddbegin \begin{remark}[Return and arrival readings of the deadline]\label{rem:reading}
\DIFadd{Definition~\ref{def:mwhed} counts a site as on time when the vehicle's
}\emph{\DIFadd{completion}} \DIFadd{of the whole $p_i$-long occupancy---the round trip,
including the return---happens by $d_i$ (the }\emph{\DIFadd{return reading}}\DIFadd{).
When what matters is that the vehicle }\emph{\DIFadd{arrive}} \DIFadd{at the site before
the hazard does, only the outbound part $a_i\le p_i$ of the occupancy
must fit: site $i$ is protected iff $C_i-p_i+a_i\le d_i^{\mathrm{haz}}$,
where $d_i^{\mathrm{haz}}$ is the hazard-arrival time. This is the same
test with deadline $d_i:=d_i^{\mathrm{haz}}+(p_i-a_i)$, which for a
symmetric trip is $d_i^{\mathrm{haz}}+p_i/2$. Hence the }\emph{\DIFadd{arrival
reading}} \DIFadd{is an instance of Definition~\ref{def:mwhed} with shifted
deadlines, and every result below applies to it unchanged; only the
input $d_i$ differs. The two readings can give different answers on the
same data, so a practitioner must state which one the hazard-arrival
times refer to; Section~\ref{sec:campfire} uses the arrival reading.
}\end{remark}

\DIFaddend The following lemma is the structural fact both the dynamic program of
Section~\ref{sec:theory} and the tractability results depend on. It is
the transportation-scheduling analogue of the classical fact that
earliest-due-date order minimizes maximum lateness on a single machine.

\begin{lemma}[Earliest-deadline-first feasibility]\label{lem:edd}
Fix a subset $S \subseteq \{1,\dots,n\}$. There exists a dispatch order
under which every site in $S$ is on time if and only if, when the sites
of $S$ are dispatched first (before any site outside $S$) in order of
increasing deadline $d_i$, every site in $S$ is on time.
\end{lemma}

\begin{proof}
Sufficiency is immediate: the stated order is one such dispatch order.
For necessity, suppose $S$ is feasible under some order $\sigma$ (every
site of $S$ on time). We transform $\sigma$ into the stated order by two
kinds of exchanges, each of which preserves feasibility of $S$.

First, suppose some site $j \notin S$ precedes some site $i \in S$ in
$\sigma$. Moving $j$ to immediately after the last element of $S$
(and shifting the intervening sites earlier by $p_j$) strictly decreases
the completion time of every site of $S$ that came after $j$, and
leaves the completion time of every site of $S$ that came before $j$
unchanged; it therefore preserves feasibility of every site in $S$, and
after repeating this for every $j \notin S$ that precedes some element
of $S$, all of $S$ forms a prefix of the schedule.

Second, with $S$ now a prefix, suppose two adjacent sites $i, i' \in S$
appear in the order with $i$ immediately before $i'$, yet $d_i >
d_{i'}$---the later-deadline site is scheduled first. Let $t_0$ be the
cumulative dispatch time immediately before this pair is served
(unaffected by how the pair is internally ordered, since it depends
only on the sites preceding them). Under the current order, $i$
completes at $t_0+p_i$ and $i'$ completes at $t_0+p_i+p_{i'}$; since
the order is feasible for $S$, in particular $t_0+p_i+p_{i'} \le
d_{i'}$. Swap the pair so that $i'$ precedes $i$: every other site in
$S$ still finishes at the same time as before, $i'$ now completes at
$t_0+p_{i'}$, and $i$ now completes at $t_0+p_{i'}+p_i$---the same
total as before, since the pair's combined processing time is
unchanged by reordering it. Because $p_i>0$, $t_0+p_{i'} <
t_0+p_i+p_{i'} \le d_{i'}$, so $i'$ remains on time; and
$t_0+p_i+p_{i'} \le d_{i'} < d_i$, so $i$ remains on time as well.
Repeating this adjacent-transposition argument (a bubble sort on the
prefix) sorts $S$ into increasing-$d_i$ order without ever violating a
previously-satisfied deadline.
\end{proof}

Lemma~\ref{lem:edd} lets us restrict attention, without loss of
generality, to dispatch orders that serve some subset $S$ first, in
increasing-deadline order, followed by the remaining (necessarily
exposed) sites in any order; the optimization problem becomes that of
choosing $S$ to maximize $\sum_{i\in S} w_i$ subject to the feasibility
condition of Lemma~\ref{lem:edd}.

\begin{example}[A running illustration]\label{ex:running}
Throughout this section we illustrate each result on a single small
instance: four sites with round-trip dispatch times $p=(2,3,1,4)$,
hazard deadlines $d=(4,6,2,9)$, and criticality weights $w=(5,8,3,10)$,
all satisfying Assumption~\ref{as:feasible}. Sorting by increasing
deadline gives the order $3,1,2,4$ (deadlines $2,4,6,9$);
Lemma~\ref{lem:edd} says we may restrict to dispatch orders that serve
some subset $S$ first, in exactly this relative order.

Dispatching every site in this order---the naive policy of simply
visiting sites earliest-deadline-first, with no further
optimization---gives completion times $1,3,6,10$ for sites $3,1,2,4$
respectively; the first three are on time, but site $4$, with deadline
$9$, arrives at cumulative time $10$ and is exposed, losing its full
weight $w_4=10$. This naive order captures only $w_3+w_1+w_2 =
3+5+8=16$ of the total weight $26$.

The optimal choice of $S$ is instead $\{1,2,4\}$, excluding site $3$:
dispatched in deadline order $1,2,4$, completion times are $2,5,9$,
each on time, for a total on-time weight of $W^\ast = 5+8+10=23$---
strictly better than the naive $16$, by deliberately sacrificing site
$3$ (deadline $2$, weight only $3$) to make room for the far more
valuable site $4$. No dispatch order can do better, as
Example~\ref{ex:dp}'s dynamic program table confirms exhaustively. This
gap between naive earliest-deadline dispatch and the true optimum is
\DIFdelbegin \DIFdel{exactly the phenomenon }\DIFdelend \DIFaddbegin \DIFadd{what }\DIFaddend Section~\ref{sec:experiments} \DIFdelbegin \DIFdel{'s numerical study
documents at scale}\DIFdelend \DIFaddbegin \DIFadd{measures on random instances}\DIFaddend .
\end{example}

\begin{figure}[h]
\centering
\includegraphics[width=0.75\textwidth]{figures/running_example_schematic.pdf}
\caption{A schematic illustration of Example~\ref{ex:running}. Site
positions are laid out at a radius proportional to each site's
round-trip dispatch time $p_i$, purely for visual spacing---this is an
abstract schematic of the instance, not a real road network or a
geographic map (MWHED's theory is instance-agnostic and this paper
uses no real-world coordinate data). Marker size is proportional to
criticality weight $w_i$; green markers and arrows show the optimal
dispatch $S^\ast=\{1,2,4\}$ and each site's completion time; the gray
marker is site $3$, sacrificed to make room for the higher-value site
$4$.}\label{fig:running-example}
\end{figure}

% ══════════════════════════════════════════════════════════════════
\section{Complexity and algorithms}\label{sec:theory}

Table~\ref{tab:summary} previews \DIFdelbegin \DIFdel{how the five }\DIFdelend \DIFaddbegin \DIFadd{the }\DIFaddend results of this section \DIFdelbegin \DIFdel{fit together, before we state and prove each in turn}\DIFdelend \DIFaddbegin \DIFadd{and marks
which are classical (stated here in transportation language, with
complete proofs) and which we believe are new}\DIFaddend .

\begin{table}[h]
\caption{Summary of this section's results. $n$ is the number of
sites, $P=\sum_i p_i$, and $\epsilon\in(0,1)$ is the FPTAS's accuracy
parameter.}\label{tab:summary}
\small
\DIFdelbeginFL %DIFDELCMD < \begin{tabular}{@{}p{2.7cm}p{3.2cm}p{2.2cm}p{2.6cm}@{}}
%DIFDELCMD < %%%
\DIFdelendFL \DIFaddbeginFL \begin{tabular}{@{}p{2.55cm}p{2.5cm}p{1.45cm}p{1.5cm}p{1.75cm}@{}}
\DIFaddendFL \toprule
Setting & Status & Result & Runtime \DIFaddbeginFL & \DIFaddFL{Origin }\DIFaddendFL \\
\midrule
General \DIFdelbeginFL \DIFdelFL{instance }\DIFdelendFL & weakly NP-hard & Thm.~\ref{thm:hardness} & --- \DIFaddbeginFL & \DIFaddFL{classical }\DIFaddendFL \\
General \DIFdelbeginFL \DIFdelFL{instance }\DIFdelendFL & solved exactly & Thm.~\ref{thm:dp} & $O(nP)$ \DIFaddbeginFL & \DIFaddFL{Lawler--Moore }\DIFaddendFL \\
General \DIFdelbeginFL \DIFdelFL{instance }\DIFdelendFL & $(1-\epsilon)$\DIFdelbeginFL \DIFdelFL{-approximated }\DIFdelendFL \DIFaddbeginFL \DIFaddFL{-approx. }\DIFaddendFL & Thm.~\ref{thm:fptas} & $O(n^3/\epsilon)$ \DIFaddbeginFL & \DIFaddFL{classical type }\DIFaddendFL \\
Equal deadlines \DIFdelbeginFL \DIFdelFL{($d_i{=}D$) }\DIFdelendFL & weakly NP-hard ($=$ knapsack) & Thm.~\ref{thm:hardness} & --- \DIFaddbeginFL & \DIFaddFL{classical }\DIFaddendFL \\
Equal \DIFdelbeginFL \DIFdelFL{cost ($p_i{=}p$}\DIFdelendFL \DIFaddbeginFL \DIFaddFL{weights }& \DIFaddFL{solved exactly }& \DIFaddFL{Thm.~\ref{thm:classification}(a}\DIFaddendFL ) & \DIFaddbeginFL \DIFaddFL{$O(n\log n)$ }& \DIFaddFL{Moore }\\
\DIFaddFL{Equal cost, $m=1$ }& \DIFaddendFL solved exactly & Thm.~\ref{thm:equalp} & $O(n\log n)$ \DIFaddbeginFL & \DIFaddFL{classical }\\
\DIFaddFL{Equal cost, $m\ge2$ identical vehicles }& \DIFaddFL{solved exactly }& \DIFaddFL{Thm.~\ref{thm:equalp} }& \DIFaddFL{$O(n\log n)$ }& \DIFaddFL{new (proof) }\\
\DIFaddFL{Which of $p,w,d$ are constant }& \DIFaddFL{polynomial iff $p$ or $w$ constant }& \DIFaddFL{Thm.~\ref{thm:classification} }& \DIFaddFL{--- }& \DIFaddFL{new framing }\DIFaddendFL \\
Naive EDD, \DIFdelbeginFL \DIFdelFL{no repair }\DIFdelendFL \DIFaddbeginFL \DIFaddFL{EDD with skipping }\DIFaddendFL & unbounded \DIFdelbeginFL \DIFdelFL{worst-case }\DIFdelendFL ratio & Prop.~\ref{prop:naive} & $O(n\log n)$ \DIFaddbeginFL & \DIFaddFL{new }\\
\DIFaddFL{Weighted greedy repair }& \DIFaddFL{unbounded ratio }& \DIFaddFL{Prop.~\ref{prop:repair} }& \DIFaddFL{--- }& \DIFaddFL{new }\\
\DIFaddFL{FPTAS analysis }& \DIFaddFL{tight to $1-\epsilon^2$ }& \DIFaddFL{Prop.~\ref{prop:tight} }& \DIFaddFL{--- }& \DIFaddFL{new }\DIFaddendFL \\
\bottomrule
\end{tabular}
\end{table}

\subsection{Hardness}

\DIFaddbegin \DIFadd{We first isolate the one fact about completion times on which the
hardness argument rests. Recall that, by the defining property of MWHED,
the occupancy $p_i$ of the vehicle at site $i$ does not depend on which
site was served before it.
}

\begin{fact}[Feasibility under a common deadline]\label{lem:common}
\DIFadd{Suppose $d_i=D$ for every site, and let $S\subseteq\{1,\dots,n\}$. There
is a dispatch order under which every site of $S$ is on time if and only
if $\sum_{i\in S}p_i\le D$.
}\end{fact}

\begin{proof}
\DIFadd{By Lemma~\ref{lem:edd} we may serve the sites of $S$ first, in any order
(all deadlines are equal, so the earliest-deadline-first order is
arbitrary among them). The vehicle is never idle and the occupancies
are fixed numbers, so the $k$-th site served completes at the partial
sum of the first $k$ occupancies, and the }\emph{\DIFadd{last}} \DIFadd{site of $S$
completes at the sum $T(S):=\sum_{i\in S}p_i$ of }\emph{\DIFadd{all}} \DIFadd{of their
occupancies, whatever the order. This is the only additivity used: it is
a statement about the completion time of the last site, not about every
site's completion time being additive.
}

\DIFadd{($\Rightarrow$) If every site of $S$ is on time, then in particular the
last one is, so $T(S)\le D$. ($\Leftarrow$) If $T(S)\le D$, then every
site of $S$ completes at a partial sum of positive numbers adding up to
$T(S)$, hence at some time $\le T(S)\le D$, so every site is on time.
}\end{proof}

\DIFaddend \begin{theorem}[NP-hardness of MWHED]\label{thm:hardness}
MWHED is NP-hard. This holds even \DIFdelbegin \DIFdel{restricted to instances in which
}\DIFdelend \DIFaddbegin \DIFadd{when }\DIFaddend every site shares a common
hazard deadline, \DIFdelbegin \DIFdel{$d_i = D$ }\DIFdelend \DIFaddbegin \DIFadd{$d_i=D$ }\DIFaddend for all $i$\DIFaddbegin \DIFadd{, and even when $w_i=p_i$ for all
$i$. Precisely, the decision problem ``is $W^\ast\ge k$?'' is
NP-complete}\DIFaddend .
\end{theorem}

\begin{proof}
\DIFdelbegin \DIFdel{Consider the restriction $d_i = D$ }\DIFdelend \DIFaddbegin \DIFadd{Membership in NP is clear: a dispatch order is a certificate and
$W(\sigma)$ is computed in $O(n)$ time. For hardness we reduce from
}\textsc{\DIFadd{Partition}} \DIFadd{\mbox{%DIFAUXCMD
\citep{gareyjohnson1979}}\hskip0pt%DIFAUXCMD
: given positive integers
$a_1,\dots,a_n$ with $A:=\sum_ia_i$ even, decide whether some
$S\subseteq\{1,\dots,n\}$ has $\sum_{i\in S}a_i=A/2$. If some $a_i>A/2$
the answer is no and we output a fixed no-instance, so assume
$a_i\le A/2$ }\DIFaddend for all $i$. \DIFdelbegin \DIFdel{By
Lemma~\ref{lem:edd}, a subset }\DIFdelend \DIFaddbegin \DIFadd{Build the MWHED instance with
}\[
\DIFadd{p_i=w_i=a_i,\qquad d_i=D:=A/2\qquad (i=1,\dots,n),
}\]
\DIFadd{which satisfies Assumption~\ref{as:feasible} and is built in polynomial
time. By Fact~\ref{lem:common} a set }\DIFaddend $S$ \DIFdelbegin \DIFdel{is feasible }\DIFdelend \DIFaddbegin \DIFadd{can be served entirely on
time iff $\sum_{i\in S}a_i\le A/2$, and then its on-time weight is
$\sum_{i\in S}w_i=\sum_{i\in S}a_i$. Since only the sites of a feasible
set can be on time, $W^\ast=\max\{\sum_{i\in S}a_i:\ \sum_{i\in S}a_i\le
A/2\}\le A/2$, with equality }\DIFaddend if and only if \DIFdelbegin \DIFdel{, dispatched
in any order (deadlines being tied, the earliest-deadline-first
condition is
vacuous), every site in $S$ completes by $D$; since
completion times along a fixed set are additive, this holds exactly
when $\sum_{i\in S} p_i \le D$.
The optimization problem
}\[
\DIFdel{\max_{S \subseteq \{1,\dots,n\}} \ \sum_{i\in S} w_i \quad \text{s.t.} \quad
\sum_{i \in S} p_i \le D
}\]%DIFAUXCMD
\DIFdel{is precisely }\DIFdelend \DIFaddbegin \DIFadd{some subset sums to exactly
$A/2$. Thus $W^\ast\ge A/2$ holds iff the }\textsc{\DIFadd{Partition}} \DIFadd{instance is
a yes-instance.
}\end{proof}

\DIFadd{With general weights the common-deadline restriction is exactly }\DIFaddend the 0/1
knapsack problem \DIFdelbegin \DIFdel{, with $p_i$ the item sizes , $w_i$
the item values , and $D$ the capacity . The (decision version of the) 0/1
knapsack problem is NP-complete \mbox{%DIFAUXCMD
\citep{gareyjohnson1979}}\hskip0pt%DIFAUXCMD
, by
a
polynomial reduction from PARTITION. Hence MWHED restricted to a common
deadline is NP-hard, and since this is a special case of the
general
problem, MWHED itself is NP-hard. }%DIFDELCMD < \end{proof}
%DIFDELCMD < 

%DIFDELCMD < %%%
\DIFdel{Theorem~\ref{thm:hardness}'s reduction is deliberately the simplest one
available}\DIFdelend \DIFaddbegin \DIFadd{with sizes $p_i$, values $w_i$ and capacity $D$, by
Fact~\ref{lem:common}; this is why the equal-deadline case is the
natural source of hardness, and why the dynamic program of the next
subsection (which is pseudo-polynomial) shows that the hardness is only
}\emph{\DIFadd{weak}}\DIFadd{. The reduction deliberately uses the simplest possible
structure}\DIFaddend : it isolates \DIFdelbegin \emph{\DIFdel{exactly}} %DIFAUXCMD
\DIFdel{the source of hardness, namely
}\DIFdelend heterogeneous dispatch costs weighed against
heterogeneous criticality under a shared time budget \DIFdelbegin \DIFdel{, and reduces the question of
what other
structure might rescue tractability to what Theorem~\ref{thm:equalp}
answers}\DIFdelend \DIFaddbegin \DIFadd{as the source of
difficulty, and Theorem~\ref{thm:classification} shows that this is the
only one}\DIFaddend .

\begin{example}[The running example under a common deadline]\label{ex:knapsack}
Suppose the four sites of Example~\ref{ex:running} were instead all
assigned the same hazard deadline $D=6$ (same dispatch times and
weights). By \DIFdelbegin \DIFdel{Theorem~\ref{thm:hardness}'s reduction }\DIFdelend \DIFaddbegin \DIFadd{Fact~\ref{lem:common} }\DIFaddend this is exactly a $0/1$ knapsack
instance with capacity $6$, item sizes $p=(2,3,1,4)$, and item values
$w=(5,8,3,10)$. Checking every subset whose total dispatch time is at
most $6$, the best is $\{1,2,3\}$, costing $2+3+1=6$ and worth
$5+8+3=16$; keeping the single most valuable site instead does worse
within the same budget---$\{1,4\}$ costs $6$ but is worth only $15$, and
$\{3,4\}$ costs $5$ but is worth only $13$. Which sites to sacrifice to
make room, by weight against size, cannot be read off from either
quantity alone; this trade-off is exactly the source of MWHED's
hardness.
\end{example}

\subsection{An exact pseudo-polynomial algorithm}

\begin{theorem}[Exact dynamic program]\label{thm:dp}
MWHED can be solved exactly in $O(n \cdot P)$ time and space, where
$P := \sum_{i=1}^n p_i$.
\end{theorem}

\begin{proof}
Relabel sites so that $d_1 \le d_2 \le \dots \le d_n$ (Lemma~\ref{lem:edd}
justifies working in this order). Define, for $i = 0,\dots,n$ and $t =
0,\dots,P$,
\[
f(i,t) := \text{the maximum value of } \sum_{j \in S} w_j \text{ over }
S \subseteq \{1,\dots,i\} \text{ with } \sum_{j\in S} p_j = t \text{ and
$S$ feasible.}
\]
(feasible meaning: dispatching $S$ in increasing-deadline order, every
member is on time). Set $f(0,0) = 0$ and $f(0,t) = -\infty$ for $t > 0$.
For $i \ge 1$,
\begin{equation}\label{eq:dp}
f(i,t) = \max\Bigl(f(i-1,t),\ \bigl[t \le d_i\bigr]\cdot\bigl(f(i-1,t-p_i)
+ w_i\bigr)\Bigr),
\end{equation}
where the second term is taken as $-\infty$ if $t < p_i$ or $t > d_i$.

We verify \eqref{eq:dp} by induction on $i$. The first term corresponds
to excluding site $i$ from $S$. For the second term, suppose $S'
\subseteq \{1,\dots,i-1\}$ achieves $f(i-1, t-p_i)$ and is feasible;
appending site $i$ after all of $S'$ (site $i$ has the largest deadline
among $\{1,\dots,i\}$, so this is consistent with increasing-deadline
order) does not change the completion time of any member of $S'$, so
$S'$'s feasibility is preserved, and site $i$'s own completion time is
exactly $t = (t - p_i) + p_i$, on time precisely when $t \le d_i$.
Conversely, any feasible $S \subseteq \{1,\dots,i\}$ with $i \in S$ and
total time $t$ restricts, by the same argument, to a feasible $S
\setminus \{i\} \subseteq \{1,\dots,i-1\}$ with total time $t - p_i$
achieving at least $f(i,t) - w_i$ at index $i-1$; hence $f(i-1,t-p_i)
\ge f(i,t) - w_i$, giving the reverse inequality and confirming
\eqref{eq:dp} computes exactly $f(i,t)$.

The optimal value is $\max_{t=0}^{P} f(n,t)$, and the achieving subset
is recovered by standard DP backtracking. The table has $O(n\cdot P)$
entries, each computed in $O(1)$ time from~\eqref{eq:dp}, giving
$O(nP)$ total time and space.
\end{proof}

Theorem~\ref{thm:dp} is pseudo-polynomial, not polynomial: $P$ can be
exponential in the input's bit length. This is consistent with
Theorem~\ref{thm:hardness}\DIFdelbegin \DIFdel{'s reduction from knapsack, which is itself
}\DIFdelend \DIFaddbegin \DIFadd{, whose hard instances need large numbers
(the problem is }\DIFaddend only weakly NP-hard\DIFaddbegin \DIFadd{)}\DIFaddend .

\DIFaddbegin \begin{remark}[Relation to Lawler--Moore]\label{rem:lawlermoore}
\DIFadd{This is the Lawler--Moore dynamic program \mbox{%DIFAUXCMD
\citep{lawlermoore1969} }\hskip0pt%DIFAUXCMD
for
$1\|\sum w_jU_j$: jobs in due-date order, state equal to the total
processing time of the jobs accepted so far, and the single test
$t\le d_i$ for an accepted job. We do not claim it as new. What the
proof above adds is a self-contained derivation from the
earliest-deadline-first lemma in the notation of the dispatch problem.
Refinements of the $O(nP)$ bound exist (see
Section~\ref{sec:related}); they apply verbatim to MWHED.
}\end{remark}

\DIFaddend \begin{example}[The dynamic program on the running instance]\label{ex:dp}
Table~\ref{tab:dpwalkthrough} traces $f(i,t)$ for
Example~\ref{ex:running}, relabeled in increasing-deadline order
$3,1,2,4$ so that index $i=1$ is site $3$, $i=2$ is site $1$, $i=3$ is
site $2$, and $i=4$ is site $4$, matching the order Lemma~\ref{lem:edd}
fixes. Reading the last row, the largest finite entry is $f(4,9)=23$;
backtracking \eqref{eq:dp} from this cell recovers exactly $S^\ast =
\{1,2,4\}$ in the original site labels, matching
Example~\ref{ex:running}. A dash marks \DIFdelbegin \DIFdel{$t > d_i$ for the site just
added at row $i$: once the cumulative time needed to include that site would already exceed }\DIFdelend \DIFaddbegin \DIFadd{$-\infty$: either the time $t$ cannot be reached by any
feasible subset of the first $i$ sites (for instance $f(1,2)$: site $3$
alone costs $1$), or including site $i$ at that column would violate }\DIFaddend its
own deadline \DIFdelbegin \DIFdel{, }\DIFdelend \DIFaddbegin \DIFadd{$t\le d_i$ (for instance $f(4,10)$, since $d_4=9$); in the
latter case }\DIFaddend the recursion can only carry the previous row's value
forward\DIFdelbegin \DIFdel{, never improve on it, at that column
and every column beyond it in the same row}\DIFdelend .
\end{example}

\begin{table}[h]
\caption{The dynamic program's table $f(i,t)$ for
Example~\ref{ex:running} ($i=1$: site 3, $p=1,d=2,w=3$; $i=2$: site 1,
$p=2,d=4,w=5$; $i=3$: site 2, $p=3,d=6,w=8$; $i=4$: site 4, $p=4,d=9,
w=10$). A dash marks $-\infty$ (infeasible or unreachable at that
cell).}\label{tab:dpwalkthrough}
\begin{tabular}{@{}lccccccccccc@{}}
\toprule
$t$ & 0 & 1 & 2 & 3 & 4 & 5 & 6 & 7 & 8 & 9 & 10\\
\midrule
$f(0,t)$ & 0 & -- & -- & -- & -- & -- & -- & -- & -- & -- & --\\
$f(1,t)$ & 0 & 3 & -- & -- & -- & -- & -- & -- & -- & -- & --\\
$f(2,t)$ & 0 & 3 & 5 & 8 & -- & -- & -- & -- & -- & -- & --\\
$f(3,t)$ & 0 & 3 & 5 & 8 & 11 & 13 & 16 & -- & -- & -- & --\\
$f(4,t)$ & 0 & 3 & 5 & 8 & 11 & 13 & 16 & 18 & 21 & \textbf{23} & --\\
\bottomrule
\end{tabular}
\end{table}

Algorithm~\ref{alg:exact} states Theorem~\ref{thm:dp}'s construction
explicitly, including the backtracking step used to recover an
achieving subset $S^\ast$, not merely its value.

\begin{algorithm}
\caption{Exact dynamic program for MWHED (Theorem~\ref{thm:dp})}\label{alg:exact}
\begin{algorithmic}[1]
\Require sites \DIFdelbegin \DIFdel{$1,\dots,n$ with $d_1 \le d_2 \le \dots \le d_n$ (relabeled per Lemma~\ref{lem:edd}), }\DIFdelend \DIFaddbegin \DIFadd{with }\DIFaddend dispatch times $p_i$, \DIFaddbegin \DIFadd{deadlines $d_i$, }\DIFaddend weights $w_i$
\Ensure optimal on-time weight $W^\ast$ and an achieving subset $S^\ast$
\State \DIFaddbegin \DIFadd{delete every site with $p_i > d_i$ (Assumption~\ref{as:feasible}); relabel the rest so $d_1 \le \dots \le d_n$ (Lemma~\ref{lem:edd})
}\State \DIFaddend $P \gets \sum_{i=1}^n p_i$
\State $f[0][0] \gets 0$; \ $f[0][t] \gets -\infty$ for $t=1,\dots,P$
\For{$i = 1,\dots,n$}
  \For{$t = 0,\dots,P$}
    \State $f[i][t] \gets f[i-1][t]$ \Comment{exclude site $i$}
    \If{$t \ge p_i$ \textbf{and} $t \le d_i$ \textbf{and} $f[i-1][t-p_i] > -\infty$}
      \If{$f[i-1][t-p_i] + w_i > f[i][t]$}
        \State $f[i][t] \gets f[i-1][t-p_i] + w_i$
        \State $\mathit{choice}[i][t] \gets \textsc{true}$ \Comment{site $i$ included to reach this value}
      \EndIf
    \EndIf
  \EndFor
\EndFor
\State $t^\ast \gets \arg\max_{t \in \{0,\dots,P\}} f[n][t]$; \ $W^\ast \gets f[n][t^\ast]$
\State $S^\ast \gets \emptyset$; \ $t \gets t^\ast$
\For{$i = n,\dots,1$} \Comment{backtrack}
  \If{$\mathit{choice}[i][t]$}
    \State $S^\ast \gets S^\ast \cup \{i\}$; \ $t \gets t - p_i$
  \EndIf
\EndFor
\State \Return $W^\ast$, $S^\ast$
\end{algorithmic}
\end{algorithm}

\subsection{A fully polynomial-time approximation scheme}

\begin{theorem}[FPTAS]\label{thm:fptas}
For every $\epsilon \in (0,1)$, a dispatch order $\sigma$ with
$W(\sigma) \ge (1-\epsilon) \cdot W^\ast$ (where $W^\ast$ is the true
optimum) can be computed in time polynomial in $n$ and $1/\epsilon$.
\end{theorem}

\begin{proof}
We adapt the standard value-scaling technique used for the 0/1 knapsack
FPTAS \citep{ibarrakim1975} to the dynamic program of
Theorem~\ref{thm:dp}. Fix $\epsilon \in (0,1)$ and let $w_{\max} :=
\max_i w_i$. Define the scaling factor $K := \max(1, \epsilon\, w_{\max}
/ n)$ and scaled weights $w_i' := \lfloor w_i / K \rfloor$, so that
$w_i' \le n/\epsilon$ for every $i$ and hence $\sum_i w_i' \le n^2 /
\epsilon$.

Run the dynamic program of Theorem~\ref{thm:dp}, but in its
\emph{value-indexed} dual form: for $i = 0,\dots,n$ and $v =
0,\dots,\lfloor n^2/\epsilon\rfloor$, let
\DIFdelbegin \[
\DIFdel{g(i,v) := \text{the minimum total dispatch time } \sum_{j\in S} p_j
\text{ over feasible } S \subseteq \{1,\dots,i\} \text{ with }
\sum_{j\in S} w_j' = v,
}\]%DIFAUXCMD
\DIFdelend \DIFaddbegin \begin{multline*}
\DIFadd{g(i,v) := \text{the minimum total dispatch time } \sum_{j\in S} p_j}\\
\DIFadd{\text{ over feasible } S \subseteq \{1,\dots,i\} \text{ with }
\sum_{j\in S} w_j' = v,
}\end{multline*}\DIFaddend 
with the same feasibility notion and the same recursion structure as
\eqref{eq:dp}, transposed to track minimum time for a target scaled
value rather than maximum value for a target time; the recursion is
\[
g(i,v) = \min\Bigl(g(i-1,v),\ \bigl[g(i-1,v-w_i') + p_i \le d_i\bigr]
\cdot \bigl(g(i-1,v-w_i') + p_i\bigr)\Bigr),
\]
by the same argument as in Theorem~\ref{thm:dp}'s proof, with the roles
of the time and value axes exchanged. This table has $O(n \cdot
n^2/\epsilon) = O(n^3/\epsilon)$ entries, each computed in $O(1)$ time,
so the whole computation runs in time polynomial in $n$ and $1/\epsilon$.

Let $\hat S$ be the feasible subset of largest scaled value $v$ with
$g(n,v) < \infty$, and let $S^\ast$ be the true optimal subset (value
$W^\ast = \sum_{i\in S^\ast} w_i$). We consider two cases, according to
which argument of $K=\max(1,\epsilon w_{\max}/n)$ attains the maximum.

\emph{Case $\epsilon w_{\max}/n < 1$.} Then $K=1$, so $w_i' = \lfloor
w_i \rfloor = w_i$ for every $i$ (the $w_i$ are integers by
Definition~\ref{def:mwhed}): the scaling is vacuous, the dynamic
program above is exactly Theorem~\ref{thm:dp}'s dynamic program, and it
returns $\hat S = S^\ast$, so $W(\hat S) = W^\ast \ge (1-\epsilon)
W^\ast$ trivially (as $\epsilon \in (0,1)$ and $W^\ast \ge 0$).

\emph{Case $\epsilon w_{\max}/n \ge 1$.} Then $K = \epsilon w_{\max}/n$
exactly (not merely an upper bound), so $Kn = \epsilon w_{\max}$
exactly. Since $w_i' \le w_i/K$, we have $\sum_{i \in S^\ast} w_i' \ge
\sum_{i\in S^\ast}(w_i/K - 1) \ge W^\ast/K - n$, and $S^\ast$ is itself
a feasible subset achieving this scaled value or better, so the
returned $\hat S$ satisfies $\sum_{i \in \hat S} w_i' \ge W^\ast/K - n$.
Converting back, since $w_i \le K(w_i' + 1)$,
\begin{align*}
W(\hat S) = \sum_{i\in \hat S} w_i
&\ \ge\ K \sum_{i \in \hat S} w_i'
\ \ge\ K\Bigl(\frac{W^\ast}{K} - n\Bigr) = W^\ast - Kn \\
&\ =\ W^\ast - \epsilon\, w_{\max}
\ \ge\ W^\ast - \epsilon\, W^\ast = (1-\epsilon) W^\ast,
\end{align*}
using $Kn = \epsilon w_{\max}$ (this case) and $w_{\max} \le W^\ast$
(Assumption~\ref{as:feasible}). Either case gives $W(\hat S) \ge
(1-\epsilon)W^\ast$.
\end{proof}

\begin{example}[The FPTAS on the running instance]\label{ex:fptas}
Apply Theorem~\ref{thm:fptas} to Example~\ref{ex:running} with a
deliberately loose $\epsilon = 0.5$. The scaling factor is $K =
\max(1,\, 0.5 \cdot 10/4) = 1.25$, and the scaled weights are $w' =
(\lfloor 5/1.25\rfloor, \lfloor 8/1.25\rfloor, \lfloor 3/1.25\rfloor,
\lfloor 10/1.25\rfloor) = (4,6,2,8)$. Re-running the dynamic program on
these scaled weights (tracking minimum dispatch time per scaled value,
as in the proof) again selects $S=\{1,2,4\}$, recovering the true
optimum $W^\ast=23$ exactly---comfortably inside the guaranteed bound
$(1-\epsilon)W^\ast = 11.5$. \DIFdelbegin \DIFdel{This mirrors Section~\ref{sec:experiments}'s
finding at scale: the FPTAS's worst-case bound is rarely tight in
practice, even at a loose $\epsilon$}\DIFdelend \DIFaddbegin \DIFadd{Here the
scaling is active ($K>1$) yet harmless; Section~\ref{sec:scaling} shows
that this is typical of random instances, while
Proposition~\ref{prop:tight} gives instances on which the guarantee is
nearly attained}\DIFaddend .
\end{example}

Algorithm~\ref{alg:fptas} states the FPTAS's value-indexed construction
explicitly. It has exactly the same shape as Algorithm~\ref{alg:exact},
with the roles of the time axis and the (scaled) value axis exchanged,
which is why its table size no longer depends on $P$.

\begin{algorithm}
\caption{FPTAS for MWHED (Theorem~\ref{thm:fptas})}\label{alg:fptas}
\begin{algorithmic}[1]
\Require sites \DIFdelbegin \DIFdel{$1,\dots,n$ with $d_1\le\dots\le d_n$, }\DIFdelend \DIFaddbegin \DIFadd{with }\DIFaddend dispatch times $p_i$, \DIFaddbegin \DIFadd{deadlines $d_i$, }\DIFaddend weights $w_i$, accuracy $\epsilon\in(0,1)$
\Ensure a dispatch order with on-time weight $\ge (1-\epsilon)W^\ast$
\State \DIFaddbegin \DIFadd{delete every site with $p_i > d_i$; relabel so $d_1\le\dots\le d_n$
}\State \DIFaddend $w_{\max} \gets \max_i w_i$; \ $K \gets \max(1,\, \epsilon\, w_{\max}/n)$
\State $w_i' \gets \lfloor w_i/K \rfloor$ for $i=1,\dots,n$; \ $V \gets \sum_{i=1}^n w_i'$
\State $g[0][0] \gets 0$; \ $g[0][v] \gets +\infty$ for $v=1,\dots,V$
\For{$i = 1,\dots,n$}
  \For{$v = 0,\dots,V$}
    \State $g[i][v] \gets g[i-1][v]$ \Comment{exclude site $i$}
    \If{$v \ge w_i'$ \textbf{and} $g[i-1][v-w_i'] < +\infty$ \textbf{and} $g[i-1][v-w_i'] + p_i \le d_i$}
      \If{$g[i-1][v-w_i'] + p_i < g[i][v]$}
        \State $g[i][v] \gets g[i-1][v-w_i'] + p_i$
        \State $\mathit{choice}[i][v] \gets \textsc{true}$
      \EndIf
    \EndIf
  \EndFor
\EndFor
\State $v^\ast \gets \max\{v : g[n][v] < +\infty\}$
\State $\hat S \gets \emptyset$; \ $v \gets v^\ast$
\For{$i = n,\dots,1$} \Comment{backtrack, as in Algorithm~\ref{alg:exact}}
  \If{$\mathit{choice}[i][v]$}
    \State $\hat S \gets \hat S \cup \{i\}$; \ $v \gets v - w_i'$
  \EndIf
\EndFor
\State \Return $\hat S$, dispatched in increasing-deadline order
\end{algorithmic}
\end{algorithm}

\DIFdelbegin %DIFDELCMD < \subsection{A second tractable special case}
%DIFDELCMD < %%%
\DIFdelend \DIFaddbegin \begin{remark}[Zero-scaled sites and an optional completion step]\label{rem:completion}
\DIFadd{When $K>1$ a site with $w_i<K$ has scaled weight $w_i'=0$ and, because
Algorithm~\ref{alg:fptas} minimizes dispatch time for each scaled value,
such a site is never selected. The bound of Theorem~\ref{thm:fptas}
accounts for the weight so lost (at most $K$ per site), and
Proposition~\ref{prop:tight} shows the loss can approach $\epsilon$ of
the optimum. A natural post-processing step, which cannot decrease the
returned weight and so preserves Theorem~\ref{thm:fptas}, is to
re-insert rejected sites in decreasing order of weight while the set
stays feasible (checkable in $O(n)$ by Lemma~\ref{lem:edd}); we use it
as the ``completed'' variant in Section~\ref{sec:scaling}.
}\end{remark}

\subsection{A second tractable special case: equal dispatch times}\label{sec:equalp}
\DIFaddend 

Theorem~\ref{thm:hardness} shows that heterogeneous dispatch costs,
weighed against heterogeneous criticality under a shared deadline,
already suffice for hardness. The next result shows \DIFaddbegin \DIFadd{that }\DIFaddend the reverse
axis---heterogeneous deadlines under a shared dispatch cost---does not.
\DIFaddbegin \DIFadd{We state it for $m\ge1$ identical vehicles. In }\emph{\DIFadd{MWHED-$m$}} \DIFadd{each of
$m$ vehicles, all starting at the depot at time $0$, serves a sequence
of round trips; each site is served by at most one vehicle, takes the
occupancy $p_i$ regardless of the vehicle, and is on time if its
completion time on its vehicle is at most $d_i$. MWHED-$1$ is MWHED.
}\DIFaddend 

\DIFdelbegin %DIFDELCMD < \begin{theorem}[Equal dispatch cost is tractable]%%%
\DIFdelend \DIFaddbegin \DIFadd{For readers who do not work with matroids we recall the notions used.
A }\emph{\DIFadd{matroid}} \DIFadd{on a finite ground set $E$ is a family $\mathcal I$ of
}\emph{\DIFadd{independent}} \DIFadd{subsets of $E$ with (i) $\emptyset\in\mathcal I$;
(ii) if $B\in\mathcal I$ and $A\subseteq B$ then $A\in\mathcal I$; and
(iii) the }\emph{\DIFadd{augmentation property}}\DIFadd{: if $A,B\in\mathcal I$ and
$|A|<|B|$, then $A\cup\{x\}\in\mathcal I$ for some $x\in B\setminus A$.
The }\emph{\DIFadd{greedy algorithm}} \DIFadd{considers the elements of $E$ in order of
non-increasing weight and keeps an element if the set kept so far plus
that element is still independent. The }\emph{\DIFadd{transversal matroid}} \DIFadd{of a
bipartite graph has as independent sets the subsets of one side that can
be matched into the other side; it is used here with sites on one side
and (vehicle, position) }\emph{\DIFadd{slots}} \DIFadd{on the other. We prove the matroid
property we need directly rather than cite it. A }\emph{\DIFadd{union--find}}
\DIFadd{(disjoint-set) structure maintains a partition of $\{0,1,\dots,n\}$ under
}\textsc{\DIFadd{find}}\DIFadd{, which returns the representative of a block, and
}\textsc{\DIFadd{union}}\DIFadd{, which merges two blocks; with path compression a
sequence of $n$ operations costs $O(n\,\alpha(n))$ time
\mbox{%DIFAUXCMD
\citep{kortevygen2018}}\hskip0pt%DIFAUXCMD
, where $\alpha$ is the inverse Ackermann
function, which grows more slowly than any iterated logarithm and is at
most $4$ for every input size that can occur in practice.
}

\begin{theorem}[Equal dispatch cost is tractable, for any number of identical vehicles]\DIFaddend \label{thm:equalp}
If \DIFdelbegin \DIFdel{$p_i = p$ for all }\DIFdelend \DIFaddbegin \DIFadd{$p_i=p$ for all }\DIFaddend $i$\DIFdelbegin \DIFdel{(all sites equidistant from the depot }\DIFdelend , \DIFdelbegin \DIFdel{in
round-trip dispatch time ), MWHED }\DIFdelend \DIFaddbegin \DIFadd{MWHED-$m$ }\DIFaddend is solvable exactly in \DIFdelbegin \DIFdel{$O(n \log n)$
time }\DIFdelend \DIFaddbegin \DIFadd{$O(n\log n)$
time for every $m\ge1$. In particular this holds for MWHED
($m=1$)}\DIFaddend .
\end{theorem}

\begin{proof}
\DIFdelbegin \DIFdel{With equal dispatch costs, the $k$}\DIFdelend \DIFaddbegin \emph{\DIFadd{Step 1: positions and slots.}} \DIFadd{Delete sites with $p>d_i$
(Assumption~\ref{as:feasible}) and put $D_i:=\min(\lfloor d_i/p\rfloor,n)\ge1$.
On any vehicle the site served $j$}\DIFaddend -th \DIFdelbegin \DIFdel{dispatched site (whichever it is) }\DIFdelend completes at time \DIFdelbegin \DIFdel{$kp$, regardless of order. Assigning a site to
``dispatch position '' $k$ is therefore feasible for that site if and only if $kp \le d_i$, i.e., $k \le \lfloor d_i/p \rfloor =: D_i$.
The problem becomes: choose an injective assignmentof a subset of
sites to positions $1,\dots,n$, each site $i$ assigned to a position at
most $D_i$, maximizing the total weight of assigned sites }\DIFdelend \DIFaddbegin \DIFadd{$jp$, so it is
on time iff $j\le D_i$. A set $S$ of sites is }\emph{\DIFadd{feasible}} \DIFadd{if all its
sites can be served on time. Call $(v,j)$ the slot of vehicle $v$ at
position $j$. Then $S$ is feasible iff its sites can be assigned to
distinct slots with each site $i$ at a position $\le D_i$: given a
schedule, assign each site to the slot it occupies; conversely, given
such an assignment, serve on each vehicle its assigned sites in
position order, shifting them to consecutive positions $1,2,\dots$,
which only lowers positions.
}

\emph{\DIFadd{Step 2: a counting criterion.}} \DIFadd{For a set $S$ and $k=1,\dots,n$ let
$N_S(k):=|\{i\in S:D_i\le k\}|$. We claim $S$ is feasible iff
$N_S(k)\le mk$ for all $k$. ($\Rightarrow$) The $N_S(k)$ sites with
$D_i\le k$ need distinct slots at positions $\le k$, and there are $mk$
such slots. ($\Leftarrow$) Order $S$ by non-decreasing $D$ as
$s_1,\dots,s_t$ and give $s_r$ the slot at position $\lceil r/m\rceil$
(filling $m$ vehicles at each position in turn). It is on time iff
$\lceil r/m\rceil\le D_{s_r}$; and $r\le N_S(D_{s_r})\le mD_{s_r}$, so
$r/m\le D_{s_r}$ and, $D_{s_r}$ being an integer,
$\lceil r/m\rceil\le D_{s_r}$}\DIFaddend .
\DIFdelbegin \DIFdel{This is the
classical problem of scheduling unit-time jobs with deadlines to maximize total profit. The family of subsets that admit such a feasible
assignment forms a matroid (the transversal matroid of the bipartite
feasibility relation between sites and positions) }\DIFdelend \DIFaddbegin 

\emph{\DIFadd{Step 3: the feasible sets form a matroid.}} \DIFadd{Properties (i) and (ii)
are clear from the criterion. For (iii) let $A,B$ be feasible with
$|A|<|B|$. Let $k$ be the largest integer in $\{0,\dots,n\}$ with
$N_B(k)\le N_A(k)$ (with $N(0):=0$, so $k$ exists); as
$N_B(n)=|B|>|A|=N_A(n)$ we have $k<n$, and by maximality
$N_B(j)\ge N_A(j)+1$ for all $j>k$. Subtracting $N_B(k)\le N_A(k)$
from $N_B(k+1)\ge N_A(k+1)+1$ shows that $B$ has more sites with
$D=k+1$ than $A$ has, so some $x\in B\setminus A$ has $D_x=k+1$. Adding
$x$ to $A$ leaves $N(j)$ unchanged for $j\le k$ and raises it by one for
$j\ge k+1$, where $N_A(j)+1\le N_B(j)\le mj$ since $B$ is feasible. By
Step 2}\DIFaddend , \DIFaddbegin \DIFadd{$A\cup\{x\}$ is feasible.
}

\emph{\DIFadd{Step 4: greedy is optimal.}} \DIFadd{Let $G=\{g_1,\dots,g_r\}$ be the
set the greedy algorithm returns (weights $w_{g_1}\ge\dots\ge w_{g_r}$,
ties broken by index), and let $O=\{o_1,\dots,o_s\}$ be any feasible set
with $w_{o_1}\ge\dots\ge w_{o_s}$. $G$ is maximal feasible, so
augmentation gives $s\le r$. We claim $w_{g_k}\ge w_{o_k}$ for all
$k\le s$, which proves $w(G)\ge w(O)$. Suppose $w_{g_k}<w_{o_k}$ for some
$k$. Apply augmentation to $A=\{g_1,\dots,g_{k-1}\}$ }\DIFaddend and
\DIFdelbegin \DIFdel{on any matroid, the greedy algorithm---here: sort }\DIFdelend \DIFaddbegin \DIFadd{$B=\{o_1,\dots,o_k\}$: some $o_j\in B\setminus A$ makes $A\cup\{o_j\}$
feasible, and $w_{o_j}\ge w_{o_k}>w_{g_k}$. The greedy algorithm
examined $o_j$ before $g_k$, when its kept set was a subset of $A$; by
property (ii) that set plus $o_j$ was feasible, so greedy kept $o_j$,
i.e.\ $o_j\in G$ and, being examined before $g_k$, $o_j\in A$---a
contradiction.
}

\emph{\DIFadd{Step 5: implementation.}} \DIFadd{Sort the }\DIFaddend sites by decreasing \DIFdelbegin \DIFdel{$w_i$,
and assign
each in turn to the latest available position
$k \le D_i$, or leave it
unassigned if noneis available---returns a maximum-weight independent
set , i. e. , an optimal solution \mbox{%DIFAUXCMD
\citep[Ch.~13]{kortevygen2018}}\hskip0pt%DIFAUXCMD
. Sorting
takes $O(n \log n)$ time, and maintaining the set of available
positions under repeated ``find andremove the latest available
position $\le D_i$'' queries can be implemented with a }\DIFdelend \DIFaddbegin \DIFadd{weight,
$O(n\log n)$. Maintain for each position $q\in\{1,\dots,n\}$ its number
of free slots, initially $m$, and a }\DIFaddend union--find \DIFdelbegin \DIFdel{structure in $O(n \alpha(n))$ additional time, giving $O(n \log n)$
overall}\DIFdelend \DIFaddbegin \DIFadd{structure over
$\{0,\dots,n\}$ in which }\textsc{\DIFadd{find}}\DIFadd{$(q)$ is the largest position
$\le q$ with a free slot (or $0$ if none): when position $q$ has no free
slot left, merge it into $q-1$. To process site $i$, let
$q:=\textsc{find}(D_i)$; if $q\ge1$ keep $i$ and use one slot of $q$,
otherwise discard $i$. This is Algorithm~\ref{alg:matroid}. It keeps
$i$ exactly when the kept set plus $i$ is feasible. If $q\ge1$ the
construction itself exhibits a feasible assignment. If $q=0$, let $q^\ast$
be the largest $q\ge D_i$ such that all slots at positions $1,\dots,q$
are full. Every kept site occupying a slot at a position $\le q^\ast$ has
$D\le q^\ast$: a kept site with $D>q^\ast$ would have been given the
latest free slot at a position $\le D$, and position $q^\ast+1$ was free
then (it is free now), so it would sit above $q^\ast$. Hence the kept
set has $N(q^\ast)\ge mq^\ast$ and, since $D_i\le q^\ast$, adding $i$
violates Step 2's criterion. The union--find operations cost
$O(n\alpha(n))$, so the total is $O(n\log n)$}\DIFaddend .
\end{proof}

\begin{example}[The matroid greedy on an equal-cost instance]\label{ex:matroid}
Consider four sites all at the same round-trip distance $p=2$ from the
depot, with deadlines $d=(3,3,5,9)$ and weights $w=(9,4,7,6)$, so the
latest feasible dispatch position for each site is $D_i = \lfloor
d_i/p\rfloor = (1,1,2,4)$. \DIFdelbegin \DIFdel{Sorted }\DIFdelend \DIFaddbegin \DIFadd{With one vehicle ($m=1$), sorted }\DIFaddend by
decreasing weight, the greedy order is site $1$ ($w=9$), site $3$
($w=7$), site $4$ ($w=6$), site $2$ ($w=4$). Site $1$ takes position
$1$\DIFdelbegin \DIFdel{(its only option)}\DIFdelend ; site $3$ takes position $2$; site $4$ takes position $4$; site $2$
needs a position $\le 1$, but position $1$ is already taken, so \DIFdelbegin \DIFdel{site $2$ }\DIFdelend \DIFaddbegin \DIFadd{it }\DIFaddend is
left unassigned. The greedy solution dispatches sites $1,3,4$ for a
total weight of $9+7+6=22$\DIFdelbegin \DIFdel{, sacrificing only the lowest-weight site whose
deadline conflicts with a higher-weight site's---exhaustive }\DIFdelend \DIFaddbegin \DIFadd{; exhaustive }\DIFaddend search over all $16$ subsets
confirms $22$ is \DIFdelbegin \DIFdel{indeed optimal. }\DIFdelend \DIFaddbegin \DIFadd{optimal. With two vehicles ($m=2$) position $1$ has two
slots, so site $2$ also fits: all four sites are served, for $26$
(vehicle~A serves site $1$ then site $3$, completing at $2$ and $4$;
vehicle~B serves site $2$ then site $4$, completing at $2$ and $4$).
}\DIFaddend \end{example}

\begin{algorithm}
\DIFdelbegin %DIFDELCMD < \caption{Matroid greedy for equal-cost MWHED (Theorem~\ref{thm:equalp})}%%%
\DIFdelend \DIFaddbegin \caption{Matroid greedy for equal-cost MWHED-$m$ (Theorem~\ref{thm:equalp})}\DIFaddend \label{alg:matroid}
\begin{algorithmic}[1]
\Require sites \DIFdelbegin \DIFdel{$1,\dots,n$}\DIFdelend \DIFaddbegin \DIFadd{with deadlines $d_i$ and weights $w_i$}\DIFaddend , common dispatch time $p$, \DIFdelbegin \DIFdel{deadlines $d_i$, weights $w_i$
}\DIFdelend \DIFaddbegin \DIFadd{number of identical vehicles $m$
}\DIFaddend \Ensure an optimal on-time subset $S^\ast$
\State \DIFdelbegin \DIFdel{$D_i \gets \lfloor d_i/p \rfloor$ for $i=1,\dots,n$ }%DIFDELCMD < \Comment{latest feasible position for site $i$}
%DIFDELCMD < \State %%%
\DIFdel{sort sites into order $\pi_1,\dots,\pi_n$ with $w_{\pi_1} \ge w_{\pi_2} \ge \dots \ge w_{\pi_n}$
}\DIFdelend \DIFaddbegin \DIFadd{delete sites with $p>d_i$; $D_i \gets \min(\lfloor d_i/p \rfloor,n)$ }\Comment{latest feasible position}
\DIFaddend \State \DIFdelbegin \DIFdel{$\mathit{available} \gets \{1,\dots,n\}$}\DIFdelend \DIFaddbegin \DIFadd{$\mathit{free}[q]\gets m$ for $q=1,\dots,n$}\DIFaddend ; \ \DIFdelbegin \DIFdel{$S^\ast \gets \emptyset$
}%DIFDELCMD < \For{$k = 1,\dots,n$}
%DIFDELCMD <   \State %%%
\DIFdel{$i \gets \pi_k$
  }%DIFDELCMD < \If{$\exists$ position $q \in \mathit{available}$ with $q \le D_i$}
%DIFDELCMD <     %%%
\DIFdelend \DIFaddbegin \DIFadd{$\textsc{find}$ over $\{0,\dots,n\}$: every $q$ its own block; \ $S^\ast\gets\emptyset$
}\For{each site $i$ in order of non-increasing $w_i$}
  \DIFaddend \State \DIFdelbegin \DIFdel{$q^\ast \gets \max\{q \in \mathit{available} : q \le D_i\}$ }%DIFDELCMD < \Comment{latest such position}
%DIFDELCMD <     %%%
\DIFdelend \DIFaddbegin \DIFadd{$q \gets \textsc{find}(D_i)$ }\Comment{largest position $\le D_i$ with a free slot, or $0$}
  \If{$q\ge1$}
    \DIFaddend \State \DIFdelbegin \DIFdel{$S^\ast \gets S^\ast \cup \{i\}$}\DIFdelend \DIFaddbegin \DIFadd{$S^\ast\gets S^\ast\cup\{i\}$}\DIFaddend ; \ \DIFdelbegin \DIFdel{$\mathit{available} \gets \mathit{available} \setminus \{q^\ast\}$
  }\DIFdelend \DIFaddbegin \DIFadd{$\mathit{free}[q]\gets\mathit{free}[q]-1$
    }\If{$\mathit{free}[q]=0$} \State \textsc{\DIFadd{union}}\DIFadd{$(q,q-1)$ }\Comment{$q$ is now represented by $q-1$} \EndIf
  \DIFaddend \EndIf
\EndFor
\State \Return $S^\ast$, dispatched in increasing-deadline order \DIFaddbegin \DIFadd{(on one vehicle) or position order (on $m$)
}\DIFaddend \end{algorithmic}
\end{algorithm}

\begin{remark}\label{rem:hardness-boundary}
Theorems~\ref{thm:hardness} and~\ref{thm:equalp} \DIFdelbegin \DIFdel{together }\DIFdelend show that neither
heterogeneous dispatch cost nor heterogeneous deadlines alone is
responsible for MWHED's hardness\DIFdelbegin \DIFdel{: it is specifically the interaction
between heterogeneous dispatch cost }\emph{\DIFdel{and}} %DIFAUXCMD
\DIFdel{heterogeneous
criticality, competing under a shared deadline structure, that produces
a knapsack-hard problem}\DIFdelend \DIFaddbegin \DIFadd{; Theorem~\ref{thm:classification}
below makes this exact}\DIFaddend . A dispatch network in which every site is
equally far from the depot---a plausible approximation for, e.g., a
ring of substations around a central facility, or short-radius urban
service territories---remains exactly solvable regardless of how
heterogeneous the sites' hazard deadlines or criticalities are\DIFaddbegin \DIFadd{, and
regardless of the number of identical vehicles}\DIFaddend .
\end{remark}

\DIFdelbegin %DIFDELCMD < \subsection{The cost of not optimizing}
%DIFDELCMD < %%%
\DIFdelend \DIFaddbegin \subsection{The cost of not reconsidering}
\DIFaddend 

Theorems~\ref{thm:dp} and~\ref{thm:fptas} give planners exact and
near-exact tools, but a planner \DIFdelbegin \DIFdel{who instead falls back on the simplest conceivable heuristic---dispatch }\DIFdelend \DIFaddbegin \DIFadd{might instead fall back on a simple
rule. The simplest is to dispatch }\DIFaddend sites in increasing-deadline order and
simply accept whichever happen to still be on time, never skipping a
site to make room for a later, more valuable \DIFdelbegin \DIFdel{one---receives no
guaranteeat all}\DIFdelend \DIFaddbegin \DIFadd{one (}\emph{\DIFadd{naive EDD}}\DIFadd{; a
site that is already late still consumes its dispatch time). A slightly
smarter rule skips a site that would arrive late (}\emph{\DIFadd{EDD with
skipping}}\DIFadd{). Neither has any guarantee}\DIFaddend .

\DIFdelbegin %DIFDELCMD < \begin{proposition}[Naive earliest-deadline dispatch has unbounded worst-case ratio]%%%
\DIFdelend \DIFaddbegin \begin{proposition}[Deadline-order dispatch has unbounded worst-case ratio]\DIFaddend \label{prop:naive}
For every $r \in (0,1)$ there is an MWHED instance, satisfying
Assumption~\ref{as:feasible}, on which \DIFdelbegin \DIFdel{the naive policy's }\DIFdelend \DIFaddbegin \DIFadd{both naive EDD and EDD with
skipping obtain }\DIFaddend on-time weight \DIFdelbegin \DIFdel{is }\DIFdelend strictly less than $r \cdot W^\ast$.
\end{proposition}

\begin{proof}
Fix $r \in (0,1)$ and choose an integer $W > 1/r$. Consider the
two-site instance with $p = (1, 2)$, $d = (1, 2)$, $w = (1, W)$; both
sites satisfy Assumption~\ref{as:feasible} ($p_i = d_i$). Since
$d_1 < d_2$, \DIFdelbegin \DIFdel{the naive policy dispatches }\DIFdelend \DIFaddbegin \DIFadd{both rules dispatch }\DIFaddend site $1$ first, completing at time
$1 \le d_1$ (on time, weight $1$)\DIFdelbegin \DIFdel{, then site }\DIFdelend \DIFaddbegin \DIFadd{. Site }\DIFaddend $2$ \DIFdelbegin \DIFdel{, completing }\DIFdelend \DIFaddbegin \DIFadd{would then complete }\DIFaddend at time
$1+2=3 > d_2=2$\DIFdelbegin \DIFdel{(exposed, losing weight $W$); its }\DIFdelend \DIFaddbegin \DIFadd{: naive EDD serves it anyway and it is exposed, and EDD
with skipping discards it; either way the }\DIFaddend total is $1$. Dispatching site
$2$ alone \DIFdelbegin \DIFdel{instead }\DIFdelend completes it at time $2 \le d_2$, \DIFdelbegin \DIFdel{on time, }\DIFdelend for weight $W$; \DIFdelbegin \DIFdel{this is feasible and, since $W > 1 \ge$
the value of any subset not containing site $2$, it }\DIFdelend \DIFaddbegin \DIFadd{since $W>1$
this }\DIFaddend is optimal, so $W^\ast = W$. The \DIFdelbegin \DIFdel{naive policy's ratio is therefore }\DIFdelend \DIFaddbegin \DIFadd{ratio is }\DIFaddend $1/W < r$.
\DIFdelbegin %DIFDELCMD < 

%DIFDELCMD < %%%
\DIFdelend \end{proof}

\DIFdelbegin \DIFdel{Proposition~\ref{prop:naive} is not merely a
boundary case: the two
sites here differ only by a factor of
two in
dispatch time,
yet the naive policy's loss can be made arbitrarily large by scaling $W$ alone.
No fixed constant-factor guaranteeis possible for a policy that never
reconsiders an early, }\DIFdelend \DIFaddbegin \DIFadd{The loss comes from never giving up an early, }\DIFaddend low-value \DIFdelbegin \DIFdel{commitment against a later,
}\DIFdelend \DIFaddbegin \DIFadd{commitment for a
later, }\DIFaddend high-value \DIFdelbegin \DIFdel{one---exactly the reconsideration that Theorem~\ref{thm:dp}'s dynamic program and }\DIFdelend \DIFaddbegin \DIFadd{one; skipping doomed sites removes only the waste of
serving a site that is already late. It is natural to ask whether a
local repair rule avoids the problem. The weighted generalization of
Moore--Hodgson's rule used as a heuristic in
Section~\ref{sec:experiments}---process sites in deadline order and,
whenever the running schedule becomes infeasible, discard the scheduled
site of smallest weight-to-time ratio---does very well on random
instances, but it too has no guarantee.
}

\begin{algorithm}
\caption{Weighted greedy repair (heuristic; no guarantee)}\label{alg:repair}
\begin{algorithmic}[1]
\Require \DIFadd{sites with $d_1\le\dots\le d_n$ (after deleting sites with $p_i>d_i$), dispatch times $p_i$, weights $w_i$
}\Ensure \DIFadd{a feasible on-time subset $\mathit{kept}$ (not necessarily optimal)
}\State \DIFadd{$\mathit{kept} \gets \emptyset$; \ $\mathit{total} \gets 0$
}\For{$i = 1,\dots,n$}
  \State \DIFadd{$\mathit{kept} \gets \mathit{kept} \cup \{i\}$; \ $\mathit{total} \gets \mathit{total} + p_i$
  }\While{$\mathit{total} > d_i$ \textbf{and} $\mathit{kept} \ne \emptyset$}
    \State \DIFadd{$j \gets \arg\min_{j' \in \mathit{kept}} w_{j'}/p_{j'}$ }\Comment{smallest weight-to-time ratio}
    \State \DIFadd{$\mathit{kept} \gets \mathit{kept} \setminus \{j\}$; \ $\mathit{total} \gets \mathit{total} - p_j$
  }\EndWhile
\EndFor
\State \Return \DIFadd{$\mathit{kept}$, dispatched in increasing-deadline order
}\end{algorithmic}
\end{algorithm}

\DIFadd{Its output is always feasible, by the same inductive argument as for the
unweighted algorithm: removing a site can only shorten the completion
time of every site that remains.
}

\begin{proposition}[Weighted greedy repair has unbounded worst-case ratio]\label{prop:repair}
\DIFadd{For every integer $k\ge2$ the two-site instance $p=(1,k)$, $d=(1,k)$,
$w=(2,2k-1)$ satisfies Assumption~\ref{as:feasible} and the weighted
greedy repair of Algorithm~\ref{alg:repair} returns on-time weight $2$,
while $W^\ast=2k-1$; the ratio $2/(2k-1)\to0$.
}\end{proposition}

\begin{proof}
\DIFadd{The algorithm keeps site $1$ (total $1\le d_1$). On adding site $2$ the
total becomes $1+k>d_2=k$, so it discards the kept site of smallest
ratio $w_j/p_j$: site $1$ has ratio $2$ and site $2$ has ratio
$(2k-1)/k<2$, so site $2$ is discarded and the result is $\{1\}$ with
weight $2$. Site $2$ alone completes at $k\le d_2$, so $W^\ast\ge2k-1$,
and the two sites cannot both be on time, so $W^\ast=2k-1$ for $k\ge2$.
}\end{proof}

\DIFadd{These two propositions say that some form of reconsideration governed
by the weights }\emph{\DIFadd{and}} \DIFadd{the dispatch times is needed for any
worst-case guarantee. They do not say that the dynamic program or the
FPTAS is the only way to obtain one: other polynomial algorithms with
guarantees may exist, and we do not claim otherwise.
}

\subsection{Further structural results}\label{sec:structure}

\paragraph{\DIFadd{Tightness of the FPTAS analysis.}} \DIFadd{The proof of
}\DIFaddend Theorem~\ref{thm:fptas} \DIFdelbegin \DIFdel{'s FPTAS perform, }\DIFdelend \DIFaddbegin \DIFadd{loses at most $Kn=\epsilon w_{\max}$ and then
bounds this by $\epsilon W^\ast$. The next result shows this is nearly
the truth for Algorithm~\ref{alg:fptas} as stated.
}

\begin{proposition}[Tightness of the value-scaling analysis]\label{prop:tight}
\DIFadd{Let $\epsilon\in(0,1)$, let $n\ge2$ and let $M\ge 2n/\epsilon$ be an
integer. There is an instance with $n$ sites, satisfying
Assumption~\ref{as:feasible}, on which Algorithm~\ref{alg:fptas}
returns weight $M$ while
$W^\ast = M+(n-1)(\lceil K\rceil-1)$ with $K=\epsilon M/n$. Hence its
ratio is at most
$\bigl(1+\epsilon\tfrac{n-1}{n}-\tfrac{n-1}{M}\bigr)^{-1}$, which tends to
$1/(1+\epsilon)$ as $n,M\to\infty$. Theorem~\ref{thm:fptas} guarantees
$1-\epsilon$, and $\frac{1/(1+\epsilon)}{1-\epsilon}=\frac{1}{1-\epsilon^2}$,
so the guarantee is tight up to a factor $1-\epsilon^2$.
}\end{proposition}

\begin{proof}
\DIFadd{Take $p_i=1$ }\DIFaddend and \DIFdelbegin \DIFdel{that Example~\ref{ex:running}
already previewed at a modest scale.
Section~\ref{sec:experiments}
confirms the same qualitative gap persists, if less extremely, on
random instances of
realistic size. }\DIFdelend \DIFaddbegin \DIFadd{$d_i=n$ for all $i$, $w_1=M$, and $w_i=\lceil K\rceil-1
\ (\ge1)$ for $i\ge2$. All $n$ sites can be served (total time $n$,
every deadline $n$), so $W^\ast=M+(n-1)(\lceil K\rceil-1)$. The scaling
factor is $K=\epsilon M/n\ge2>1$, site $1$ has scaled weight
$\lfloor n/\epsilon\rfloor$, and every other site has $w_i<K$, hence
scaled weight $0$. Algorithm~\ref{alg:fptas} minimizes the dispatch time
for each scaled value, so it never selects a zero-scaled site (adding it
only increases the time at an unchanged value) and returns $\{1\}$, of
weight $M$. Finally $\lceil K\rceil-1\ge K-1$ gives the stated bound.
}\end{proof}
\DIFaddend 

%DIF <  ══════════════════════════════════════════════════════════════════
\DIFdelbegin %DIFDELCMD < \section{Numerical illustration}\label{sec:experiments}
%DIFDELCMD < %%%
\DIFdelend \DIFaddbegin \DIFadd{The optional completion step of Remark~\ref{rem:completion} repairs this
family completely, which is why we report both variants in
Section~\ref{sec:scaling}.
}\DIFaddend 

\DIFdelbegin \DIFdel{We illustrate Theorems~\ref{thm:dp} and~\ref{thm:fptas} and Proposition~\ref{prop:naive} in five complementary experiments: an
accuracy comparison across four algorithms on synthetic instances ; a
runtime comparison isolating the pseudo-polynomial (}\DIFdelend \DIFaddbegin \paragraph{\DIFadd{Which heterogeneity causes hardness?}} \DIFadd{Theorems~\ref{thm:hardness}
and~\ref{thm:equalp} and Moore's algorithm together settle the
complexity of every restriction obtained by making some of the three
data vectors constant. For $T\subseteq\{p,w,d\}$ let $\mathcal C_T$ be
the class of MWHED instances in which each vector named in $T$ is
constant ($p_i=p$, $w_i=w$, $d_i=D$ for all $i$).
}

\begin{theorem}[Classification by homogeneity]\label{thm:classification}
\DIFadd{Let $T\subseteq\{p,w,d\}$.
}\begin{enumerate}
\item[\DIFadd{(a)}] \DIFadd{If $p\in T$ or $w\in T$, MWHED restricted to $\mathcal C_T$
is solvable in $O(n\log n)$ time.
}\item[\DIFadd{(b)}] \DIFadd{If $T\subseteq\{d\}$, MWHED restricted to $\mathcal C_T$ is
weakly NP-hard; in particular it is NP-hard for $T=\{d\}$.
}\end{enumerate}
\DIFadd{Thus an MWHED restriction is polynomial exactly when the dispatch times
or the weights are constant, and constancy of the deadlines never
helps.
}\end{theorem}

\begin{proof}
\DIFadd{(a) If $p\in T$, apply Theorem~\ref{thm:equalp} with $m=1$. If $w\in T$,
all sites have the same weight $w$, so maximizing on-time weight means
maximizing the }\emph{\DIFadd{number}} \DIFadd{of on-time sites, which is the problem
$1\|\sum U_j$ solved in $O(n\log n)$ time by Moore's algorithm
\mbox{%DIFAUXCMD
\citep{moore1968}}\hskip0pt%DIFAUXCMD
. (b) $\mathcal C_{\{d\}}\supseteq$ the instances of
Theorem~\ref{thm:hardness}, so it is NP-hard, and it is weakly NP-hard
because }\DIFaddend Theorem~\ref{thm:dp} \DIFdelbegin \DIFdel{)
versus truly polynomial (Theorem~\ref{thm:fptas})dependence on the instance's numerical magnitude; a sensitivity sweep over }\DIFdelend \DIFaddbegin \DIFadd{solves it in pseudo-polynomial time;
$\mathcal C_\emptyset$ contains $\mathcal C_{\{d\}}$.
}\end{proof}

\DIFadd{The classification is elementary once Theorems~\ref{thm:hardness}
and~\ref{thm:equalp} are available; its value is interpretive. Hardness
needs }\emph{\DIFadd{two independent}} \DIFadd{sources of heterogeneity, in the dispatch
time }\emph{\DIFadd{and}} \DIFadd{in the criticality, which compete for the shared time
budget; heterogeneity in the hazard deadlines, the part of the instance
produced by the hazard model, plays no role. In network terms: if every
site is equally far from the depot, or every site is equally critical,
the dispatch problem is easy however the hazard spreads.
}

%DIF >  ══════════════════════════════════════════════════════════════════
\section{Numerical study}\label{sec:experiments}

\DIFadd{The study has four purposes: to check every algorithm against an
independent exhaustive search (Section~\ref{sec:verify}); to compare the
algorithms and the heuristics of Section~\ref{sec:theory} across
instance sizes and generators, and to show, in a regime where }\DIFaddend the
FPTAS's \DIFdelbegin \DIFdel{accuracy parameter }\DIFdelend \DIFaddbegin \DIFadd{rounding is genuinely active, how its accuracy and cost behave
as }\DIFaddend $\epsilon$ \DIFdelbegin \DIFdel{; a robustness check repeating the
accuracy comparison under two alternative synthetic instance-generation
regimes; and
, finally, a small instance built directly from public
records of }\DIFdelend \DIFaddbegin \DIFadd{shrinks (Sections~\ref{sec:acc}--\ref{sec:robust}); and
to illustrate the model on }\DIFaddend a real wildfire\DIFaddbegin \DIFadd{, with a sensitivity analysis
over its uncertain inputs }\DIFaddend (Section~\ref{sec:campfire})\DIFdelbegin \DIFdel{, to confirm
the synthetic findings are not an artifact of how the synthetic
instances are generated}\DIFdelend \DIFaddbegin \DIFadd{. All code is
self-contained and deterministic given its seeds (see the Data
Availability statement); every table below is regenerated by
}\texttt{\DIFadd{python3 experiment.py}} \DIFadd{(about 15 seconds) and
}\texttt{\DIFadd{python3 case\_study\_campfire.py}}\DIFadd{.
}

\subsection{Verification against exhaustive search}\label{sec:verify}

\DIFadd{Before any performance claim we verify the implementations
(}\texttt{\DIFadd{verify\_small.py}}\DIFadd{). (i) On $400$ random instances with $n\le7$
the dynamic program's optimum equals the optimum over }\emph{\DIFadd{all}} \DIFadd{$n!$
dispatch orders, which checks Lemma~\ref{lem:edd} and
Theorem~\ref{thm:dp} from the definition. (ii) On $1{,}500$ random
instances with $n\le14$ (a third of them with weights up to $10^6$ and
all with individually infeasible sites left in) the dynamic program
equals the best feasible subset found by enumeration, and the FPTAS
returns a feasible set of weight at least $(1-\epsilon)W^\ast$ for
$\epsilon\in\{0.5,0.2,0.05\}$, with and without the completion step, with
no violation in $4{,}500$ runs ($2{,}923$ of which had active scaling).
(iii) The equal-cost greedy of Algorithm~\ref{alg:matroid} equals the
dynamic program on $300$ instances for $m=1$, and equals exhaustive
search over all assignments to $m=2,3$ vehicles on the instances with
$n\le7$. (iv) Proposition~\ref{prop:repair}'s ratio $2/(2k-1)$ is
reproduced for $k=3,10,100$}\DIFaddend .

\subsection{Accuracy across instance sizes}\DIFaddbegin \label{sec:acc}
\DIFaddend 

For \DIFdelbegin \DIFdel{$n \in \{10,15,20,25,30,50,75,100\}$ we generate 20 }\DIFdelend \DIFaddbegin \DIFadd{$n \in \{10,15,20,25,30,50,75,100,200,400\}$ we generate $50$
}\DIFaddend random instances each ($p_i$ uniform on $\{1,\dots,20\}$, $w_i$ uniform
on $\{1,\dots,100\}$, $d_i$ uniform on $\{1,\dots,\sum_i p_i\}$\DIFdelbegin \DIFdel{),
and
compare four algorithms}\DIFdelend \DIFaddbegin \DIFadd{,
individually infeasible sites deleted as in
Assumption~\ref{as:feasible}) and compare}\DIFaddend : the true optimum
\DIFdelbegin \DIFdel{via }\DIFdelend \DIFaddbegin \DIFadd{(}\DIFaddend Theorem~\ref{thm:dp}\DIFdelbegin \DIFdel{'s
exact dynamic program}\DIFdelend \DIFaddbegin \DIFadd{)}\DIFaddend ; the FPTAS \DIFdelbegin \DIFdel{of Theorem~\ref{thm:fptas} at $\epsilon \in \{0.1,0.2\}$,
implemented exactly }\DIFdelend \DIFaddbegin \DIFadd{at $\epsilon\in\{0.1,0.2\}$,
implemented }\DIFaddend as in the proof\DIFdelbegin \DIFdel{(the
value-indexed dual DP, not a shortcut); a weighted Moore--Hodgson-style
greedy repair
heuristic (dispatch in earliest-deadline order, and
whenever the running schedule becomes infeasible, repeatedly drop the
scheduled site with the smallest weight-to-time ratio until feasible
again---a natural, always-feasible heuristic generalizing
\mbox{%DIFAUXCMD
\citeauthor{moore1968}}\hskip0pt%DIFAUXCMD
's unweighted rule, for which we claim no
approximation guarantee}\DIFdelend \DIFaddbegin \DIFadd{; the weighted greedy repair
(Algorithm~\ref{alg:repair}}\DIFaddend ); and \DIFdelbegin \DIFdel{the naive
baseline of
}\DIFdelend \DIFaddbegin \DIFadd{two deadline-order baselines, naive
EDD and EDD with skipping (}\DIFaddend Proposition~\ref{prop:naive}\DIFdelbegin \DIFdel{(dispatch in earliest-deadline order,
simply keep whichever sites happen to still be on time, with no
repair at all)}\DIFdelend \DIFaddbegin \DIFadd{). Naive EDD
keeps dispatching to sites that are already late, so it is a weak
anchor; EDD with skipping is the fairer simple baseline}\DIFaddend .
Table~\ref{tab:illustration} reports \DIFdelbegin \DIFdel{the ratio of each
method's on-time weight to the true optimum, averaged over the 20
instances at each size. Algorithm~\ref{alg:repair} states the greedy
repair heuristic precisely; feasibility of its output is guaranteed by
the same inductive argument used for the unweighted Moore--Hodgson
algorithm (each repair step can only shorten, never lengthen, the
completion time of every remaining scheduled site), but---unlike
Algorithms~\ref{alg:exact}--\ref{alg:matroid}---we prove no
approximation ratio for it; it is included as a realistic,
easy-to-implement point of comparison, not as a fifth theorem}\DIFdelend \DIFaddbegin \DIFadd{mean ratios to the optimum, and
Figure~\ref{fig:accuracy} plots them}\DIFaddend .

\DIFdelbegin %DIFDELCMD < \begin{algorithm}
%DIFDELCMD < \caption{Weighted greedy repair (heuristic; no proven guarantee)}\label{alg:repair}
%DIFDELCMD < \begin{algorithmic}[1]
%DIFDELCMD < \Require %%%
\DIFdel{sites $1,\dots,n$ with $d_1\le\dots\le d_n$, dispatch times $p_i$, weights $w_i$
}%DIFDELCMD < \Ensure %%%
\DIFdel{a feasible on-time subset $\mathit{kept}$ (not necessarily optimal)
}%DIFDELCMD < \State %%%
\DIFdel{$\mathit{kept} \gets \emptyset$; \ $\mathit{total} \gets 0$
}%DIFDELCMD < \For{$i = 1,\dots,n$}
%DIFDELCMD <   \State %%%
\DIFdel{$\mathit{kept} \gets \mathit{kept} \cup \{i\}$; \ $\mathit{total} \gets \mathit{total} + p_i$
  }%DIFDELCMD < \While{$\mathit{total} > d_i$ \textbf{and} $\mathit{kept} \ne \emptyset$}
%DIFDELCMD <     \State %%%
\DIFdel{$j \gets \arg\min_{j' \in \mathit{kept}} w_{j'}/p_{j'}$ }%DIFDELCMD < \Comment{smallest weight-to-time ratio}
%DIFDELCMD <     \State %%%
\DIFdel{$\mathit{kept} \gets \mathit{kept} \setminus \{j\}$; \ $\mathit{total} \gets \mathit{total} - p_j$
  }%DIFDELCMD < \EndWhile
%DIFDELCMD < \EndFor
%DIFDELCMD < \State \Return %%%
\DIFdel{$\mathit{kept}$, dispatched in increasing-deadline order
}%DIFDELCMD < \end{algorithmic}
%DIFDELCMD < \end{algorithm}
%DIFDELCMD < 

%DIFDELCMD < %%%
\DIFdelend \begin{table}[h]
\DIFdelbeginFL %DIFDELCMD < \caption{Mean ratio to the true optimum, by instance size ($n$), over
%DIFDELCMD < 20 random instances per row. ``M--H style'' is the weighted
%DIFDELCMD < Moore--Hodgson-style greedy repair heuristic.}%%%
\DIFdelendFL \DIFaddbeginFL \caption{Mean ratio to the true optimum, by instance size $n$, over $50$
random instances per column (weights $\le100$). The last three rows give
the worst ratio over the $50$ instances for greedy repair and for naive
EDD, and the share of instances on which the FPTAS's scaling factor was
$K=1$ (no rounding at all), for $\epsilon=0.2$ and $\epsilon=0.1$.}\DIFaddendFL \label{tab:illustration}
\small
\DIFdelbeginFL %DIFDELCMD < \begin{tabular}{@{}lcccccccc@{}}
%DIFDELCMD < %%%
\DIFdelendFL \DIFaddbeginFL \begin{tabular}{@{}lcccccc@{}}
\DIFaddendFL \toprule
$n$ & 10 & \DIFdelbeginFL \DIFdelFL{15 }%DIFDELCMD < & %%%
\DIFdelendFL 20 & \DIFdelbeginFL \DIFdelFL{25 }%DIFDELCMD < & %%%
\DIFdelFL{30 }%DIFDELCMD < & %%%
\DIFdelendFL 50 & \DIFdelbeginFL \DIFdelFL{75 }%DIFDELCMD < & %%%
\DIFdelendFL 100 \DIFaddbeginFL & \DIFaddFL{200 }& \DIFaddFL{400 }\DIFaddendFL \\
\midrule
FPTAS ($\epsilon=0.2$) & 1.000 & 1.000 & 1.000 & 1.000 & 1.000 & 1.000 \DIFdelbeginFL %DIFDELCMD < & %%%
\DIFdelFL{1.000 }%DIFDELCMD < & %%%
\DIFdelFL{1.000 }\DIFdelendFL \\
FPTAS ($\epsilon=0.1$) & 1.000 & 1.000 & 1.000 & 1.000 & 1.000 & 1.000 \DIFdelbeginFL %DIFDELCMD < & %%%
\DIFdelFL{1.000 }%DIFDELCMD < & %%%
\DIFdelFL{1.000 }\DIFdelendFL \\
Greedy repair \DIFdelbeginFL \DIFdelFL{(M--H style) }\DIFdelendFL & 0.993 & \DIFdelbeginFL \DIFdelFL{0.997 }%DIFDELCMD < & %%%
\DIFdelendFL 0.996 & \DIFdelbeginFL \DIFdelFL{0.993 }%DIFDELCMD < & %%%
\DIFdelFL{0.997 }%DIFDELCMD < & %%%
\DIFdelFL{0.998 }\DIFdelendFL \DIFaddbeginFL \DIFaddFL{0.999 }\DIFaddendFL & 0.999 & 0.999 \DIFaddbeginFL & \DIFaddFL{1.000 }\\
\DIFaddFL{EDD with skipping }& \DIFaddFL{0.933 }& \DIFaddFL{0.941 }& \DIFaddFL{0.954 }& \DIFaddFL{0.972 }& \DIFaddFL{0.981 }& \DIFaddFL{0.983 }\DIFaddendFL \\
Naive EDD \DIFdelbeginFL \DIFdelFL{(no repair ) }\DIFdelendFL & \DIFdelbeginFL \DIFdelFL{0.380 }\DIFdelendFL \DIFaddbeginFL \DIFaddFL{0.714 }\DIFaddendFL & \DIFdelbeginFL \DIFdelFL{0.343 }\DIFdelendFL \DIFaddbeginFL \DIFaddFL{0.653 }\DIFaddendFL & \DIFdelbeginFL \DIFdelFL{0.622 }\DIFdelendFL \DIFaddbeginFL \DIFaddFL{0.595 }\DIFaddendFL & \DIFdelbeginFL \DIFdelFL{0.369 }\DIFdelendFL \DIFaddbeginFL \DIFaddFL{0.573 }\DIFaddendFL & \DIFdelbeginFL \DIFdelFL{0.323 }\DIFdelendFL \DIFaddbeginFL \DIFaddFL{0.669 }\DIFaddendFL & \DIFdelbeginFL \DIFdelFL{0.482 }\DIFdelendFL \DIFaddbeginFL \DIFaddFL{0.583 }\\
\midrule
\DIFaddFL{worst ratio, greedy repair }\DIFaddendFL & \DIFdelbeginFL \DIFdelFL{0.515 }\DIFdelendFL \DIFaddbeginFL \DIFaddFL{0.929 }\DIFaddendFL & \DIFdelbeginFL \DIFdelFL{0.439 }\DIFdelendFL \DIFaddbeginFL \DIFaddFL{0.963 }& \DIFaddFL{0.976 }& \DIFaddFL{0.992 }& \DIFaddFL{0.995 }& \DIFaddFL{0.999 }\\
\DIFaddFL{worst ratio, naive EDD }& \DIFaddFL{0.065 }& \DIFaddFL{0.132 }& \DIFaddFL{0.032 }& \DIFaddFL{0.110 }& \DIFaddFL{0.027 }& \DIFaddFL{0.034 }\\
\DIFaddFL{share with $K=1$, $\epsilon=0.2$ / $0.1$ }& \DIFaddFL{0/.54 }& \DIFaddFL{.60/1 }& \DIFaddFL{1/1 }& \DIFaddFL{1/1 }& \DIFaddFL{1/1 }& \DIFaddFL{1/1 }\DIFaddendFL \\
\bottomrule
\end{tabular}
\end{table}

\begin{figure}[h]
\centering
\includegraphics[width=0.72\textwidth]{figures/accuracy_vs_n.pdf}
\DIFdelbeginFL %DIFDELCMD < \caption{Table~\ref{tab:illustration} plotted: mean ratio to the true
%DIFDELCMD < optimum against instance size $n$. The FPTAS and greedy-repair curves
%DIFDELCMD < sit on top of each other at $1.0$; naive EDD fluctuates with no
%DIFDELCMD < trend of improving as $n$ grows, exactly as
%DIFDELCMD < Proposition~\ref{prop:naive} predicts.}%%%
\DIFdelendFL \DIFaddbeginFL \caption{Table~\ref{tab:illustration} plotted. The FPTAS and greedy
repair sit at $1.0$; EDD with skipping improves slowly with $n$; naive
EDD fluctuates around $0.6$ with no trend and, on individual instances,
captures only a few percent of the optimum.}\DIFaddendFL \label{fig:accuracy}
\end{figure}

\DIFdelbegin \DIFdel{At both approximation levels tested, the }\DIFdelend \DIFaddbegin \DIFadd{The }\DIFaddend FPTAS recovers the \DIFdelbegin \DIFdel{true
optimum on essentially every one of the 160 random instances (mean
ratio rounds to $1.000$ throughout), far outperforming its }\DIFdelend \DIFaddbegin \DIFadd{optimum on every instance, but this is
}\emph{\DIFadd{not}} \DIFadd{evidence about the approximation scheme. With weights at most
$100$ the scaling factor $K=\max(1,\epsilon w_{\max}/n)$ equals $1$ for
$\epsilon=0.1$ at every $n\ge15$ and for $\epsilon=0.2$ at every
$n\ge25$ (last row of Table~\ref{tab:illustration}), so no rounding
takes place and the FPTAS coincides with the exact value-indexed
dynamic program; a ratio of $1.000$ is then guaranteed, not
informative. The next subsection therefore uses weights large enough
that the rounding is active. Among the heuristics, the greedy repair
captures $99.3$--$100\%$ of the optimum on average but has a worst
ratio of $0.93$ at $n=10$, consistent with its having no guarantee
(Proposition~\ref{prop:repair}); naive EDD shows no trend with $n$, as
a rule that never reconsiders would not, though a }\DIFaddend worst-case \DIFdelbegin \DIFdel{$(1-\epsilon)$ guarantee, consistent with the well-documented empirical
looseness of knapsack-style FPTAS }\DIFdelend \DIFaddbegin \DIFadd{result such
as Proposition~\ref{prop:naive} does not by itself predict average
behavior on random instances.
}

\subsection{The FPTAS with active scaling}\label{sec:scaling}

\DIFadd{We now take weights up to $10^6$ with $p_i\le20$, so that the
time-indexed exact dynamic program stays cheap (its table depends on
$P$, not on the weights) and serves as the reference, while the
FPTAS's scale factor $K$ is in the hundreds to thousands. Three
families are used, with $n=100$ and $30$ instances per $\epsilon$ for
the first two. }\emph{\DIFadd{Random:}} \DIFadd{weights uniform on $\{1,\dots,10^6\}$.
}\emph{\DIFadd{Strongly correlated:}} \DIFadd{a common deadline and $w_i=1000p_i+500$ with
$p_i$ up to $1000$, the hard family for knapsack-type schemes
(\mbox{%DIFAUXCMD
\citealp{hejl2022} }\hskip0pt%DIFAUXCMD
study this regime for exact algorithms).
}\emph{\DIFadd{Adversarial:}} \DIFadd{the family of Proposition~\ref{prop:tight}, with
one dominant site and $n-1$ sites just below the scaling granule. For
each, Table~\ref{tab:scaling} reports the worst ratio over the instances
for Algorithm~\ref{alg:fptas} as stated and, for the adversarial family,
with the completion step of Remark~\ref{rem:completion}.
}

\begin{table}[h]
\caption{FPTAS accuracy with active scaling ($n=100$, $p_i\le20$ or as
stated). Entries are the \emph{worst} ratio to the optimum over the
instances of each family ($30$ per $\epsilon$ for the first two). ``Share $<1$''
is the fraction of random instances on which the FPTAS was not exactly
optimal. Every entry respects the guarantee $1-\epsilon$.}\label{tab:scaling}
\small
\begin{tabular}{@{}lcccccc@{}}
\toprule
\DIFaddFL{$\epsilon$ }& \DIFaddFL{$1-\epsilon$ }& \DIFaddFL{Random }& \DIFaddFL{Share $<1$ }& \DIFaddFL{Str.\ corr. }& \DIFaddFL{Adversarial }& \DIFaddFL{Adv.\ + completion }\\
\midrule
\DIFaddFL{0.50 }& \DIFaddFL{0.50 }& \DIFaddFL{0.9998 }& \DIFaddFL{0.30 }& \DIFaddFL{0.9993 }& \DIFaddFL{0.669 }& \DIFaddFL{1.000 }\\
\DIFaddFL{0.30 }& \DIFaddFL{0.70 }& \DIFaddFL{0.99996 }& \DIFaddFL{0.10 }& \DIFaddFL{0.9994 }& \DIFaddFL{0.771 }& \DIFaddFL{1.000 }\\
\DIFaddFL{0.20 }& \DIFaddFL{0.80 }& \DIFaddFL{0.99996 }& \DIFaddFL{0.10 }& \DIFaddFL{0.9994 }& \DIFaddFL{0.835 }& \DIFaddFL{1.000 }\\
\DIFaddFL{0.10 }& \DIFaddFL{0.90 }& \DIFaddFL{0.99999 }& \DIFaddFL{0.03 }& \DIFaddFL{0.9996 }& \DIFaddFL{0.910 }& \DIFaddFL{1.000 }\\
\DIFaddFL{0.05 }& \DIFaddFL{0.95 }& \DIFaddFL{1.00000 }& \DIFaddFL{0.00 }& \DIFaddFL{0.9996 }& \DIFaddFL{0.953 }& \DIFaddFL{1.000 }\\
\bottomrule
\end{tabular}
\end{table}

\begin{table}[h]
\caption{Cost of the FPTAS as $\epsilon$ shrinks (random weights up to
$10^6$, $n=100$): mean scaling factor $K$, mean number of dynamic-program
cells, and mean wall-clock time. A value-indexed exact table would need
about $5\times10^{7}$ columns per row at these weights, against the
$10^{4}$--$10^{5}$ used here.}\label{tab:epscost}
\small
\begin{tabular}{@{}lccccc@{}}
\toprule
\DIFaddFL{$\epsilon$ }& \DIFaddFL{0.50 }& \DIFaddFL{0.30 }& \DIFaddFL{0.20 }& \DIFaddFL{0.10 }& \DIFaddFL{0.05 }\\
\midrule
\DIFaddFL{mean $K$ }& \DIFaddFL{4}{\DIFaddFL{,}}\DIFaddFL{996 }& \DIFaddFL{2}{\DIFaddFL{,}}\DIFaddFL{997 }& \DIFaddFL{1}{\DIFaddFL{,}}\DIFaddFL{998 }& \DIFaddFL{999 }& \DIFaddFL{500 }\\
\DIFaddFL{cells ($\times10^6$) }& \DIFaddFL{0.99 }& \DIFaddFL{1.65 }& \DIFaddFL{2.47 }& \DIFaddFL{4.95 }& \DIFaddFL{9.90 }\\
\DIFaddFL{time (ms) }& \DIFaddFL{3.4 }& \DIFaddFL{4.9 }& \DIFaddFL{6.7 }& \DIFaddFL{14.4 }& \DIFaddFL{33.1 }\\
\bottomrule
\end{tabular}
\end{table}

\begin{figure}[h]
\centering
\includegraphics[width=0.95\textwidth]{figures/epsilon_sensitivity.pdf}
\caption{Tables~\ref{tab:scaling} and~\ref{tab:epscost} plotted. Left:
the worst ratio on the adversarial family tracks the guarantee
$1-\epsilon$ (it stays above it and, as Proposition~\ref{prop:tight}
shows, approaches $1/(1+\epsilon)$); on random instances the FPTAS is
essentially exact, and the completion step makes the adversarial
family exact. Right: cost grows in proportion to $1/\epsilon$, as the
$O(n^3/\epsilon)$ bound says.}\label{fig:epsilon}
\end{figure}

\DIFadd{Two things are visible. First, on random and on strongly correlated
instances the FPTAS is nearly exact even at $\epsilon=0.5$ (worst ratio
$0.9993$), far above its guarantee: the bound $\epsilon w_{\max}$
charges every selected site a full rounding error $K$ and compares the
total to $w_{\max}$, whereas on these instances the optimum contains
many sites and $W^\ast\gg w_{\max}$. This is why the ratio is close to
$1$, and it is the answer to the question of why accuracy is near $1$ in
all earlier tables. Second, the guarantee is }\emph{\DIFadd{not}} \DIFadd{vacuous: on the
adversarial family, which Proposition~\ref{prop:tight} shows is the
worst case, the ratio falls to $0.669$ at $\epsilon=0.5$ and follows
the guarantee as $\epsilon$ shrinks, so a planner who needs a }\DIFaddend worst-case
\DIFdelbegin \DIFdel{bounds
\mbox{%DIFAUXCMD
\citep{ibarrakim1975}}\hskip0pt%DIFAUXCMD
: the scaling construction is designed to
guarantee the worst case, not
to describe the typical one. The greedy
repair heuristic, despite carrying no proven guarantee, is a
substantial practical improvement over doing nothing---capturing
99.3--99.9\% of the optimum---because reconsidering even a single
locally-worst commitment when infeasibility first arises already
recovers most of the value a full optimization would. The naive EDD baseline, which never reconsiders anything, captures only
32--62\% of the
optimal on-time weight and shows no trend of improving
with
instance size, exactly as Proposition~\ref{prop:naive} predicts:
its worst-case ratio does not shrink with $n$, it is simply unbounded
below on adversarial instances and highly variable on random ones}\DIFdelend \DIFaddbegin \DIFadd{bound cannot do better than $\epsilon$-dependent cost. The inexpensive
completion step removes this failure on the family. Cost grows in
proportion to $1/\epsilon$ (Table~\ref{tab:epscost}), as predicted, and
is independent of $P$ (next subsection)}\DIFaddend .

\DIFdelbegin %DIFDELCMD < \subsection{Runtime: pseudo-polynomial versus polynomial in practice}
%DIFDELCMD < %%%
\DIFdelend \DIFaddbegin \subsection{Runtime: time-indexed versus value-indexed}\label{sec:runtime}
\DIFaddend 

Theorem~\ref{thm:dp}'s \DIFdelbegin \DIFdel{exact DP }\DIFdelend \DIFaddbegin \DIFadd{dynamic program }\DIFaddend runs in $O(nP)$ time, \DIFdelbegin \DIFdel{where $P=\sum_i
p_i$ can be }\DIFdelend \DIFaddbegin \DIFadd{with
$P=\sum_ip_i$ }\DIFaddend exponential in the input's bit length; \DIFdelbegin \DIFdel{Theorem~\ref{thm:fptas}'s FPTAS runs in time polynomial in $n$ and $1/\epsilon$, independent of
}\DIFdelend \DIFaddbegin \DIFadd{the FPTAS's table
does not depend on }\DIFaddend $P$. To see this\DIFdelbegin \DIFdel{gap empirically rather than only asymptotically, we
}\DIFdelend \DIFaddbegin \DIFadd{, }\DIFaddend fix $n=40$\DIFdelbegin \DIFdel{and }\DIFdelend \DIFaddbegin \DIFadd{, }\DIFaddend $\epsilon=0.1$\DIFdelbegin \DIFdel{and }\DIFdelend \DIFaddbegin \DIFadd{, draw
weights up to $10^6$ (so the scaling is active at every setting), }\DIFaddend vary
the range $p_{\max}$ from which \DIFdelbegin \DIFdel{processing }\DIFdelend \DIFaddbegin \DIFadd{dispatch }\DIFaddend times are drawn\DIFdelbegin \DIFdel{(so $P$ grows roughly linearly with
$p_{\max}$), timing }\DIFdelend \DIFaddbegin \DIFadd{, and time }\DIFaddend both
algorithms on \DIFdelbegin \DIFdel{5 }\DIFdelend \DIFaddbegin \DIFadd{$5$ }\DIFaddend random instances per setting.
\DIFdelbegin \DIFdel{Table~\ref{tab:runtime} reports the mean wall-clock time.
}\DIFdelend 

\begin{table}[h]
\DIFdelbeginFL %DIFDELCMD < \caption{Mean wall-clock time (milliseconds), $n=40$, $\epsilon=0.1$,
%DIFDELCMD < over 5 random instances per row, as the processing-time range
%DIFDELCMD < $p_{\max}$ grows (so $P=\sum_i p_i$ grows roughly proportionally).}%%%
\DIFdelendFL \DIFaddbeginFL \caption{Mean wall-clock time (milliseconds), $n=40$, $\epsilon=0.1$,
weights up to $10^6$, over 5 random instances per row, as $p_{\max}$
grows (so $P$ grows roughly proportionally).}\DIFaddendFL \label{tab:runtime}
\begin{tabular}{@{}lccccc@{}}
\toprule
$p_{\max}$ & 50 & 200 & 800 & 3{,}200 & 12{,}800 \\
\midrule
Exact DP (Theorem~\ref{thm:dp}) & \DIFdelbeginFL \DIFdelFL{4.7 }\DIFdelendFL \DIFaddbeginFL \DIFaddFL{0.75 }\DIFaddendFL & \DIFdelbeginFL \DIFdelFL{16.8 }\DIFdelendFL \DIFaddbeginFL \DIFaddFL{0.93 }\DIFaddendFL & \DIFdelbeginFL \DIFdelFL{65.3 }\DIFdelendFL \DIFaddbeginFL \DIFaddFL{2.75 }\DIFaddendFL & \DIFdelbeginFL \DIFdelFL{275.7 }\DIFdelendFL \DIFaddbeginFL \DIFaddFL{9.16 }\DIFaddendFL & \DIFdelbeginFL \DIFdelFL{1137.9 }\DIFdelendFL \DIFaddbeginFL \DIFaddFL{55.1 }\DIFaddendFL \\
FPTAS (Theorem~\ref{thm:fptas}) & \DIFdelbeginFL \DIFdelFL{8.3 }\DIFdelendFL \DIFaddbeginFL \DIFaddFL{1.18 }\DIFaddendFL & \DIFdelbeginFL \DIFdelFL{8.4 }\DIFdelendFL \DIFaddbeginFL \DIFaddFL{1.07 }\DIFaddendFL & \DIFdelbeginFL \DIFdelFL{7.8 }\DIFdelendFL \DIFaddbeginFL \DIFaddFL{1.17 }\DIFaddendFL & \DIFdelbeginFL \DIFdelFL{7.7 }\DIFdelendFL \DIFaddbeginFL \DIFaddFL{1.13 }\DIFaddendFL & \DIFdelbeginFL \DIFdelFL{7.6 }\DIFdelendFL \DIFaddbeginFL \DIFaddFL{1.64 }\DIFaddendFL \\
\bottomrule
\end{tabular}
\end{table}

\begin{figure}[h]
\centering
\includegraphics[width=0.68\textwidth]{figures/runtime_scaling.pdf}
\DIFdelbeginFL %DIFDELCMD < \caption{Table~\ref{tab:runtime} plotted on log-log axes: the exact
%DIFDELCMD < DP's pseudo-polynomial $O(nP)$ dependence on $p_{\max}$ is visibly
%DIFDELCMD < linear, while the FPTAS's $O(n^3/\epsilon)$ runtime---independent of
%DIFDELCMD < $p_{\max}$---stays flat.}%%%
\DIFdelendFL \DIFaddbeginFL \caption{Table~\ref{tab:runtime} on log-log axes: the exact DP's time
grows with $p_{\max}$ (roughly linearly or slightly faster once the
fixed overhead is exceeded), while the FPTAS's stays flat.}\DIFaddendFL \label{fig:runtime}
\end{figure}

The exact \DIFdelbegin \DIFdel{DP's runtime grows essentially linearly with $p_{\max}$, as
its $O(nP)$ bound predicts, increasing roughly $240$}\DIFdelend \DIFaddbegin \DIFadd{dynamic program grows about $74$}\DIFaddend -fold as $p_{\max}$ \DIFdelbegin \DIFdel{increases }\DIFdelend \DIFaddbegin \DIFadd{grows
}\DIFaddend $256$-fold \DIFdelbegin \DIFdel{; the FPTAS 's runtime stays flat , within
measurement noise, across the same range, since its table size
$O(n^3/\epsilon)$ never depends on $P$ at all}\DIFdelend \DIFaddbegin \DIFadd{(the smallest settings are dominated by fixed overhead),
whereas the FPTAS is flat up to noise}\DIFaddend .
At the smallest $p_{\max}$ \DIFdelbegin \DIFdel{tested here the exact DP is actually }\DIFdelend \DIFaddbegin \DIFadd{the exact program is }\DIFaddend faster in absolute
terms, \DIFdelbegin \DIFdel{since its (pseudo-polynomial but small) table is cheaper than
the FPTAS's fixed $O(n^3/\epsilon)$ overhead---a }\DIFdelend \DIFaddbegin \DIFadd{a }\DIFaddend reminder that an \DIFdelbegin \DIFdel{FPTAS's asymptotic advantage only pays off once the
instance's
}\DIFdelend \DIFaddbegin \DIFadd{approximation scheme pays off only when the
}\DIFaddend numerical magnitude is large\DIFdelbegin \DIFdel{enough, exactly the regime
Theorem~\ref{thm:hardness}'s weak-NP-hardness result says the exact DP
cannot handle efficiently.
}%DIFDELCMD < 

%DIFDELCMD < \subsection{Sensitivity to the accuracy parameter}
%DIFDELCMD < 

%DIFDELCMD < %%%
\DIFdel{Theorem~\ref{thm:fptas}'s table has $O(n^3/\epsilon)$ entries, so
shrinking $\epsilon$ should, in principle, grow both the accuracy and
the running time of the FPTAS. We fix $n=30$ and sweep $\epsilon \in
\{0.5, 0.3, 0.2, 0.1, 0.05, 0.01\}$ on the same 20 random instances at
each value, reporting the mean ratio to the true optimum and the mean
wall-clock time in Table~\ref{tab:epsilon}.
}%DIFDELCMD < 

%DIFDELCMD < \begin{table}[h]
%DIFDELCMD < \caption{FPTAS accuracy and runtime as $\epsilon$ shrinks, $n=30$, mean
%DIFDELCMD < over the same 20 random instances at each $\epsilon$.}\label{tab:epsilon}
%DIFDELCMD < \begin{tabular}{@{}lcccccc@{}}
%DIFDELCMD < \toprule
%DIFDELCMD < %%%
\DIFdelFL{$\epsilon$ }%DIFDELCMD < & %%%
\DIFdelFL{0.50 }%DIFDELCMD < & %%%
\DIFdelFL{0.30 }%DIFDELCMD < & %%%
\DIFdelFL{0.20 }%DIFDELCMD < & %%%
\DIFdelFL{0.10 }%DIFDELCMD < & %%%
\DIFdelFL{0.05 }%DIFDELCMD < & %%%
\DIFdelFL{0.01 }%DIFDELCMD < \\
%DIFDELCMD < \midrule
%DIFDELCMD < %%%
\DIFdelFL{Mean ratio to optimum }%DIFDELCMD < & %%%
\DIFdelFL{0.9999 }%DIFDELCMD < & %%%
\DIFdelFL{1.0000 }%DIFDELCMD < & %%%
\DIFdelFL{1.0000 }%DIFDELCMD < & %%%
\DIFdelFL{1.0000 }%DIFDELCMD < & %%%
\DIFdelFL{1.0000 }%DIFDELCMD < & %%%
\DIFdelFL{1.0000 }%DIFDELCMD < \\
%DIFDELCMD < %%%
\DIFdelFL{Mean time (ms) }%DIFDELCMD < & %%%
\DIFdelFL{2.5 }%DIFDELCMD < & %%%
\DIFdelFL{3.9 }%DIFDELCMD < & %%%
\DIFdelFL{3.9 }%DIFDELCMD < & %%%
\DIFdelFL{3.8 }%DIFDELCMD < & %%%
\DIFdelFL{4.2 }%DIFDELCMD < & %%%
\DIFdelFL{3.8 }%DIFDELCMD < \\
%DIFDELCMD < \bottomrule
%DIFDELCMD < \end{tabular}
%DIFDELCMD < \end{table}
%DIFDELCMD < 

%DIFDELCMD < \begin{figure}[h]
%DIFDELCMD < \centering
%DIFDELCMD < \includegraphics[width=0.95\textwidth]{figures/epsilon_sensitivity.pdf}
%DIFDELCMD < \caption{Table~\ref{tab:epsilon} plotted: accuracy (left) is
%DIFDELCMD < essentially flat at $1.0$ across the entire tested range of
%DIFDELCMD < $\epsilon$, while runtime (right) grows only mildly as $\epsilon$
%DIFDELCMD < shrinks two orders of magnitude---both axes with $\epsilon$
%DIFDELCMD < decreasing left to right.}\label{fig:epsilon}
%DIFDELCMD < \end{figure}
%DIFDELCMD < 

%DIFDELCMD < %%%
\DIFdel{Even at the loosest setting tested, $\epsilon=0.5$ (a $50\%$ worst-case
loss allowed), the FPTAS already recovers $99.99\%$ of the optimum on
average, and by $\epsilon=0.3$ it matches the optimum on every instance
in this sample; consistent with Example~\ref{ex:fptas} and
Section~\ref{sec:experiments}'s accuracy table, the worst-case bound
$(1-\epsilon)W^\ast$ is a safety margin that is rarely approached in
practice, not a description of typical behavior. The runtime grows only
mildly over this range at $n=30$ (a two-order-of-magnitude shrinkage in
$\epsilon$, from $0.5$ to $0.01$, produces roughly a $50\%$ increase in
running time, not the $50\times$ the $O(n^3/\epsilon)$ bound alone
would suggest as an upper bound), because the table's row dimension is
capped by $\sum_i w_i'$, which grows sub-linearly in $1/\epsilon$ once
the scaled weights $w_i' = \lfloor w_i/K\rfloor$ saturate at small
integer values; the
$O(n^3/\epsilon)$ bound is a worst-case ceiling; on these instances the effective table is considerably smaller}\DIFdelend \DIFaddbegin \DIFadd{. We stress that this comparison is between
a time-indexed and a value-indexed table, and that the dependence on the
weights has simply moved: a value-indexed }\emph{\DIFadd{exact}} \DIFadd{program would
depend on $\sum_iw_i$ in the same way}\DIFaddend .

\subsection{Robustness to the instance-generation regime}\DIFaddbegin \label{sec:robust}
\DIFaddend 

\DIFdelbegin \DIFdel{The accuracy comparison of Section~\ref{sec:experiments}.1 uses
uniformly distributed weights and a deadline horizon spanning the full
range $[1,\sum_i p_i]$. }\DIFdelend To check that \DIFdelbegin \DIFdel{its qualitative pattern }\DIFdelend \DIFaddbegin \DIFadd{the pattern of Table~\ref{tab:illustration} }\DIFaddend is not an
artifact of \DIFdelbegin \DIFdel{that specific choice, }\DIFdelend \DIFaddbegin \DIFadd{the generator }\DIFaddend we repeat the comparison at $n=30$ \DIFdelbegin \DIFdel{under two different regimes: }\DIFdelend \DIFaddbegin \DIFadd{($20$
instances each) under }\DIFaddend heavy-tailed weights (\DIFdelbegin \DIFdel{drawn from an
exponentialdistribution rather than uniform, so a few sites are far
more critical than the rest}\DIFdelend \DIFaddbegin \DIFadd{exponential, mean $20$}\DIFaddend ) and
tight deadlines (the \DIFdelbegin \DIFdel{deadline }\DIFdelend horizon narrowed to \DIFdelbegin \DIFdel{$0.5\sum_i p_i$, so on average only about half the sites
can be served regardless of order).
Table~\ref{tab:robustness} reports
the same three ratios as before under each regime.
}\DIFdelend \DIFaddbegin \DIFadd{$0.5\sum_ip_i$).
}\DIFaddend 

\begin{table}[h]
\DIFdelbeginFL %DIFDELCMD < \caption{Mean ratio to the true optimum under two alternative
%DIFDELCMD < instance-generation regimes, $n=30$, 20 random instances per
%DIFDELCMD < row.}%%%
\DIFdelendFL \DIFaddbeginFL \caption{Mean ratio to the true optimum under two alternative
instance-generation regimes, $n=30$, 20 random instances per row.
``Repair'' is the greedy repair of Algorithm~\ref{alg:repair}; ``Skip'' is
EDD with skipping; the FPTAS uses $\epsilon=0.1$.}\DIFaddendFL \label{tab:robustness}
\DIFdelbeginFL %DIFDELCMD < \begin{tabular}{@{}lccc@{}}
%DIFDELCMD < %%%
\DIFdelendFL \DIFaddbeginFL \begin{tabular}{@{}lcccc@{}}
\DIFaddendFL \toprule
Regime & FPTAS \DIFdelbeginFL \DIFdelFL{($\epsilon{=}0.1$) }\DIFdelendFL & Repair \DIFdelbeginFL \DIFdelFL{(M--H) }\DIFdelendFL \DIFaddbeginFL & \DIFaddFL{Skip }\DIFaddendFL & Naive EDD \\
\midrule
Heavy-tailed weights & 1.000 & \DIFdelbeginFL \DIFdelFL{0.998 }\DIFdelendFL \DIFaddbeginFL \DIFaddFL{0.999 }\DIFaddendFL & \DIFdelbeginFL \DIFdelFL{0.274 }\DIFdelendFL \DIFaddbeginFL \DIFaddFL{0.944 }& \DIFaddFL{0.537 }\DIFaddendFL \\
Tight deadlines ($0.5\times$ horizon) & 1.000 & 0.988 & \DIFdelbeginFL \DIFdelFL{0.026 }\DIFdelendFL \DIFaddbeginFL \DIFaddFL{0.799 }& \DIFaddFL{0.125 }\DIFaddendFL \\
\bottomrule
\end{tabular}
\end{table}

The \DIFdelbegin \DIFdel{FPTAS and the greedy repair heuristic are essentially unaffected :
both continue to recover almost all of the optimum, confirming that
Section~\ref{sec:experiments}.1's pattern is not an artifact of uniform
weights or a loose deadline horizon. The naive EDD baseline, by
contrast, gets }\emph{\DIFdel{worse}} %DIFAUXCMD
\DIFdel{under both regimes---down to $2.6\%$ of
optimum under tight deadlines---which is exactly what
Proposition~\ref{prop:naive} predicts: tightening deadlinesor
increasing weight heterogeneity both widen the gap between the single
early commitment naive EDD makes and the value it forecloses, since
}\DIFdelend \DIFaddbegin \DIFadd{greedy repair is essentially unaffected and EDD with skipping
degrades under tight deadlines, where }\DIFaddend there is less slack to absorb a
poor early \DIFdelbegin \DIFdel{choice and more value
concentrated in whichever sites get displaced by it .
}\DIFdelend \DIFaddbegin \DIFadd{commitment; naive EDD collapses to $12.5\%$ of the optimum,
largely because it also wastes time on sites that are already late.
(As in Section~\ref{sec:acc}, the FPTAS entries are $1.000$ partly
because the scaling is vacuous at weights $\le100$.)
}\DIFaddend 

\subsection{A real-world illustration: the 2018 Camp Fire}\label{sec:campfire}

The instances above are synthetic by \DIFdelbegin \DIFdel{design---MWHED }\DIFdelend \DIFaddbegin \DIFadd{design: MWHED }\DIFaddend is a new problem with
no \DIFdelbegin \DIFdel{existing }\DIFdelend benchmark set, and the \DIFdelbegin \DIFdel{robustness check just given
specifically needs controllable distributional properties that only a
generator, not a single fixed dataset, can provide}\DIFdelend \DIFaddbegin \DIFadd{experiments need controllable distributions}\DIFaddend .
To ground the model's scale \DIFdelbegin \DIFdel{and behavior }\DIFdelend in a real event \DIFdelbegin \DIFdel{, we also }\DIFdelend \DIFaddbegin \DIFadd{we }\DIFaddend build one small instance
\DIFdelbegin \DIFdel{directly }\DIFdelend from public records of the 2018 Camp Fire (Butte County, California), the
deadliest wildfire in California history (85 fatalities, over 18,000
structures destroyed). We \DIFdelbegin \DIFdel{emphasize which
numbers below are real measurements and which are }\DIFdelend \DIFaddbegin \DIFadd{separate real measurements from }\DIFaddend disclosed
modeling assumptions, \DIFdelbegin \DIFdel{rather than blending the two silently}\DIFdelend \DIFaddbegin \DIFadd{and we present the case as an }\emph{\DIFadd{illustration of
the model and its sensitivity}}\DIFadd{, not as an operational recommendation}\DIFaddend .

\textbf{Real inputs.} The fire's ignition point (a PG\&E transmission
tower near Pulga, CA) and its progression timeline---in particular, fire
reaching Concow approximately $52$ minutes after \DIFdelbegin \DIFdel{ignition}\DIFdelend \DIFaddbegin \DIFadd{the time zero of the
timeline}\DIFaddend , and spot fires igniting in Paradise approximately $71$
minutes after \DIFdelbegin \DIFdel{ignition---are as documented in }\DIFdelend \DIFaddbegin \DIFadd{it---are taken from }\DIFaddend a National Institute of Standards and
Technology forensic reconstruction \citep{nist2021campfire}. We take
four real, named communities threatened by the fire (Concow, Paradise,
Magalia, and Yankee Hill, Butte County) with their real geographic
coordinates and 2010 U.S.\ Census populations (the census immediately
preceding the fire, used as a proxy criticality weight $w_i$; later
counts reflect the \DIFdelbegin \DIFdel{fire's }\DIFdelend aftermath, not the population at risk). A real
regional fire-agency depot (Oroville, CA) anchors the dispatch side.

\textbf{Disclosed modeling assumptions.} \DIFdelbegin \DIFdel{Two quantities are not
directly measured. First, }\DIFdelend \DIFaddbegin \DIFadd{(i) }\DIFaddend $p_i$ is \DIFdelbegin \DIFdel{computed from }\DIFdelend \DIFaddbegin \DIFadd{twice }\DIFaddend the real
great-circle distance between the depot and \DIFdelbegin \DIFdel{each site (a real }\DIFdelend \DIFaddbegin \DIFadd{site $i$ (a }\DIFaddend lower bound on
\DIFdelbegin \DIFdel{true
}\DIFdelend road distance) divided by an assumed average response speed; we report
\DIFdelbegin \DIFdel{results at two illustrative speeds, }\DIFdelend $50$~km/h (conservative, winding mountain roads) and $80$~km/h
(\DIFdelbegin \DIFdel{an optimistichighway-speed upper
bound). Second, }\DIFdelend \DIFaddbegin \DIFadd{optimistic). (ii) }\DIFaddend Magalia and Yankee Hill have no \DIFdelbegin \DIFdel{directly documented fire-arrival }\DIFdelend \DIFaddbegin \DIFadd{documented arrival
}\DIFaddend timestamp in the source\DIFdelbegin \DIFdel{we use}\DIFdelend , so their \DIFdelbegin \DIFdel{deadlines $d_i$
}\DIFdelend \DIFaddbegin \DIFadd{arrival times }\DIFaddend are extrapolated: \DIFdelbegin \DIFdel{we compute the
real fire-spread }\DIFdelend \DIFaddbegin \DIFadd{the
spread }\DIFaddend rate implied by Concow's and Paradise's real \DIFdelbegin \DIFdel{distance-from-ignition and real documented }\DIFdelend \DIFaddbegin \DIFadd{distance from
ignition and real }\DIFaddend arrival time ($12.1$ and $14.2$~km/h\DIFdelbegin \DIFdel{respectively, averaging
}\DIFdelend \DIFaddbegin \DIFadd{, average
}\DIFaddend $13.1$~km/h) \DIFdelbegin \DIFdel{, then apply this averaged rate to Magalia's and Yankee
Hill's own real }\DIFdelend \DIFaddbegin \DIFadd{is applied to their own straight-line }\DIFaddend distance from the
ignition point\DIFdelbegin \DIFdel{. Table~\ref{tab:campfire}
reports the resulting instance and Table~\ref{tab:campfire-results} the outcome of running Algorithms~\ref{alg:exact}--\ref{alg:repair}
unmodified on it (code and full source citations:
}\texttt{\DIFdel{case\_study\_campfire.py}}%DIFAUXCMD
\DIFdel{) }\DIFdelend \DIFaddbegin \DIFadd{; this extrapolation is isotropic, whereas the real spread
was wind-driven and directional. (iii) The crew can leave at the time
zero of the timeline, with no detection or dispatch delay. (iv) The
}\emph{\DIFadd{arrival reading}} \DIFadd{of Remark~\ref{rem:reading} is used: a site is
protected if the crew }\emph{\DIFadd{arrives}} \DIFadd{before the fire, so the MWHED
deadline is $d_i=d_i^{\mathrm{haz}}+p_i/2$. Table~\ref{tab:campfire}
gives the instance}\DIFaddend .

\begin{table}[h]
\DIFdelbeginFL %DIFDELCMD < \caption{The Camp Fire instance. $p_i$ (round-trip dispatch minutes
%DIFDELCMD < from the Oroville depot) is shown at both response-speed assumptions;
%DIFDELCMD < $d_i$ (minutes after ignition) is documented for Concow and Paradise,
%DIFDELCMD < extrapolated for the other two (see text); $w_i$ is 2010 census
%DIFDELCMD < population.}%%%
\DIFdelendFL \DIFaddbeginFL \caption{The Camp Fire instance. $p_i$ is the round-trip occupancy from
the Oroville depot at the two response speeds; $d^{\mathrm{haz}}$ is the
fire-arrival minute (documented for Concow and Paradise, extrapolated
for the others); the MWHED deadline under the arrival reading is
$d_i=d^{\mathrm{haz}}+p_i/2$ (rounded), shown for both speeds; $w_i$ is
the 2010 census population.}\DIFaddendFL \label{tab:campfire}
\small
\DIFdelbeginFL %DIFDELCMD < \begin{tabular}{@{}lcccc@{}}
%DIFDELCMD < %%%
\DIFdelendFL \DIFaddbeginFL \begin{tabular}{@{}lccccccc@{}}
\DIFaddendFL \toprule
Site & $p_i$ (50\DIFdelbeginFL \DIFdelFL{km/h}\DIFdelendFL ) & $p_i$ (80\DIFdelbeginFL \DIFdelFL{km/h}\DIFdelendFL \DIFaddbeginFL \DIFaddFL{) }& \DIFaddFL{$d^{\mathrm{haz}}$ }& \DIFaddFL{$d_i$ (50}\DIFaddendFL ) & $d_i$ \DIFaddbeginFL \DIFaddFL{(80) }\DIFaddendFL & $w_i$ \\
\midrule
Concow & 63 & 39 & 52 & \DIFaddbeginFL \DIFaddFL{84 }& \DIFaddFL{72 }& \DIFaddendFL 710 \\
Magalia & 89 & 55 & 58 & \DIFaddbeginFL \DIFaddFL{102 }& \DIFaddFL{86 }& \DIFaddendFL 11{,}310 \\
Yankee Hill & 54 & 34 & 64 & \DIFaddbeginFL \DIFaddFL{91 }& \DIFaddFL{81 }& \DIFaddendFL 333 \\
Paradise & 70 & 44 & 71 & \DIFaddbeginFL \DIFaddFL{106 }& \DIFaddFL{93 }& \DIFaddendFL 26{,}218 \\
\bottomrule
\end{tabular}
\end{table}

\begin{figure}[h]
\centering
\includegraphics[width=0.75\textwidth]{figures/campfire_map.pdf}
\DIFdelbeginFL %DIFDELCMD < \caption{The real geography behind Table~\ref{tab:campfire}: the
%DIFDELCMD < depot, the fire's real ignition point, and the four real communities,
%DIFDELCMD < plotted from their real coordinates (dotted lines show each
%DIFDELCMD < community's real distance from the ignition point, the basis for the
%DIFDELCMD < documented and extrapolated $d_i$ values). The solid arrow is the
%DIFDELCMD < optimal dispatch, in both response-speed scenarios: straight to
%DIFDELCMD < Paradise, bypassing every other community.}%%%
\DIFdelendFL \DIFaddbeginFL \caption{The real geography behind Table~\ref{tab:campfire}: the depot,
the fire's ignition point, and the four communities, plotted from their
real coordinates (dotted lines: each community's real distance from the
ignition point, the basis for the documented and extrapolated arrival
times). The solid arrow, dispatch to Paradise, is in the optimum at both
response speeds; at $80$~km/h the optimum also serves Concow, first
(dashed arrow).}\DIFaddendFL \label{fig:campfire-map}
\end{figure}

\begin{table}[h]
\DIFdelbeginFL %DIFDELCMD < \caption{Outcome of each algorithm on the Camp Fire instance, at both
%DIFDELCMD < response-speed assumptions.}%%%
\DIFdelendFL \DIFaddbeginFL \caption{On-time weight (population) of each algorithm on the Camp Fire
instance, arrival reading; the sites served are in parentheses.}\DIFaddendFL \label{tab:campfire-results}
\DIFaddbeginFL \small
\DIFaddendFL \begin{tabular}{@{}lcc@{}}
\toprule
 & $50$~km/h & $80$~km/h \\
\midrule
Exact optimum (Thm.~\ref{thm:dp}) & \DIFdelbeginFL \DIFdelFL{$26{,}218$ }\DIFdelendFL \DIFaddbeginFL \DIFaddFL{26}{\DIFaddFL{,}}\DIFaddFL{218 }\DIFaddendFL (Paradise) & \DIFdelbeginFL \DIFdelFL{$26{,}218$ (}\DIFdelendFL \DIFaddbeginFL \DIFaddFL{26}{\DIFaddFL{,}}\DIFaddFL{928 (Concow, }\DIFaddendFL Paradise) \\
FPTAS, $\epsilon=0.1$ (Thm.~\ref{thm:fptas}) & \DIFdelbeginFL \DIFdelFL{$26{,}218$ }\DIFdelendFL \DIFaddbeginFL \DIFaddFL{26}{\DIFaddFL{,}}\DIFaddFL{218 }\DIFaddendFL & \DIFdelbeginFL \DIFdelFL{$26{,}218$ }\DIFdelendFL \DIFaddbeginFL \DIFaddFL{26}{\DIFaddFL{,}}\DIFaddFL{928 }\DIFaddendFL \\
Greedy repair \DIFdelbeginFL \DIFdelFL{(M--H style) }\DIFdelendFL & \DIFdelbeginFL \DIFdelFL{$26{,}218$ }\DIFdelendFL \DIFaddbeginFL \DIFaddFL{26}{\DIFaddFL{,}}\DIFaddFL{218 }\DIFaddendFL & \DIFdelbeginFL \DIFdelFL{$26{,}218$ }\DIFdelendFL \DIFaddbeginFL \DIFaddFL{26}{\DIFaddFL{,}}\DIFaddFL{218 (Paradise) }\\
\DIFaddFL{EDD with skipping }& \DIFaddFL{710 (Concow) }& \DIFaddFL{1}{\DIFaddFL{,}}\DIFaddFL{043 (Concow, Yankee Hill) }\DIFaddendFL \\
Naive EDD \DIFdelbeginFL \DIFdelFL{(no repair) }\DIFdelendFL & \DIFdelbeginFL \DIFdelFL{$0$ }\DIFdelendFL \DIFaddbeginFL \DIFaddFL{710 (Concow) }\DIFaddendFL & \DIFdelbeginFL \DIFdelFL{$710$ (Concow}\DIFdelendFL \DIFaddbeginFL \DIFaddFL{1}{\DIFaddFL{,}}\DIFaddFL{043 (Concow, Yankee Hill}\DIFaddendFL ) \\
\bottomrule
\end{tabular}
\end{table}

\DIFdelbegin \DIFdel{At the conservative }\DIFdelend \DIFaddbegin \textbf{\DIFadd{Results.}} \DIFadd{At }\DIFaddend $50$~km/h \DIFdelbegin \DIFdel{response speed, Concow and Magalia are
individually infeasible---even dispatched first and alone, a crew from
the Oroville depot cannot reach them before the fire does, a sobering
but real reflection of how far this fire's actual spread outran any
single-vehicle response. The optimal policy recognizes this
immediately and }\DIFdelend \DIFaddbegin \DIFadd{the optimum }\DIFaddend sends the crew \DIFdelbegin \DIFdel{directly }\DIFdelend \DIFaddbegin \DIFadd{straight }\DIFaddend to
Paradise, \DIFdelbegin \DIFdel{saving all
}\DIFdelend \DIFaddbegin \DIFadd{protecting its weight of }\DIFaddend $26{,}218$\DIFdelbegin \DIFdel{people at risk there; naive earliest-deadline dispatch, recall, tries Concow and Magalia first }\emph{\DIFdel{because}} %DIFAUXCMD
\DIFdel{their deadlines
are earlier, exhausts the time budget on two sites it can never
reach, and arrives }\DIFdelend \DIFaddbegin \DIFadd{; any detour first, even the
shortest, uses time Paradise's deadline cannot absorb. Deadline-order
rules go to Concow first because its deadline is earlier, protect only
its $710$ residents, and arrive }\DIFaddend at Paradise too late\DIFdelbegin \DIFdel{to help anyone---an
on-the-nose illustration of Proposition~\ref{prop:naive}'s unbounded
worst-case loss, not a constructed adversarial example. At the more
optimistic $80$}\DIFdelend \DIFaddbegin \DIFadd{. At $80$~km/h the
shorter trips leave room to serve Concow }\emph{\DIFadd{and then}} \DIFadd{Paradise, for
$26{,}928$; the weighted greedy repair misses the Concow detour and
returns $26{,}218$, $2.6\%$ less than the optimum, so even on this
four-site instance a heuristic with no guarantee is not exact, while the
FPTAS and the exact program agree. Under the stricter }\emph{\DIFadd{return
reading}} \DIFadd{(the crew must be back before the fire arrives) the answer
changes: at $50$}\DIFaddend ~km/h \DIFdelbegin \DIFdel{speed all four sites are individually reachable, yet the optimal policy }\DIFdelend \DIFaddbegin \DIFadd{Concow and Magalia are individually infeasible
(round trips of $63$ and $89$ minutes against arrival minutes $52$ and
$58$), although the crew would }\DIFaddend \emph{\DIFdelbegin \DIFdel{still}\DIFdelend \DIFaddbegin \DIFadd{reach}\DIFaddend } \DIFdelbegin \DIFdel{goes straight to Paradise alone rather than visiting Concow, Magalia, or Yankee Hill first: any
detour through the
smaller communities consumes time Paradise's
deadline cannot absorb. Naive EDD, dispatching earliest-deadline-first
regardless, manages to save only Concow 's $710$ residents and misses
Paradise entirely. Both the exact DP and the FPTAS recover the true
optimum }\DIFdelend \DIFaddbegin \DIFadd{them in $31.5$ and
$44.5$ minutes; the optimum is Paradise alone at both speeds. This is
why Remark~\ref{rem:reading} asks the user to state the reading.
}

\textbf{\DIFadd{Sensitivity to the uncertain inputs.}} \DIFadd{The inputs are uncertain
in several ways: the response speed, the arrival minutes (documented
ones carry reconstruction error, extrapolated ones more), and the
dispatch delay. We draw $5{,}000$ scenarios with speed uniform on
$[40,90]$ km/h, arrival minutes scaled by independent factors uniform
on $[0.85,1.15]$ (documented sites) or $[0.70,1.30]$ (extrapolated
sites), and a departure delay uniform on $[0,20]$ minutes, and solve each
exactly. Paradise is in the optimum in $100\%$ of the draws; the optimal
set is Paradise alone in $65.3\%$, Concow with Paradise in $18.4\%$,
Yankee Hill with Paradise in $13.9\%$ and Magalia with Paradise }\DIFaddend in
\DIFdelbegin \DIFdel{every case}\DIFdelend \DIFaddbegin \DIFadd{$2.3\%$. The plan that is optimal for the nominal $80$~km/h instance
(Concow, then Paradise) is fully on time in only $20.7\%$ of the draws
and retains on average $26.6\%$ of the drawn optimum (fifth percentile
$2.7\%$), because a delay or a slower road lets the Concow detour make
Paradise late. The simpler plan ``Paradise only'' is on time in all
draws and retains $98.6\%$ of the drawn optimum on average (fifth
percentile $97.4\%$). The nominal optimum is therefore brittle}\DIFaddend , and the
\DIFdelbegin \DIFdel{greedy repair heuristic matches them
here as well.
This single real }\DIFdelend \DIFaddbegin \DIFadd{brittleness comes from adding a low-weight site ahead of a dominant one;
this is the kind of effect that a model with uncertain parameters should
capture (Section~\ref{sec:disc-stochastic}).
}

\DIFadd{This single }\DIFaddend instance cannot substitute for the synthetic \DIFdelbegin \DIFdel{sweep}\DIFdelend \DIFaddbegin \DIFadd{study}\DIFaddend 's
statistical evidence, \DIFdelbegin \DIFdel{but it shows the same
qualitative pattern is not a synthetic-data artifact, and the $50$~km/h scenario's already-infeasible sites are exactly the kind of
situation Section~\ref{sec:disc-multi}'s discussion of multiple
vehicles is motivated by: a single depot's single crew has a hard
limit on what it can save, regardless of how it is scheduled}\DIFdelend \DIFaddbegin \DIFadd{and its outputs depend on the assumptions above;
its role is to show the model, the algorithms, and the reading and
uncertainty issues on real geography}\DIFaddend .

% ══════════════════════════════════════════════════════════════════
\section{Discussion and extensions}\label{sec:discussion}

The results of Section~\ref{sec:theory} settle the single-depot,
single-vehicle, deterministic-deadline case completely. Three natural
directions relax one assumption at a time; we discuss, for each, what
carries over from our analysis and what specifically breaks, rather
than simply listing them as unexplored future work.

\subsection{Multiple depots or vehicles}\label{sec:disc-multi}

Suppose $m>1$ vehicles are available, each \DIFdelbegin \DIFdel{able to serve a disjoint
subset of
sites, still under the same round-trip dispatch-time and
deadline structure. Call this }\DIFdelend \DIFaddbegin \DIFadd{serving its own sequence of
round trips, as in }\DIFaddend MWHED-$m$ \DIFdelbegin \DIFdel{(MWHED itself is MWHED-$1$) .
Two facts are immediate.
First, }\DIFdelend \DIFaddbegin \DIFadd{of Section~\ref{sec:equalp}. Four facts
describe what changes.
}

\emph{\DIFadd{Hardness persists for every fixed $m$.}} \DIFadd{Take a common-deadline
instance with deadline $D$ and total weight $\Omega=\sum_iw_i$ (the
hard instances of Theorem~\ref{thm:hardness}) and add $m-1$ ``blocker''
sites with $p=d=D$ and weight $\Omega+1$. A blocker occupies a vehicle
for all of $[0,D]$: if anything is served on that vehicle before it, the
blocker is late, and if anything is served after it, that site is late.
So each served blocker uses a vehicle exclusively, dropping a blocker
costs $\Omega+1$ and can free a vehicle worth at most $\Omega$, and an
optimum serves all $m-1$ blockers and solves the original instance on the
one remaining vehicle. Hence }\DIFaddend MWHED-$m$ is NP-hard for every fixed \DIFdelbegin \DIFdel{$m\ge 1$: the $m=1$ instance embeds into any MWHED-}\DIFdelend $m$\DIFdelbegin \DIFdel{instance by
adding $m-1$ vehicles with no sites to serve, so Theorem~\ref{thm:hardness}'s
hardness persists.Second, }\DIFdelend \DIFaddbegin \DIFadd{.
(Simply adding idle vehicles to an instance would not do: extra vehicles
can serve sites, which changes the problem.)
}

\emph{\DIFadd{The exact dynamic program grows with $m$.}} \DIFaddend Theorem~\ref{thm:dp}'s
\DIFdelbegin \DIFdel{dynamic program does
}\emph{\DIFdel{not}} %DIFAUXCMD
\DIFdel{extend at the same complexity: the natural generalization
tracks a completion-time state per vehicle}\DIFdelend \DIFaddbegin \DIFadd{recursion extends by tracking a completion time per vehicle. Assign sites
in deadline order (Lemma~\ref{lem:edd}, applied on each vehicle) and let
$f(i,t_1,\dots,t_m)$ be the maximum on-time weight using sites
$1,\dots,i$ when vehicle $k$ has accumulated exactly $t_k$ occupancy;
site $i$ may be appended to any vehicle $k$ with $t_k\le d_i$}\DIFaddend ,
\DIFdelbegin \DIFdel{giving a table of size
$O(n \cdot P^m)$ rather than $O(nP)$---pseudo-polynomial }\DIFdelend \DIFaddbegin \begin{multline*}
\DIFadd{f(i,t_1,\dots,t_m)=\max\Bigl(f(i-1,t_1,\dots,t_m),}\\
\DIFadd{\max_{k:\,p_i\le t_k\le d_i}
\bigl(f(i-1,t_1,\dots,t_k-p_i,\dots,t_m)+w_i\bigr)\Bigr).
}\end{multline*}
\DIFadd{The table has $O(nP^m)$ cells, pseudo-polynomial }\DIFaddend for fixed $m$ \DIFdelbegin \DIFdel{,
}\DIFdelend but with
an exponent that grows with the number of vehicles\DIFdelbegin \DIFdel{, not merely
a constant-factor slowdown. Whether a smarter exact algorithm avoids
this blow-up, or whether the $m^{\text{th}}$ power is
unavoidable , is open ; this connects MWHED-$m$ to }\DIFdelend \DIFaddbegin \DIFadd{. Whether this is
unavoidable is open here; }\DIFaddend the parallel-machine scheduling literature
\DIFdelbegin \DIFdel{'s own maximum-weighted-on-time-jobs problem , whose
complexity landscape beyond a single machine is }\DIFdelend \DIFaddbegin \DIFadd{treats the problem as }\DIFaddend considerably less settled than the single-machine
case \DIFdelbegin \DIFdel{we resolve here
\mbox{%DIFAUXCMD
\citep{lenstra1977}}\hskip0pt%DIFAUXCMD
. Recent evidence confirms this: \mbox{%DIFAUXCMD
\citet{heegermolter2025}
}\hskip0pt%DIFAUXCMD
show that even the parallel-machine analogue of our }\emph{\DIFdel{tractable}}
%DIFAUXCMD
\DIFdel{equal-cost special case---identical processing times across identical
machines, $P|r_j, p_j{=}p|\sum w_jU_j$---is already NP-hard, and $W[2]$-hard
parameterized by the number of machines, once release dates are
allowed. Translated into our setting, this suggests that
Theorem~\ref{thm:equalp}'s tractability for MWHED-$1$ is unlikely to
survive as cleanly into MWHED-}\DIFdelend \DIFaddbegin \DIFadd{\mbox{%DIFAUXCMD
\citep{lenstra1977}}\hskip0pt%DIFAUXCMD
. We conjecture, but have not verified, that
value-scaling extends to an }\DIFaddend $m$\DIFdelbegin \DIFdel{for $m\ge 2$, even restricted to the
equal-dispatch-cost case---the multi-vehicle relaxation appears to cost more than just an exponent in the DP's table size.
We conjecture the FPTAS of Theorem~\ref{thm:fptas}
extends to }\DIFdelend \DIFaddbegin \DIFadd{-dimensional state.
}

\emph{\DIFadd{Equal dispatch times stay tractable.}} \DIFadd{Theorem~\ref{thm:equalp}
solves }\DIFaddend MWHED-$m$ \DIFdelbegin \DIFdel{for fixed $m$ via the same value-scaling idea
applied to an }\DIFdelend \DIFaddbegin \DIFadd{with equal dispatch times in $O(n\log n)$ time for
every }\DIFaddend $m$\DIFdelbegin \DIFdel{-dimensional state (the scaled-weight axis is untouched }\DIFdelend \DIFaddbegin \DIFadd{. This is the multi-vehicle analogue of the single-vehicle
special case and holds without release dates.
}

\emph{\DIFadd{Release dates break it.}} \DIFadd{\mbox{%DIFAUXCMD
\citet{heegermolter2025} }\hskip0pt%DIFAUXCMD
show that the
unweighted parallel-machine problem with identical processing times
}\emph{\DIFadd{and release dates}}\DIFadd{, $P|r_j,p_j{=}p|\sum U_j$, is NP-hard and
$W[2]$-hard parameterized }\DIFaddend by the number of \DIFdelbegin \DIFdel{vehicles), but we have
not verified this in detail and do not claim it as a result here.
}%DIFDELCMD < 

%DIFDELCMD < %%%
\DIFdel{To see concretely where the $P^m$ term comes from: assign sites to
vehicles greedily by deadline as in Lemma~\ref{lem:edd}, but now track, for each vehicle $k=1,\dots,m$, its own cumulative dispatch time
$t_k$. The natural generalization of $f(i,t)$ in Theorem~\ref{thm:dp}'s
proof becomes $f(i, t_1,\dots,t_m)$: the
maximum on-time weight
achievable using sites $1,\dots,i$, where vehicle $k$ has accumulated
exactly $t_k$ dispatch time . Site $i$ can be assigned to whichever
vehicle $k$ currently has $t_k + p_i \le d_i$, giving the recursion
}\begin{align*}
\DIFdel{f(i, t_1,\dots,t_m) = \max\Bigl(}&\DIFdel{f(i-1,t_1,\dots,t_m),}\\
&\DIFdel{\max_{k :\, t_k \ge p_i,\, t_k \le d_i} \bigl(f(i-1, t_1,\dots,t_k -
p_i,\dots,t_m) + w_i\bigr)\Bigr),
}\end{align*}%DIFAUXCMD
\DIFdel{a direct generalization of \eqref{eq:dp} with $m$ time axes instead of
one. Since each axis ranges over $0,\dots,P$, the table has
$O(P^m)$ cells per site, $O(nP^m)$ overall---exactly the blow-up
claimed above, and the reason a genuinely new idea, not just a bigger
table, would be needed to bring MWHED-$m$'s exact algorithm back to a
complexity independent of $m$ in the exponent}\DIFdelend \DIFaddbegin \DIFadd{machines. Release dates have
no counterpart in MWHED as defined here, since every vehicle is at the
depot at time $0$ and every site can be served from the start. They
would arise if sites become reachable, or crews available, only at given
times (for instance, a road cleared by a firebreak). The equal-cost
tractability is therefore robust to adding vehicles but not to adding
release dates, and extending the model in that direction is a natural
open problem}\DIFaddend .

\subsection{A relocating depot or growing-radius hazard}\label{sec:disc-relocate}

Definition~\ref{def:mwhed} takes each site's round-trip dispatch time
$p_i$ as a fixed constant, appropriate when the depot is stationary. If
the depot itself can relocate between dispatches---a mobile command
post following a fire line, for instance---$p_i$ becomes a function of
the depot's current position, which in turn depends on the dispatch
order already chosen. Lemma~\ref{lem:edd}'s exchange argument relies
critically on $p_i$ being order-independent (Section~\ref{sec:model}'s
observation that every visit is a self-contained round trip); with a
relocating depot this fails, and the problem acquires genuine routing
structure on top of its scheduling structure. We expect this variant to
be at least as hard as MWHED (it contains MWHED as the special case of
a depot that chooses not to move) and likely strictly harder, since the
order in which sites are visited now affects future dispatch costs, not
just which deadlines are met.

It is worth being precise about what ``growing-radius hazard'' does
and does not already mean for a \emph{stationary} depot, since the two
relaxations named in this subsection's title are not the same
complication. If the hazard simply expands outward from a fixed origin
at a fixed, known rate, its arrival time at each site is already fully
captured by the deadline $d_i$ in Definition~\ref{def:mwhed}---this is
exactly the reduction described in Section~\ref{sec:intro}, and
requires no extension to our model at all. The genuine second
relaxation, distinct from a relocating depot, arises if the spreading
hazard itself degrades the road or access network the depot's vehicle
uses---for instance, a wildfire making a route impassable once it
crosses that route, or rising floodwater cutting off a bridge---so that
$p_i$ is not merely unknown but a function of \emph{when} site $i$ is
reached, since a later attempt may have to detour around
hazard-affected infrastructure that an earlier attempt would not have.
This again breaks Lemma~\ref{lem:edd}'s order-independence assumption,
but through the road network's degradation rather than the depot's own
movement, and the two effects would compound if both are present
simultaneously.

\DIFdelbegin %DIFDELCMD < \subsection{Uncertain hazard arrival times}%%%
\DIFdelend \DIFaddbegin \subsection{Uncertain parameters}\DIFaddend \label{sec:disc-stochastic}

Definition~\ref{def:mwhed} takes \DIFdelbegin \DIFdel{each deadline }\DIFdelend \DIFaddbegin \DIFadd{the dispatch times $p_i$, the
deadlines }\DIFaddend $d_i$ \DIFaddbegin \DIFadd{and the weights $w_i$ }\DIFaddend as known exactly. \DIFdelbegin \DIFdel{If $d_i$ is instead a random variable (as it would be if derived from a stochastic hazard-spread simulation), the natural analogue replaces the
hard feasibility indicator $\mathbb 1[C_i(\sigma)\le d_i]$ in
\eqref{eq:objective} with $\Pr[C_i(\sigma) \le d_i]$, turning MWHED into
a chance-constrained or expected-value scheduling problem. Lemma~\ref{lem:edd}'s
proof does not survive this change unmodified: the exchange argument
compares two fixed completion times against two fixed deadlines, and with random deadlines the ordering that maximizes }\DIFdelend \DIFaddbegin \DIFadd{In practice
each is an estimate: arrival times come from a fire-spread, flood or
plume simulator and carry its forecast error; occupancies depend on road
conditions that the hazard itself may degrade; weights are proxies for
population or asset exposure. The optimal solutions and the structural
statements above are conditional on these inputs, and the case study of
Section~\ref{sec:campfire} shows that the nominal optimum can be brittle
under plausible perturbations. We discuss three ways to model the
uncertainty, and which of our results survive each.
}

\emph{\DIFadd{Scenario-based (stochastic) models.}} \DIFadd{Let uncertain parameters
take values $\xi_s=(p^s,d^s,w^s)$ with probabilities $\pi_s$ for
$s=1,\dots,S$, and let the dispatch order $\sigma$ be chosen before the
scenario is revealed. The natural objective is the }\DIFaddend expected on-time
weight \DIFdelbegin \DIFdel{need not sort sites by any single deterministic statistic of
$d_i$,
for the same reason that earliest-due-date sequencing is not
generally optimal under stochastic due dates in the classical
scheduling literature}\DIFdelend \DIFaddbegin \DIFadd{$\sum_s\pi_s\sum_iw_i^s\,\mathbb 1[C_i^s(\sigma)\le d_i^s]$,
which for $S=1$ is MWHED and is therefore NP-hard. Lemma~\ref{lem:edd}
holds scenario by scenario but not simultaneously: no single order sorts
all scenarios' deadlines, and the exchange argument fails}\DIFaddend . A small
example \DIFdelbegin \DIFdel{makes this concrete: two }\DIFdelend \DIFaddbegin \DIFadd{shows that sorting by expected deadline can be nearly twice as
bad as the optimum. Two }\DIFaddend sites with $p=1$\DIFdelbegin \DIFdel{each, }\DIFdelend \DIFaddbegin \DIFadd{: }\DIFaddend site $A$ with $w_A=1$ and
\DIFdelbegin \DIFdel{a deterministic deadline $d_A=2$}\DIFdelend \DIFaddbegin \DIFadd{deadline $2$}\DIFaddend , and site $B$ with $w_B=10$ and a random deadline \DIFdelbegin \DIFdel{$d_B=1$ or $d_B=100$, each }\DIFdelend \DIFaddbegin \DIFadd{equal to
$1$ or $100$ }\DIFaddend with probability $\tfrac12$ \DIFdelbegin \DIFdel{, so $\mathbb E[d_B]=50.5 >
\mathbb E[d_A]=2$. Sorting by increasing $\mathbb E[d_i]$ dispatches
}\DIFdelend \DIFaddbegin \DIFadd{each, so $\mathbb E[d_B]=50.5>
\mathbb E[d_A]=2$. Serving }\DIFaddend $A$ first \DIFdelbegin \DIFdel{: $A$ completes at $1\le 2$ (always on time) , and $B$
completes at $2$, on time only when $d_B=100$, giving expected value
$1 + 10\cdot\tfrac12 = 6$. Dispatching }\DIFdelend \DIFaddbegin \DIFadd{(the expected-deadline order) gives
$1+10\cdot\tfrac12=6$; serving }\DIFaddend $B$ \DIFdelbegin \DIFdel{first instead: }\DIFdelend \DIFaddbegin \DIFadd{first gives $10+1=11$, because }\DIFaddend $B$\DIFdelbegin \DIFdel{completes
at $1$, on time whenever $d_B\ge 1$, which holds for }\emph{\DIFdel{both}}
%DIFAUXCMD
\DIFdel{possible values of $d_B$, and $A$ completes at $2\le 2$, always on
time, giving expected value $10 + 1 = 11$---nearly double, by visiting
the site with the }\emph{\DIFdel{larger}} %DIFAUXCMD
\DIFdel{expected deadline first , because $d_B$}\DIFdelend 's
tight realization \DIFdelbegin \DIFdel{($d_B=1$) is only avoidable }\DIFdelend \DIFaddbegin \DIFadd{is avoidable only }\DIFaddend by acting immediately\DIFdelbegin \DIFdel{,
information the expected value alone discards. This would couple
MWHED to exactly the
kind of sequential-decision and statistical
monitoring questions studied elsewhere in the stochastic
vehicle-routing literature, rather than to a direct generalization of Theorems~\ref{thm:dp}--\ref{thm:fptas}}\DIFdelend \DIFaddbegin \DIFadd{. This is the
setting of two-stage stochastic programming, in which a first-stage plan
is evaluated against scenario-dependent recourse; evacuation planning
under uncertain demand and network damage
\mbox{%DIFAUXCMD
\citep{wang2020evac} }\hskip0pt%DIFAUXCMD
and facility location under disruptions
\mbox{%DIFAUXCMD
\citep{koca2023} }\hskip0pt%DIFAUXCMD
are close relatives in the disaster-response
literature. A sample-average formulation of the expected-weight
objective is a mixed-integer program whose combinatorial core is
MWHED with scenario-dependent deadlines; a Benders-type decomposition by
scenario is natural, and exploiting Lemma~\ref{lem:edd} within a scenario
to generate cuts is, we believe, a promising direction that we have not
pursued.
}

\emph{\DIFadd{Robust feasibility.}} \DIFadd{Feasibility of a set of sites is monotone in
the parameters: it is preserved when $d_i$ increases or $p_i$
decreases. A set that is on time for }\emph{\DIFadd{every}} \DIFadd{realization in a box
$[p_i^-,p_i^+]\times[d_i^-,d_i^+]$ is therefore exactly a feasible set
of the MWHED instance $(p^+,d^-,w)$, and every result above applies to
that worst-case instance. This is cheap but conservative; in the case
study the worst corner of the sampled ranges (slow roads, early arrival,
a $20$-minute delay) makes even Paradise individually infeasible, so
instance-wise robustness would return nothing. Chance-constrained or
scenario-averaged objectives are the useful formulations there.
}

\emph{\DIFadd{Calibration and sensitivity.}} \DIFadd{Because the inputs are data-driven,
reporting how the optimum moves with them, as in
Section~\ref{sec:campfire}, is part of using the model. The case study
suggests a practical rule that scenario models would formalize: when one
site dominates the weight, protect it first and add other sites only if
they remain feasible under the plausible range of delays}\DIFaddend .

% ══════════════════════════════════════════════════════════════════
\section{Conclusion}\label{sec:conclusion}

We have \DIFdelbegin \DIFdel{given a complete complexity and algorithmic characterization of
}\DIFdelend \DIFaddbegin \DIFadd{studied }\DIFaddend Minimum Weighted Hazard-Exposure Dispatch, a scheduling
problem that arises \DIFdelbegin \DIFdel{naturally }\DIFdelend whenever a single depot must choose a dispatch order
against a spreading hazard with \DIFdelbegin \DIFdel{known, }\DIFdelend site-specific arrival deadlines. The
problem is exactly the classical weighted on-time-jobs \DIFdelbegin \DIFdel{scheduling
problem
}\DIFdelend \DIFaddbegin \DIFadd{problem
$1\|\sum w_jU_j$ }\DIFaddend in transportation language\DIFdelbegin \DIFdel{, and is weakly NP-hard already in
its simplest heterogeneous case (}\DIFdelend \DIFaddbegin \DIFadd{; consequently its weak
NP-hardness (already at }\DIFaddend equal deadlines, \DIFdelbegin \DIFdel{Theorem
\ref{thm:hardness}); an exact }\DIFdelend \DIFaddbegin \DIFadd{by a direct reduction from
}\textsc{\DIFadd{Partition}}\DIFadd{), the }\DIFaddend pseudo-polynomial dynamic program\DIFdelbegin \DIFdel{(Theorem~\ref{thm:dp}) and an FPTAS }\DIFdelend \DIFaddbegin \DIFadd{, and the
approximation scheme are classical, and we have presented them with
complete proofs rather than as new techniques. What we believe the
transportation framing adds is a precise account of }\emph{\DIFadd{which
heterogeneity makes dispatch hard}}\DIFadd{: the problem is polynomial exactly
when the dispatch times or the weights are constant and weakly NP-hard
otherwise }\DIFaddend (Theorem~\DIFdelbegin \DIFdel{\ref{thm:fptas})give
planners exact and provably near-exact tools regardless; and a second, structurally distinct special case---equal dispatch cost, Theorem
\ref{thm:equalp}---remains }\DIFdelend \DIFaddbegin \DIFadd{\ref{thm:classification}), heterogeneous hazard
deadlines never cause hardness, and the equal-dispatch-time case remains
}\DIFaddend exactly solvable in \DIFdelbegin \DIFdel{near-linear time ,
delineating precisely which combination of network heterogeneity
causes the difficulty. These guarantees are not merely convenient: the
natural do-nothing fallback of dispatching sites in deadline order with no reconsideration has }\DIFdelend \DIFaddbegin \DIFadd{$O(n\log n)$ time for any number of identical
vehicles (Theorem~\ref{thm:equalp}), though not once release dates are
added. Natural heuristics---deadline order with or without skipping, and
a weighted greedy repair---have }\DIFaddend unbounded worst-case \DIFdelbegin \DIFdel{loss relative to the
optimum (Proposition~\ref{prop:naive}), so some form of optimization is necessary for any performance guarantee at all, and our numerical
illustration confirms this qualitative gap persists }\DIFdelend \DIFaddbegin \DIFadd{ratio, so a
guarantee requires a rule that reconsiders earlier commitments in light
of both weights and dispatch times; the dynamic program and the FPTAS
are two such rules, not the only ones. The FPTAS's analysis is tight to
within a factor $1-\epsilon^2$, although }\DIFaddend on random instances \DIFdelbegin \DIFdel{of realistic size, alongside a weighted Moore--Hodgson-style
greedy repair heuristic that---while carrying no proven guarantee---
recovers substantially more of the optimum than doing nothing at all}\DIFdelend \DIFaddbegin \DIFadd{it is
essentially exact; this is only visible once the instance's weights are
large enough for its rounding to be active}\DIFaddend .

\DIFdelbegin \DIFdel{Section~\ref{sec:discussion} examines three natural relaxations---multiple
depots or vehicles, a relocating depot, and uncertain hazard-arrival
times---identifying precisely which part of our analysis (the exchange argument, the
dynamic program's state space, or both) each
relaxation breaks, rather than leaving them as an unexamined list. }%DIFDELCMD < 

%DIFDELCMD < %%%
\DIFdel{More broadly, MWHED is a small, self-contained illustration of a pattern that recurs across combinatorial optimization: a
transportation-network reparametrization of }\DIFdelend \DIFaddbegin \DIFadd{The numerical study and the case study carry two lessons beyond the
theorems. The reading of the deadline (the crew must arrive before the
hazard, or return before it) changes the answer on the same data and
must be stated. And the optimum of a model with estimated inputs can be
brittle: on the Camp Fire instance the nominal optimum serves a small
site ahead of the dominant one and fails under modest delays, whereas
protecting the dominant site first is robust. Section~\ref{sec:discussion}
discusses how scenario-based and robust formulations would capture this,
and which parts of the analysis survive: Lemma~\ref{lem:edd} holds
within a scenario but not across scenarios, which is what makes the
stochastic version a genuinely different problem. Natural next steps are
a stochastic or chance-constrained MWHED with a decomposition that uses
the earliest-deadline lemma within scenarios, release dates and
heterogeneous vehicles, and }\DIFaddend a \DIFdelbegin \DIFdel{decades-old scheduling
problem reveals both that the classical hardness and algorithmic
machinery transfer over essentially unchanged (Theorems~\ref{thm:hardness}--\ref{thm:fptas}), and that the specific structural question of which restriction
restores tractability (Theorem~\ref{thm:equalp}) has a clean,
network-interpretable answer---sites equidistant from the depot---that
was not visible in the original scheduling formulation's language of
abstract jobs and machines. We see thisas the paper's main
methodological contribution alongside its specific results: reframing
a classical scheduling question in transportation terms did not just
relabel it, it surfaced a distinction (dispatch-cost heterogeneity
versus deadline heterogeneity) that determines tractability, and that
distinction is exactly what Section~\ref{sec:discussion}'s three
extensions each perturb in a different way}\DIFdelend \DIFaddbegin \DIFadd{relocating depot}\DIFaddend .

\begin{appendices}

\section{Exhaustive verification of the running example}\label{app:exhaustive}

Examples~\ref{ex:running}--\ref{ex:matroid} rely on
Example~\ref{ex:running}'s instance being small enough to check by
exhaustive search, not only by the dynamic program itself. Since
Lemma~\ref{lem:edd} lets us restrict to dispatching each candidate
subset $S$ in increasing-deadline order, Table~\ref{tab:exhaustive}
lists all $2^4=16$ subsets of Example~\ref{ex:running}'s four sites
($p=(2,3,1,4)$, $d=(4,6,2,9)$, $w=(5,8,3,10)$), each dispatched in that
order, together with its total dispatch time and feasibility. Every
feasible subset's weight is at most $23$, and $\{1,2,4\}$ attains it,
confirming $W^\ast=23$ independently of Algorithm~\ref{alg:exact}.

\begin{table}[h]
\caption{All $16$ subsets of Example~\ref{ex:running}'s running
instance, dispatched in increasing-deadline order. ``Sum $p$'' is the
subset's total dispatch time; a subset is feasible iff every member
completes on time in this order (Lemma~\ref{lem:edd}).}\label{tab:exhaustive}
\small
\begin{tabular}{@{}lccc@{}}
\toprule
Subset $S$ & Order & Sum $p$ & Feasible? \quad $W(S)$ \\
\midrule
$\emptyset$ & --- & 0 & yes \quad 0 \\
$\{1\}$ & 1 & 2 & yes \quad 5 \\
$\{2\}$ & 2 & 3 & yes \quad 8 \\
$\{3\}$ & 3 & 1 & yes \quad 3 \\
$\{4\}$ & 4 & 4 & yes \quad 10 \\
$\{1,2\}$ & 1,2 & 5 & yes \quad 13 \\
$\{1,3\}$ & 3,1 & 3 & yes \quad 8 \\
$\{1,4\}$ & 1,4 & 6 & yes \quad 15 \\
$\{2,3\}$ & 3,2 & 4 & yes \quad 11 \\
$\{2,4\}$ & 2,4 & 7 & yes \quad 18 \\
$\{3,4\}$ & 3,4 & 5 & yes \quad 13 \\
$\{1,2,3\}$ & 3,1,2 & 6 & yes \quad 16 \\
$\{1,2,4\}$ & 1,2,4 & 9 & yes \quad \textbf{23} \\
$\{1,3,4\}$ & 3,1,4 & 7 & yes \quad 18 \\
$\{2,3,4\}$ & 3,2,4 & 8 & yes \quad 21 \\
$\{1,2,3,4\}$ & 3,1,2,4 & 10 & no (site 4 completes at $10>d_4{=}9$) \\
\bottomrule
\end{tabular}
\end{table}

Only the full set $\{1,2,3,4\}$ is infeasible, and it is infeasible for
exactly the reason Example~\ref{ex:running} describes: dispatching all
four sites, even in the best possible (deadline) order, makes site $4$
complete one time unit after its deadline. Every other subset is
feasible, so the optimization reduces to choosing, among the $15$
feasible subsets, the one of largest weight---which Table~\ref{tab:exhaustive}
confirms is $\{1,2,4\}$, at $W=23$, exactly matching
Example~\ref{ex:running}, Example~\ref{ex:dp}'s dynamic-program table,
and Example~\ref{ex:fptas}'s FPTAS instantiation.

\section{Detailed derivation of the FPTAS's approximation ratio}\label{app:fptas-detail}

Theorem~\ref{thm:fptas}'s proof compresses the approximation-ratio
derivation into a single chain of inequalities. We spell out each step
separately here, for a reader who wants to check the bound without
re-deriving it.

Recall the setup: $w_{\max} = \max_i w_i$, $K = \max(1, \epsilon
w_{\max}/n)$, $w_i' = \lfloor w_i/K \rfloor$, $S^\ast$ is the true
optimal subset with $W^\ast = \sum_{i \in S^\ast} w_i$, and $\hat S$ is
the subset the FPTAS returns (the feasible subset of largest scaled
value $v$ with $g(n,v)<\infty$). As in the main proof, we split on
which argument of $K=\max(1,\epsilon w_{\max}/n)$ attains the maximum;
the case below covers the case $\epsilon w_{\max}/n < 1$ (where the
scaling is vacuous and $\hat S = S^\ast$ exactly) in one line, and
spells out the remaining case in full.

\textbf{Vacuous-scaling case, $\epsilon w_{\max}/n < 1$.} Then $K=1$,
so $w_i' = \lfloor w_i \rfloor = w_i$ (the $w_i$ are already integers),
the value-indexed dynamic program is identical to Theorem~\ref{thm:dp}'s
program, and it returns $\hat S = S^\ast$, so $W(\hat S) = W^\ast \ge
(1-\epsilon)W^\ast$ immediately, since $\epsilon \in (0,1)$.

The remaining steps assume $\epsilon w_{\max}/n \ge 1$, so that $K =
\epsilon w_{\max}/n$ \emph{exactly} (not merely as an upper bound) ---
this exact equality, rather than an inequality, is what Step 5 below
needs, and is the reason the case split above is necessary rather than
cosmetic.

\textbf{Step 1: relating $w_i$ and $w_i'$.} By definition of the floor
function, $w_i' \le w_i/K < w_i' + 1$, so
\begin{equation}\label{eq:app-step1}
K w_i' \ \le\ w_i \ <\ K(w_i'+1).
\end{equation}
The right-hand inequality of \eqref{eq:app-step1}, rearranged, gives
$w_i \le K(w_i'+1)$, used in Step 4 below; the left-hand inequality
gives $w_i' \le w_i/K$, used in Step 2.

\textbf{Step 2: $S^\ast$'s scaled value is close to $W^\ast/K$.}
Summing the left-hand inequality of \eqref{eq:app-step1}'s consequence
$w_i' \le w_i/K$ is the wrong direction for a lower bound on $\sum_{i\in
S^\ast} w_i'$, so instead we use $w_i/K - 1 < w_i'$ (from $w_i/K < w_i'
+ 1$, i.e. the same right-hand inequality of \eqref{eq:app-step1}
divided by $K$ and rearranged), summed over $S^\ast$:
\begin{equation}\label{eq:app-step2}
\sum_{i \in S^\ast} w_i' \ >\ \sum_{i \in S^\ast}\Bigl(\frac{w_i}{K} -
1\Bigr) \ =\ \frac{W^\ast}{K} - |S^\ast| \ \ge\ \frac{W^\ast}{K} - n.
\end{equation}

\textbf{Step 3: the FPTAS's returned value is at least as good as
$S^\ast$'s scaled value.} $S^\ast$ is itself a feasible subset (it is
feasible by definition, being the true optimum), so it attains some
scaled value $\sum_{i\in S^\ast} w_i' =: v_0$ with $g(n,v_0) < \infty$.
Since $\hat S$ is defined as the feasible subset of \emph{largest}
scaled value among all $v$ with $g(n,v)<\infty$, its scaled value is at
least $v_0$:
\begin{equation}\label{eq:app-step3}
\sum_{i \in \hat S} w_i' \ \ge\ \sum_{i \in S^\ast} w_i' \ \overset{\eqref{eq:app-step2}}{>}\ \frac{W^\ast}{K} - n.
\end{equation}

\textbf{Step 4: converting $\hat S$'s scaled value back to true
weight.} Using the right-hand inequality of \eqref{eq:app-step1}
($w_i \le K(w_i'+1)$) would overcount by $Kn$ if applied termwise and
summed; instead we directly lower-bound $\hat S$'s true weight using
the \emph{left}-hand inequality of \eqref{eq:app-step1} ($K w_i' \le
w_i$), summed over $\hat S$, combined with \eqref{eq:app-step3}:
\begin{equation}\label{eq:app-step4}
W(\hat S) = \sum_{i \in \hat S} w_i \ \ge\ K \sum_{i\in \hat S} w_i' \
\overset{\eqref{eq:app-step3}}{>}\ K\Bigl(\frac{W^\ast}{K} - n\Bigr) =
W^\ast - Kn.
\end{equation}

\textbf{Step 5: $Kn$ equals $\epsilon w_{\max}$ exactly.} In this case
$K = \epsilon w_{\max}/n$ exactly (established above), so
\begin{equation}\label{eq:app-step5}
Kn \ =\ \epsilon\, w_{\max}.
\end{equation}
This is an equality, not merely a bound---the reason the case split
matters: in the vacuous-scaling case $K$ can instead exceed $\epsilon
w_{\max}/n$ (since $K=1$ there), which is exactly why that case is
handled separately rather than folded into this one.

\textbf{Step 6: bounding $w_{\max}$ by $W^\ast$.} By
Assumption~\ref{as:feasible}, the single highest-weight site is
individually feasible ($p_i \le d_i$), so dispatching it alone is a
feasible subset of value $w_{\max}$; since $W^\ast$ is the optimum over
\emph{all} feasible subsets, $w_{\max} \le W^\ast$.

\textbf{Combining Steps 4--6:}
\[
W(\hat S) \ >\ W^\ast - Kn \ =\ W^\ast - \epsilon\, w_{\max} \ \ge\
W^\ast - \epsilon W^\ast \ =\ (1-\epsilon) W^\ast,
\]
which gives $W(\hat S) > (1-\epsilon)W^\ast$ in this case, and
$W(\hat S) = W^\ast$ in the vacuous-scaling case; either way,
$W(\hat S) \ge (1-\epsilon)W^\ast$, matching Theorem~\ref{thm:fptas}.

\end{appendices}

\backmatter

\section*{Statements and Declarations}

\subsection*{Funding}
The authors declare that no funds, grants, or other support were
received during the preparation of this manuscript.

\subsection*{Competing Interests}
The authors have no relevant financial or non-financial interests to
disclose.

\subsection*{Author Contributions}
Q.V.D.\ conceived the problem and developed the theoretical results
(the hardness reduction, the exact dynamic program, the FPTAS, the
matroid-greedy special case, and the naive-policy lower bound).
H.V.V.\ implemented the algorithms, ran the numerical experiments and
the real-world case study, and produced the figures. M.N.D.\ conducted
the literature review and helped position the work relative to prior
scheduling and disaster-response research. P.S.N.\ designed the
real-world case study and its practical interpretation. Q.V.D.\ wrote
the first draft of the manuscript, and all authors commented on
previous versions, reviewed, and approved the final manuscript.

\subsection*{Data Availability}
The synthetic instance generators\DIFdelbegin \DIFdel{and all code used to produce every table }\DIFdelend \DIFaddbegin \DIFadd{, all algorithms, the exhaustive-search
verification, and the scripts that regenerate every table and figure }\DIFaddend in
Section~\ref{sec:experiments} \DIFaddbegin \DIFadd{(}\texttt{\DIFadd{experiment.py}}\DIFadd{,
}\texttt{\DIFadd{verify\_small.py}}\DIFadd{, }\texttt{\DIFadd{case\_study\_campfire.py}}\DIFadd{, and the
scripts in }\texttt{\DIFadd{figures/}}\DIFadd{) }\DIFaddend are included with this submission and are
self-contained \DIFdelbegin \DIFdel{, except for
}\DIFdelend \DIFaddbegin \DIFadd{and deterministic given their random seeds. The
real-world instance of }\DIFaddend Section~\ref{sec:campfire} \DIFdelbegin \DIFdel{'s real-world instance, whose inputs are
public-domain }\DIFdelend \DIFaddbegin \DIFadd{is derived from
public }\DIFaddend figures (a NIST technical report and U.S.\ Census data, both
cited in the text) rather than data collected by the authors; the
script that derives \DIFdelbegin \DIFdel{the instance from those public }\DIFdelend \DIFaddbegin \DIFadd{it from those }\DIFaddend figures and reproduces
Tables~\ref{tab:campfire} and~\ref{tab:campfire-results} \DIFaddbegin \DIFadd{and the
sensitivity analysis }\DIFaddend is also included\DIFdelbegin \DIFdel{with this submission}\DIFdelend .

\subsection*{Use of Large Language Models}
An AI coding assistant (Claude, Anthropic) was used to assist with
LaTeX formatting and to run the numerical illustration code reported in
Section~\ref{sec:experiments}. All theorem statements, proofs, and the
manuscript text were authored, verified, and are taken responsibility
for by the authors.

\bibliography{references}

\end{document}
